% Capítulo acadêmico para inclusão em documento principal por meio de:
% % Capítulo acadêmico para inclusão em documento principal por meio de:
% % Capítulo acadêmico para inclusão em documento principal por meio de:
% % Capítulo acadêmico para inclusão em documento principal por meio de:
% \input{metodologia_treinamento_bert_v01}
%
% Pacotes recomendados no preâmbulo do documento principal:
% \usepackage[T1]{fontenc}
% \usepackage[utf8]{inputenc} % dispensável em distribuições LaTeX recentes
% \usepackage[brazil]{babel}
% \usepackage{amsmath,amssymb}
% \usepackage{booktabs,longtable,array}
% \usepackage{graphicx}
% \usepackage{url}
% \usepackage[hidelinks]{hyperref}
% \usepackage[alf]{abntex2cite}
%
% As referências foram incorporadas ao final deste arquivo para que o capítulo
% não dependa de um arquivo .bib externo. Caso o documento principal já possua
% bibliografia própria, os itens podem ser transferidos para ela sem alterar as
% chaves usadas nas citações.

\chapter{Metodologia de treinamento dos modelos BERT}
\label{cap:metodologia-bert-v01}

\section{Visão geral e objetivos}
\label{sec:visao-geral-bert}

Este capítulo estabelece o protocolo para construção, treinamento, seleção e
avaliação dos primeiros classificadores supervisionados do projeto Radar Ódio
a partir do \textit{Dataset} v0.1. Essa versão será formada pelas avaliações
humanas resolvidas que correspondam a, no mínimo, 10\% do universo elegível. O
objetivo não é tratar a v0.1 como versão definitiva do conjunto de dados, mas
como o primeiro marco reprodutível de um processo contínuo de aprendizagem e
avaliação, a ser ampliado nas versões posteriores.

A estratégia emprega modelos pré-treinados da família BERT e derivados. O uso
de representações contextuais pré-treinadas permite adaptar um modelo de língua
a uma tarefa supervisionada mediante a adição de uma camada de classificação e
o ajuste conjunto dos parâmetros, procedimento conhecido como \textit{fine-tuning}
\cite{devlin2019bert}. Para o português brasileiro, os modelos BERTimbau
constituem referências monolíngues relevantes, tendo sido pré-treinados em um
grande corpus nacional e avaliados em diferentes tarefas de processamento de
linguagem natural \cite{souza2020bertimbau}. Modelos multilíngues serão incluídos
como controles comparativos, permitindo verificar se a especialização no idioma
produz vantagem no domínio de textos oriundos de redes sociais.

O protocolo contempla duas tarefas relacionadas, porém treinadas e avaliadas
separadamente:

\begin{enumerate}
    \item \textbf{detecção binária de discurso de ódio}: classificação de cada
    texto como \texttt{hate} ou \texttt{nao\_hate};
    \item \textbf{classificação multirrótulo dos tipos de ódio}: atribuição de
    zero ou mais categorias temáticas aos registros cujo rótulo gold seja
    \texttt{hate}.
\end{enumerate}

Essa separação define a arquitetura hierárquica candidata: o primeiro modelo
decide se o conteúdo contém discurso de ódio; quando a decisão é positiva, o
segundo identifica os tipos aplicáveis. Sua adoção não será tratada como uma
vantagem presumida. Ela será comparada a uma arquitetura conjunta e terá a
propagação de erros quantificada, conforme a Seção~\ref{sec:arquiteturas}.

\section{Perguntas de pesquisa, hipóteses e contribuição}
\label{sec:perguntas-hipoteses}

Os modelos são empregados como instrumentos para investigar generalização,
desbalanceamento, calibração e propagação de erros, e não apenas como alternativas
de implementação. As perguntas de pesquisa são:

\begin{description}
    \item[RQ1:] um encoder monolíngue pré-treinado em português brasileiro
    apresenta desempenho superior a encoders multilíngues comparáveis na
    detecção de discurso de ódio?
    \item[RQ2:] a ponderação da função de perda melhora a identificação da
    classe minoritária \texttt{hate} e dos tipos menos frequentes?
    \item[RQ3:] a calibração pós-processada reduz o erro das probabilidades
    estimadas sem deteriorar a discriminação?
    \item[RQ4:] o desempenho varia entre plataformas e se mantém quando uma
    plataforma inteira é excluída do treinamento?
    \item[RQ5:] qual é a perda de desempenho causada pela propagação de erros
    entre a detecção binária e a classificação dos tipos?
    \item[RQ6:] o desempenho se mantém em dados coletados depois do período
    representado pela v0.1?
    \item[RQ7:] o pré-treinamento específico em textos de redes sociais produz
    vantagem sobre o pré-treinamento geral da web, especialmente no X/Twitter?
\end{description}

As hipóteses e os contrastes primários serão registrados antes da abertura do
teste. Para RQ1, o contraste principal será BERTimbau Base contra mBERT, pois
ambos possuem porte semelhante; BERTimbau Large e XLM-R serão comparações
secundárias. Para RQ7, o contraste secundário pré-especificado será BERTabaporu
Base contra BERTimbau Base, mantendo as mesmas partições e o mesmo protocolo de ajuste. Seja
$M$ a métrica previamente indicada para cada pergunta. Serão
testadas as seguintes hipóteses direcionais:

\begin{align}
    H_{0,1}&: M_{\mathrm{BERTimbau\ Base}}-M_{\mathrm{mBERT}} \leq 0,
    &H_{1,1}&: M_{\mathrm{BERTimbau\ Base}}-M_{\mathrm{mBERT}} > 0;\\
    H_{0,2}&: M_{\mathrm{ponderada}}-M_{\mathrm{comum}} \leq 0,
    &H_{1,2}&: M_{\mathrm{ponderada}}-M_{\mathrm{comum}} > 0;\\
    H_{0,3}&: E_{\mathrm{calibrado}}-E_{\mathrm{bruto}} \geq 0,
    &H_{1,3}&: E_{\mathrm{calibrado}}-E_{\mathrm{bruto}} < 0;\\
    H_{0,5}&: F1_{\mathrm{oracle}}-F1_{\mathrm{ponta\ a\ ponta}} \leq 0,
    &H_{1,5}&: F1_{\mathrm{oracle}}-F1_{\mathrm{ponta\ a\ ponta}} > 0.
\end{align}

Em RQ3, $E$ designa o Brier score, com o ECE como medida complementar. RQ4 e
RQ6 serão avaliadas como hipóteses de heterogeneidade e degradação de domínio,
respectivamente. RQ6 é prospectiva e somente poderá receber resposta
confirmatória quando houver um teste temporal posterior à v0.1.

Além de significância estatística, será adotada como referência uma diferença
mínima cientificamente relevante de $\delta=0{,}02$ em Macro-F1 ou F1 da classe
\texttt{hate}. Esse valor constitui uma decisão operacional prévia, e não uma
constante universal; análises de sensibilidade usarão também $\delta=0{,}01$ e
$\delta=0{,}03$. A conclusão distinguirá: ausência de evidência, diferença
estatisticamente compatível com zero e diferença superior ao limiar de relevância
prática.

\begin{longtable}{p{0.06\textwidth}p{0.25\textwidth}p{0.30\textwidth}p{0.18\textwidth}}
    \caption{Matriz pré-especificada entre pergunta, hipótese, experimento e métrica.}
    \label{tab:matriz-rq} \\
    \toprule
    \textbf{RQ} & \textbf{Contraste} & \textbf{Experimento} & \textbf{Métrica principal} \\
    \midrule
    \endfirsthead
    \toprule
    \textbf{RQ} & \textbf{Contraste} & \textbf{Experimento} & \textbf{Métrica principal} \\
    \midrule
    \endhead
    RQ1 & monolíngue versus multilíngue & BERTimbau Base versus mBERT nas mesmas partições & Macro-F1 \\
    RQ2 & perda ponderada versus comum & ablação pareada por backbone e semente & F1-\texttt{hate}; AUPRC \\
    RQ3 & probabilidade calibrada versus bruta & \textit{temperature scaling} em conjunto próprio & Brier score; ECE \\
    RQ4 & dentro versus fora do domínio & resultados por plataforma e \textit{leave-one-platform-out} & Macro-F1 \\
    RQ5 & \textit{oracle} versus ponta a ponta & execução das duas formas sobre os mesmos exemplos & $\Delta$Macro-F1 \\
    RQ6 & período histórico versus posterior & teste temporal nas versões posteriores & Macro-F1; AUPRC \\
    RQ7 & pré-treinamento social versus web & BERTabaporu Base versus BERTimbau Base, globalmente e por plataforma & Macro-F1 \\
    \bottomrule
\end{longtable}

A contribuição científica esperada é, portanto, a caracterização controlada de
como idioma de pré-treinamento, desenho amostral, desbalanceamento, calibração,
plataforma e organização hierárquica afetam classificadores de discurso de ódio
em português brasileiro. A construção do dataset e a implementação dos modelos
são meios para essa investigação. Resultados confirmatórios serão separados de
análises exploratórias, especialmente para categorias raras e plataformas com
baixa cobertura.

\section{Definição e congelamento do Dataset v0.1}
\label{sec:dataset-v01}

\subsection{Retrato da base em 10 de setembro de 2026}

Na data de referência deste protocolo, o universo continha 212.476 análises com
estado de processamento bem-sucedido. A base gold reunia 25.090 textos canônicos
com decisão humana resolvida, correspondentes a 11,81\% do universo. Desse total,
3.995 registros haviam sido classificados como \texttt{hate} e 21.095 como
\texttt{nao\_hate}. A distribuição evidencia desbalanceamento entre as classes:
aproximadamente 15,92\% dos registros gold eram positivos e 84,08\% eram
negativos.

\begin{table}[htbp]
    \centering
    \caption{Retrato da base gold em 10 de setembro de 2026.}
    \label{tab:retrato-base-gold}
    \begin{tabular}{lrr}
        \toprule
        \textbf{Indicador} & \textbf{Quantidade} & \textbf{Percentual} \\
        \midrule
        Universo elegível & 212.476 & 100,00\% \\
        Decisões gold resolvidas & 25.090 & 11,81\% \\
        Rótulo \texttt{hate} & 3.995 & 15,92\% do gold \\
        Rótulo \texttt{nao\_hate} & 21.095 & 84,08\% do gold \\
        \bottomrule
    \end{tabular}
    \par\smallskip
    {\footnotesize Fonte: elaboração própria com base nas tabelas de avaliação
    do projeto.\par}
\end{table}

As decisões gold estavam distribuídas entre sete plataformas, conforme a
Tabela~\ref{tab:distribuicao-plataforma}. Essa distribuição deverá ser recalculada
no momento do congelamento, pois o acréscimo de avaliações pode alterar tanto as
contagens quanto as proporções.

\begin{table}[htbp]
    \centering
    \caption{Distribuição preliminar das decisões gold por plataforma.}
    \label{tab:distribuicao-plataforma}
    \begin{tabular}{lr}
        \toprule
        \textbf{Plataforma} & \textbf{Quantidade} \\
        \midrule
        YouTube & 7.735 \\
        Facebook & 4.891 \\
        Reddit & 3.101 \\
        WhatsApp & 3.037 \\
        Twitter & 2.895 \\
        Telegram & 2.660 \\
        TikTok & 771 \\
        \midrule
        \textbf{Total} & \textbf{25.090} \\
        \bottomrule
    \end{tabular}
    \par\smallskip
    {\footnotesize Fonte: elaboração própria com base nas tabelas de avaliação
    do projeto.\par}
\end{table}

Entre os textos positivos, uma mesma publicação pode receber mais de um tipo de
ódio. Por isso, a soma das frequências da Tabela~\ref{tab:tipos-odio-preliminar}
é superior ao número total de textos classificados como \texttt{hate}.

\begin{table}[htbp]
    \centering
    \caption{Frequência preliminar dos tipos de ódio na base gold.}
    \label{tab:tipos-odio-preliminar}
    \begin{tabular}{llr}
        \toprule
        \textbf{Código} & \textbf{Categoria} & \textbf{Ocorrências} \\
        \midrule
        \texttt{political} & Ódio político & 1.841 \\
        \texttt{sexism} & Sexismo & 1.020 \\
        \texttt{lgbtphobia} & LGBTfobia & 725 \\
        \texttt{aporophobia} & Classismo & 316 \\
        \texttt{religious\_intolerance} & Intolerância religiosa & 255 \\
        \texttt{racism} & Racismo & 241 \\
        \texttt{xenophobia} & Xenofobia & 141 \\
        \texttt{body\_shame} & Gordofobia & 123 \\
        \texttt{ageism} & Ageísmo & 52 \\
        \texttt{capacitism} & Capacitismo & 39 \\
        \texttt{other} & Categoria histórica inativa & 1 \\
        \bottomrule
    \end{tabular}
    \par\smallskip
    {\footnotesize Fonte: elaboração própria. As frequências são multirrótulo e,
    portanto, não devem ser somadas para estimar o número de textos positivos.\par}
\end{table}

As categorias menos frequentes ainda não oferecem quantidade suficiente para
uma avaliação confirmatória estável. Esse problema não será resolvido por meio
da eliminação automática de categorias ou pela criação de exemplos artificiais
antes da separação dos dados. Em vez disso, todos os rótulos ativos serão
preservados no treinamento e os resultados das classes raras serão explicitamente
marcados como exploratórios e acompanhados de intervalos de confiança.

\subsection{Justificativa do tamanho da v0.1}

O percentual de 10\% constitui um marco operacional para iniciar o ciclo de
treinamento, e não uma justificativa estatística isolada. O retrato efetivamente
alcançado reúne 25.090 registros e 3.995 positivos, com prevalência observada de
15,92\%. Sob amostragem aleatória simples, a margem de erro aproximada de 95\%
para essa prevalência seria de 0,45 ponto percentual. Esse cálculo não será
interpretado como precisão efetiva porque a base foi estratificada, possui pesos
desiguais e contém grupos correlacionados.

Antes do congelamento será produzido um relatório de suficiência com: contagem
esperada e observada de positivos; contagem por categoria, plataforma e tipo de
conteúdo; tamanho efetivo da amostra; largura dos intervalos de confiança
bootstrap; e quantidade de positivos em cada partição. Para pesos $w_i$, o
tamanho efetivo descritivo será calculado por

\begin{equation}
    n_{\mathrm{efetivo}} =
    \frac{\left(\sum_i w_i\right)^2}{\sum_i w_i^2}.
\end{equation}

Uma análise de poder por simulação, baseada apenas na parcela de desenvolvimento
disponível, estimará a probabilidade de detectar uma diferença pareada de
$\delta=0{,}02$ nas comparações primárias, com nível de 5\% e poder-alvo de 80\%.
A simulação preservará prevalência, estratos e grupos e apresentará também a
menor diferença detectável sob o tamanho alcançado. Seus parâmetros e código
serão arquivados; nenhuma observação do teste será usada para ajustar a análise.

O estudo binário será considerado apto à análise confirmatória apenas se o teste
contiver quantidade suficiente de positivos para produzir estimativas estáveis e
se os intervalos alcançarem precisão compatível com as comparações primárias.
Para os tipos de ódio, a classificação confirmatória ou exploratória obedecerá
aos critérios da Seção~\ref{sec:tarefa-multirrotulo}. Assim, ultrapassar o limiar
de 21.248 registros autoriza o início do benchmark, mas não transforma automaticamente
todas as análises por subgrupo em confirmatórias.

\subsection{Regra de congelamento}

Dez por cento de 212.476 registros correspondem a 21.247,6 observações. Assim,
o limiar operacional da v0.1 será de 21.248 decisões gold válidas. Como as
avaliações são organizadas em lotes, o congelamento ocorrerá ao término do
primeiro lote completo que fizer a base atingir ou ultrapassar esse limiar. Não
será realizado corte dentro de um lote já iniciado, evitando uma seleção
oportunista de casos.

O congelamento deverá produzir um instantâneo imutável, acompanhado de:

\begin{itemize}
    \item data e hora da extração;
    \item consulta SQL ou código de exportação versionado;
    \item número total de registros e distribuição dos rótulos;
    \item identificador do último lote incorporado;
    \item hash SHA-256 do arquivo exportado;
    \item versão da taxonomia e do protocolo de avaliação;
    \item histórico de exclusões e correções aplicadas.
\end{itemize}

Em 10 de setembro de 2026 foi produzida uma exportação candidata com 25.090
decisões válidas. O volume já supera o limiar mínimo da v0.1. Entretanto, o lote
24 ainda constava administrativamente como \texttt{active}, embora todos os seus
953 itens possuíssem decisão resolvida. Por isso, o snapshot pode iniciar a
preparação experimental, mas seu congelamento definitivo deverá registrar o
encerramento administrativo do lote e confirmar novamente as contagens e hashes.

Uma vez congelada, a v0.1 não deverá receber alterações silenciosas. Correções
posteriores de rótulos deverão gerar uma nova versão do conjunto de dados, com
registro das diferenças em relação à versão anterior.

O manifesto de congelamento deverá identificar separadamente as versões do
dataset, da taxonomia e do manual de anotação, por exemplo:

\begin{quote}
\texttt{dataset = radarodio-v0.1}\par
\texttt{taxonomy = taxonomy-v1.2}\par
\texttt{annotation-guidelines = annotation-guidelines-v1.2}
\end{quote}

Mudanças posteriores no manual ou na taxonomia não modificarão retroativamente
o significado de decisões produzidas sob uma versão anterior.

\subsection{Qualidade da anotação e formação do gold standard}
\label{sec:qualidade-anotacao}

Não será realizada nova dupla anotação sistemática da v0.1, pois esse procedimento
é operacionalmente inviável no cronograma do projeto. O controle de confiabilidade
utilizará os arquivos históricos de dupla anotação do Reddit como estudo auxiliar,
mantido separado do teste principal. A auditoria disponível identificou 1.644
pares binários válidos, dos quais 1.512 apresentaram consenso e 132 divergência,
correspondendo a 91,97\% de concordância bruta e $\kappa$ de Cohen de aproximadamente
0,39. Os 1.512 casos consensuais somente poderão ser incorporados depois de
vinculação ao corpus, deduplicação e harmonização da taxonomia; divergências sem
adjudicação permanecerão fora do gold.

Antes da anotação definitiva, o projeto deverá arquivar uma ficha metodológica
com o número e o perfil funcional dos avaliadores, treinamento recebido, versão
do manual, exemplos de treinamento, exercício de qualificação, critérios de
aprovação e procedimento para dúvidas. Informações pessoais dos avaliadores não
entrarão no arquivo destinado ao treinamento dos modelos.

Na tarefa binária, a concordância será descrita pelo percentual bruto, pelo
kappa de Cohen e pelo alfa de Krippendorff. Na tarefa multirrótulo serão
calculados o Jaccard e o F1 entre conjuntos, além da concordância e do alfa de
Krippendorff por categoria \cite{cohen1960kappa,krippendorff2018content}. Todos
os coeficientes serão acompanhados do número de
itens, distribuição marginal e intervalo de confiança. A concordância não será
usada isoladamente como prova de validade; desacordos serão analisados por
categoria, plataforma e fenômeno linguístico.

Haverá divergência quando os avaliadores históricos discordarem sobre o rótulo
binário ou sobre ao menos um tipo ativo. Sem nova adjudicação rastreável, esses
casos não serão transformados em rótulos de treinamento. Se uma adjudicação
pontual vier a ocorrer, deverá registrar decisão, justificativa, versão das
diretrizes, data e categorias em disputa, preservando as duas respostas originais.

Na data de referência, todas as 25.090 decisões resolvidas no sistema tinham
estado \texttt{single}. A v0.1 poderá ser usada para desenvolvimento e avaliação
interna, mas a ausência de dupla anotação contemporânea e multiplataforma limita
a força das conclusões sobre confiabilidade do gold. O estudo histórico do Reddit
será reportado como evidência auxiliar, sem generalização automática para as
demais plataformas.

\subsection{Critérios de elegibilidade}

Serão incluídos somente textos canônicos cujo consenso possua um dos estados
\texttt{single}, \texttt{agreed} ou \texttt{adjudicated}. Serão excluídos:

\begin{itemize}
    \item avaliações puladas ou ainda atribuídas;
    \item conflitos que não tenham sido adjudicados;
    \item registros sem texto ou com falha de codificação não recuperável;
    \item registros sem rótulo final consistente;
    \item cópias não canônicas de textos duplicados;
    \item observações que violem as regras vigentes da taxonomia.
\end{itemize}

Na data de referência, todas as 25.090 decisões resolvidas tinham estado
\texttt{single}. Consequentemente, não é possível estimar concordância entre
avaliadores para essa parcela da base. Essa condição será registrada como uma
limitação da v0.1. Os casos históricos de revisão dupla serão mantidos em um
conjunto auxiliar com proveniência própria e não alterarão silenciosamente o
estado das decisões produzidas pelo sistema.

O código histórico \texttt{other}, atualmente inativo, não integrará o vetor de
saída da tarefa multirrótulo. O registro que ainda contém esse código poderá ser
mantido na tarefa binária, mas deverá ser excluído da tarefa de tipos até que
seja revisto e recodificado em uma categoria ativa.

Antes da exportação serão aplicadas as seguintes invariantes:

\begin{align}
    \texttt{nao\_hate} &\Rightarrow
    \texttt{gold\_hate\_types}=\varnothing,\\
    \texttt{hate} &\Rightarrow
    \left|\texttt{gold\_hate\_types}\right|\geq 1,
\end{align}

em que todos os tipos devem pertencer à versão ativa da taxonomia. A condição
\texttt{hate} sem tipo não será silenciosamente convertida em negativa: o caso
será devolvido para adjudicação ou ficará fora da tarefa multirrótulo até ser
resolvido. Também serão verificadas unicidade de \texttt{analysis\_id}, presença
dos campos obrigatórios, reprodução do \texttt{text\_hash}, ausência de tipos
inativos e inexistência de rótulos conflitantes dentro de um grupo de duplicatas.
O relatório dessas validações será armazenado com o snapshot.

\subsection{Contrato de exportação}

A extração para treinamento deverá seguir o contrato apresentado na
Tabela~\ref{tab:contrato-exportacao}. Campos administrativos e identidades dos
avaliadores não serão incorporados ao conjunto utilizado pelos modelos.

\begin{longtable}{p{0.22\textwidth}p{0.16\textwidth}p{0.48\textwidth}}
    \caption{Contrato mínimo de exportação do Dataset v0.1.}
    \label{tab:contrato-exportacao} \\
    \toprule
    \textbf{Campo} & \textbf{Tipo} & \textbf{Finalidade} \\
    \midrule
    \endfirsthead
    \toprule
    \textbf{Campo} & \textbf{Tipo} & \textbf{Finalidade} \\
    \midrule
    \endhead
    \texttt{analysis\_id} & inteiro & Identificador interno e rastreabilidade da análise. \\
    \texttt{text\_hash} & texto & Identificação de duplicatas exatas e definição dos grupos de partição. \\
    \texttt{text} & texto & Conteúdo anonimizado que será fornecido ao tokenizer. \\
    \texttt{platform} & categoria & Auditoria e avaliação estratificada por plataforma. \\
    \texttt{content\_kind} & categoria & Identificação de post ou comentário. \\
    \texttt{gold\_label} & categoria & Alvo binário: \texttt{hate} ou \texttt{nao\_hate}. \\
    \texttt{gold\_hate\_types} & lista & Conjunto dos tipos de ódio aplicáveis ao texto. \\
    \texttt{decision\_status} & categoria & Indicação de decisão simples, concordante ou adjudicada. \\
    \texttt{batch\_id} & inteiro & Identificação do lote de amostragem. \\
    \texttt{stratum} & texto & Estrato usado na seleção e no cálculo de pesos amostrais. \\
    \texttt{inclusion\_probability} & numérico & Probabilidade de inclusão, quando conhecida e auditável. \\
    \texttt{duplicate\_group\_id} & texto & Grupo transitivo de duplicatas e quase duplicatas. \\
    \texttt{context\_group\_id} & texto/nulo & Conversa, tópico ou publicação de origem compartilhada. \\
    \texttt{content\_created\_at} & data/hora/nulo & Data do conteúdo para avaliação temporal. \\
    \texttt{reviewed\_at} & data/hora & Controle da versão e avaliações temporais futuras. \\
    \texttt{taxonomy\_version} & texto & Versão da taxonomia aplicada ao rótulo. \\
    \texttt{guidelines\_version} & texto & Versão do manual de anotação. \\
    \texttt{context\_dependency} & categoria/nulo & Indício de dependência de imagem, vídeo, conversa ou evento externo. \\
    \bottomrule
\end{longtable}

Os campos \texttt{gemma\_label}, \texttt{gemma\_hate\_probability} e
\texttt{gemma\_hate\_types} poderão ser exportados em uma tabela auxiliar
restrita à auditoria, ao cálculo dos estratos e à reprodução do baseline. Eles
não poderão ser concatenados ao texto nem empregados como atributos de entrada
dos classificadores BERT. Essa separação impede que os novos modelos apenas
reproduzam os sinais do sistema que orientou a amostragem.

\section{Preparação dos textos}
\label{sec:preparacao-textos}

O pré-processamento será deliberadamente conservador. Modelos BERT utilizam
tokenização em subpalavras e foram concebidos para aprender representações a
partir de sequências próximas da linguagem observada. Remoções agressivas podem
eliminar sinais relevantes para discurso de ódio, especialmente pontuação,
emojis, hashtags, caixa, erros ortográficos e alterações intencionais de
grafia.

Cada texto passará, nesta ordem, pelas seguintes operações:

\begin{enumerate}
    \item conversão para Unicode e normalização na forma NFC;
    \item substituição de sequências inválidas de codificação quando a correção
    for inequívoca;
    \item padronização de quebras de linha e espaços consecutivos;
    \item substituição de URLs pelo marcador \texttt{<URL>};
    \item substituição de nomes de usuário e menções pelo marcador
    \texttt{<USER>};
    \item substituição de endereços de e-mail por \texttt{<EMAIL>};
    \item substituição de números de telefone por \texttt{<PHONE>};
    \item remoção de espaços nas extremidades, sem alterar o restante do texto.
\end{enumerate}

Os marcadores serão adicionados como tokens especiais indivisíveis ao tokenizer
de cada checkpoint, com redimensionamento correspondente da matriz de embeddings.
As novas linhas serão inicializadas pelo procedimento padrão definido na
configuração do checkpoint, treinadas durante o fine-tuning e controladas pelas
mesmas sementes das demais inicializações. A lista de tokens, os IDs resultantes
e o hash do tokenizer alterado serão registrados. Uma ablação no conjunto de
desenvolvimento comparará o mascaramento com a manutenção das formas originais,
após anonimização irreversível dos identificadores, para medir eventual perda de
sinais úteis.

Não serão aplicados \textit{stemming}, lematização, remoção de \textit{stopwords},
tradução, conversão geral para minúsculas ou correção ortográfica. Os modelos
sensíveis à caixa receberão a capitalização original. Expressões regulares,
dicionários e versões de qualquer reconhecedor de informação pessoal serão
versionados. Nomes de grupos sociais necessários à interpretação do alvo não
serão removidos como se fossem identificadores pessoais. O texto bruto ficará
em ambiente de acesso restrito; somente a versão anonimizada poderá integrar os
artefatos de treinamento.

O BERTabaporu é um modelo \textit{uncased}; portanto, sua normalização de caixa
será realizada pelo tokenizer do próprio checkpoint. O arquivo de entrada comum
continuará preservando a capitalização original, evitando que a inclusão desse
modelo modifique os dados recebidos pelos encoders sensíveis à caixa.

\subsection{Duplicatas e prevenção de vazamento}

Primeiramente, o mapa de duplicatas do sistema será usado para manter uma única
linha por \texttt{text\_hash} canônico. Em seguida, será criada uma representação
normalizada em 5-gramas de caracteres. Pares com similaridade de Jaccard maior
ou igual a 0,90 serão tratados como quase duplicatas para fins de partição. A
detecção de quase duplicatas não implicará exclusão automática: os registros
serão reunidos sob um mesmo identificador de grupo, garantindo que permaneçam
na mesma partição e no mesmo fold.

O limiar 0,90 será validado antes da divisão definitiva. Serão sorteados pares
dos intervalos 0,80--0,89, 0,90--0,94 e 0,95--1,00 para julgamento humano cego,
registrando precisão e revocação aproximadas do detector de quase duplicatas.
Se a inspeção indicar separação inadequada, qualquer alteração do limiar será
feita e documentada antes da abertura do teste.

O agrupamento será transitivo. Os textos serão vértices de um grafo; cada par
que satisfizer o limiar formará uma aresta; e cada componente conexo receberá
um único \texttt{duplicate\_group\_id}. Assim, se $A$ for semelhante a $B$ e
$B$ a $C$, os três permanecerão juntos mesmo quando a similaridade direta entre
$A$ e $C$ for inferior a 0,90.

Quando houver identificador confiável de conversa, tópico ou encadeamento, todos
os textos do mesmo contexto também serão mantidos no mesmo grupo. Considera-se
contexto compartilhado a mesma postagem original e seus comentários, a mesma
thread, o mesmo encadeamento de respostas ou uma conversa identificada pelo
mesmo ID de origem. Um \texttt{context\_group\_id} prevalecerá sobre a decisão
individual de partição. Essa regra reduz o risco de um comentário no treinamento
revelar o conteúdo de outro comentário relacionado presente no teste.

\subsection{Comprimento das sequências}

Antes de definir o truncamento, será calculada, separadamente para cada tokenizer,
a distribuição do comprimento dos textos. Serão reportados média, mediana,
percentis 90, 95 e 99, máximo e as proporções $P(L>128)$, $P(L>256)$ e
$P(L>512)$. O limite técnico superior será de 512 tokens, incluindo os tokens
especiais do encoder.

Três políticas serão comparadas no desenvolvimento: (A) primeiros 510 tokens de
conteúdo; (B) 255 tokens iniciais e 255 finais; e (C) janelas de até 510 tokens
com sobreposição, agregadas no nível do documento. Para a política C, tamanho da
janela, sobreposição e regra de agregação serão fixados antes do teste. Se menos
de 5\% dos textos excederem 512 tokens e não houver diferença relevante no
desenvolvimento, a política B será adotada por parcimônia. Caso contrário, a
política será tratada como hiperparâmetro de ablação. A quantidade de textos
truncados, o conteúdo médio descartado e o custo adicional do \textit{chunking}
serão registrados.

\section{Definição das tarefas de aprendizagem}
\label{sec:tarefas-aprendizagem}

\subsection{Tarefa A: detecção binária}

Para um texto $x_i$, o alvo binário será definido por

\begin{equation}
    y_i =
    \begin{cases}
        1, & \text{se o rótulo gold for \texttt{hate}},\\
        0, & \text{se o rótulo gold for \texttt{nao\_hate}}.
    \end{cases}
\end{equation}

O vetor contextual associado ao token de classificação será fornecido a uma
camada linear com duas saídas, seguida por \textit{softmax}. A função de perda
será a entropia cruzada. Serão comparadas duas condições: perda sem ponderação e
perda ponderada por classe. Na condição ponderada, o peso da classe $c$ será
calculado apenas no subconjunto de treinamento de cada fold:

\begin{equation}
    \alpha_c = \frac{N_{\mathrm{treino}}}{2N_c},
\end{equation}

em que $N_c$ é o número de exemplos da classe $c$ no fold de treinamento.

\subsection{Tarefa B: classificação multirrótulo}
\label{sec:tarefa-multirrotulo}

A tarefa multirrótulo utilizará somente os textos com rótulo gold
\texttt{hate}. Para cada texto, será construído um vetor binário
$\mathbf{y}_i \in \{0,1\}^{10}$, na seguinte ordem fixa:

\begin{enumerate}
    \item \texttt{ageism};
    \item \texttt{aporophobia};
    \item \texttt{body\_shame};
    \item \texttt{capacitism};
    \item \texttt{lgbtphobia};
    \item \texttt{political};
    \item \texttt{racism};
    \item \texttt{religious\_intolerance};
    \item \texttt{sexism};
    \item \texttt{xenophobia}.
\end{enumerate}

No banco de origem, o rótulo de sexismo ainda está armazenado sob o código
histórico \texttt{misogyny}. Durante a preparação da tarefa multirrótulo, esse
valor será mapeado deterministicamente para \texttt{sexism}, sem alterar o
arquivo bruto. A tabela de mapeamento e sua versão integrarão o manifesto do
dataset.

O cabeçote terá dez saídas independentes com ativação sigmoide. A perda será a
entropia cruzada binária com logits, implementada por
\texttt{BCEWithLogitsLoss}. Também serão comparadas uma condição sem ponderação
e outra com peso positivo por categoria. Para a categoria $\ell$, o peso será
calculado exclusivamente no fold de treinamento:

\begin{equation}
    p_{\ell} = \frac{N_{\ell}^{-}}{N_{\ell}^{+}},
\end{equation}

em que $N_{\ell}^{+}$ e $N_{\ell}^{-}$ são, respectivamente, as quantidades de
exemplos positivos e negativos para o rótulo $\ell$.

Uma categoria será considerada elegível para inferência confirmatória somente
se possuir pelo menos 100 ocorrências na v0.1 e pelo menos 20 ocorrências no
conjunto de teste. Esses valores são salvaguardas operacionais mínimas para
evitar folds praticamente vazios; não garantem, sozinhos, precisão inferencial.
Por isso, serão apresentados intervalos de confiança e uma análise de
sensibilidade exigindo 20, 30 e 50 positivos no teste. Categorias cujo intervalo
de 95\% para F1 permaneça excessivamente amplo serão descritas como exploratórias,
mesmo que ultrapassem o limite mínimo.

\section{Arquiteturas e propagação de erros}
\label{sec:arquiteturas}

O experimento principal comparará duas organizações sob o mesmo encoder,
partições, pré-processamento e orçamento de ajuste:

\begin{enumerate}
    \item \textbf{hierárquica}: um classificador binário é executado primeiro e
    o classificador multirrótulo somente produz tipos quando o primeiro estágio
    prevê \texttt{hate};
    \item \textbf{conjunta não hierárquica}: um único encoder possui onze saídas
    sigmoides, uma para \texttt{hate} e dez para os tipos. Registros
    \texttt{nao\_hate} recebem zero nos dez tipos, e nenhuma regra de bloqueio é
    aplicada antes da medição. Inconsistências, como prever um tipo sem prever
    \texttt{hate}, serão contabilizadas.
\end{enumerate}

Na arquitetura conjunta, a perda será a soma ponderada da perda binária e da
perda dos tipos,

\begin{equation}
    \mathcal{L}_{\mathrm{conjunta}} =
    \mathcal{L}_{\mathrm{binária}} +
    \gamma\mathcal{L}_{\mathrm{tipos}},
\end{equation}

com $\gamma$ definido exclusivamente no desenvolvimento. Para limitar o custo,
essa ablação será executada inicialmente com o backbone vencedor da tarefa
binária; uma extensão aos demais backbones será identificada como análise
secundária.

O classificador de tipos será avaliado tanto com a entrada binária gold
(condição \textit{oracle}) quanto no pipeline completo. A propagação de erros
será quantificada por

\begin{equation}
    \Delta F1_{\mathrm{propagação}} =
    F1_{\mathrm{oracle}}-F1_{\mathrm{ponta\ a\ ponta}},
\end{equation}

globalmente, por categoria e por plataforma. A comparação hierárquica versus
conjunta utilizará predições pareadas no mesmo teste; qualquer vantagem deverá
ser interpretada junto do custo computacional e da taxa de inconsistências.

\section{Particionamento e desenho de validação}
\label{sec:particionamento}

\subsection{Teste fixo e conjunto de desenvolvimento}

Após a formação dos grupos de duplicatas, 20\% da v0.1 serão reservados como
conjunto de teste fixo. Dos 80\% restantes, 12,5\% serão separados previamente
como conjunto de calibração, equivalentes a 10\% da base completa. Os 70\%
restantes formarão o conjunto de seleção e ajuste. A proporção global será,
portanto, 70\% para seleção/treinamento, 10\% para calibração e limiares e 20\%
para teste. Todas as divisões respeitarão os grupos e a estratificação.

O conjunto de calibração terá três usos previamente limitados: estimar o
parâmetro de calibração, definir os limiares da configuração já selecionada e
medir a estabilidade dessas escolhas. Ele não será usado para escolher backbone,
taxa de aprendizado, perda, época ou política de truncamento. O teste será
armazenado separadamente e não poderá participar de:

\begin{itemize}
    \item ajuste de hiperparâmetros;
    \item seleção de modelo ou de época;
    \item definição de pesos de classe;
    \item escolha de limiares de decisão;
    \item calibração de probabilidades ou escolha de seus métodos;
    \item análise exploratória destinada a alterar o pipeline.
\end{itemize}

A divisão será realizada com estratificação iterativa multirrótulo, abordagem
destinada a conservar distribuições de rótulos em problemas nos quais uma mesma
instância pode pertencer a várias classes \cite{sechidis2011stratification}. O
vetor de estratificação combinará o rótulo binário, os dez tipos de ódio, a
plataforma, o tipo de conteúdo e o estrato de amostragem. A unidade indivisível
será o grupo de duplicatas, e não o registro individual.

\subsection{Validação cruzada}

O conjunto de seleção e ajuste será dividido em cinco folds por meio da mesma
estratégia de agrupamento e estratificação. Em cada rodada, quatro folds serão
usados para treinamento e um para validação. Ao final das cinco rodadas, cada
registro desse conjunto terá exatamente uma predição fora da amostra,
denominada \textit{out-of-fold} (OOF).

As predições OOF serão usadas para comparar configurações e estimar sua
variabilidade sem consultar o conjunto de calibração. A métrica OOF global,
calculada após concatenar as cinco predições de validação, será o critério
principal de seleção. A média e o desvio-padrão das cinco métricas calculadas
separadamente também serão reportados, mas não substituirão o critério global.
Essa escolha é especialmente importante porque o
fine-tuning de modelos pré-treinados pode apresentar variação relevante em
função da inicialização e da ordem dos dados, sobretudo em conjuntos menores
\cite{dodge2020finetuning}.

\subsection{Correção do desenho amostral}

Os lotes de avaliação foram balanceados a partir de informações como plataforma,
rótulo previsto pelo Gemma e tipo previsto de ódio. Logo, a distribuição da
base anotada pode diferir da distribuição do universo completo. O manifesto de
amostragem registrará o universo disponível em cada lote, definição e tamanho
dos estratos, quantidade selecionada, algoritmo, semente, critérios de exclusão,
uso das previsões do Gemma e qualquer priorização manual. Para cada estrato $h$,
será calculado o peso de expansão

\begin{equation}
    w_h = \frac{N_h}{n_h},
\end{equation}

em que $N_h$ é o número de registros elegíveis do estrato no universo e $n_h$
é o número de registros selecionados do mesmo estrato. A expressão corresponde
ao inverso da probabilidade de inclusão somente quando a seleção é conhecida e
equiprovável dentro de cada estrato. Para estabilidade numérica, os pesos serão
normalizados para média 1 dentro de cada partição.

As métricas do teste serão apresentadas em duas versões: não ponderada, que
descreve diretamente o conjunto anotado, e ponderada, que busca estimar o
desempenho sob a composição do universo original. Essa segunda interpretação
depende das hipóteses de inclusão conhecível, positividade e seleção ignorável
dentro dos estratos. Seleção manual, revisão dirigida ou exclusão subjetiva
viola essas hipóteses e não é corrigida automaticamente por $w_h$; nessa situação,
a versão ponderada será qualificada como ajuste descritivo, não como estimativa
representativa da população. Os pesos não serão calculados com base no rótulo
gold do teste e jamais serão derivados de seus erros.

Sempre que operacionalmente possível, será sorteada do universo uma amostra
probabilística independente das previsões do Gemma. Ela será anotada pelo mesmo
protocolo e usada como teste externo da prevalência, dos pesos e da generalização.
Se essa amostra não existir na v0.1, a ausência será registrada como limitação e
sua construção será prioridade da versão seguinte.

\section{Modelos comparados}
\label{sec:modelos-comparados}

Cinco encoders serão submetidos ao mesmo desenho experimental:

\begin{description}
    \item[BERTimbau Base:] \texttt{neuralmind/bert-base-portuguese-cased},
    modelo BERT Base monolíngue, sensível à caixa e pré-treinado para português
    brasileiro \cite{souza2020bertimbau,neuralmind2026bertimbau};
    \item[BERTimbau Large:] \texttt{neuralmind/bert-large-portuguese-cased},
    versão de maior capacidade do mesmo projeto
    \cite{souza2020bertimbau,neuralmind2026bertimbau};
    \item[BERTabaporu Base:]
    \texttt{pablocosta/bertabaporu-base-uncased}, modelo BERT Base monolíngue
    sem distinção de caixa, pré-treinado em aproximadamente 238 milhões de
    tweets em português brasileiro, provenientes de mais de 100 mil usuários e
    totalizando cerca de 2,9 bilhões de tokens
    \cite{dacosta2023bertabaporu,dacosta2026bertabaporu};
    \item[mBERT:] \texttt{google-bert/bert-base-multilingual-cased}, modelo BERT
    multilíngue sensível à caixa, pré-treinado em 104 idiomas
    \cite{devlin2019bert,google2026mbert};
    \item[XLM-R Base:] \texttt{FacebookAI/xlm-roberta-base}, encoder
    multilíngue baseado em RoBERTa, incluído como comparador da família de
    Transformers bidirecionais derivados de BERT
    \cite{conneau2020xlmr,meta2026xlmr}.
\end{description}

O BERTabaporu permitirá testar diretamente o efeito do gênero de pré-treinamento:
BERTimbau Base representa o português geral da web, enquanto BERTabaporu Base
representa a linguagem de redes sociais. Como a v0.1 é multiplataforma, eventual
vantagem no X/Twitter não será generalizada para as demais fontes. Serão
reportados o resultado global, o resultado por plataforma e a interação entre
modelo e plataforma. A versão
\texttt{pablocosta/bertabaporu-large-uncased} ficará como extensão secundária,
condicionada ao orçamento computacional, e não substituirá a execução do Base
no benchmark principal.

Os checkpoints e tokenizers serão carregados por meio da biblioteca
Transformers \cite{wolf2020transformers}. A revisão exata de cada repositório
será fixada por hash de commit no início dos experimentos, evitando que uma
atualização remota altere silenciosamente os resultados.

\subsection{Baselines}

Três famílias de baselines serão mantidas:

\begin{enumerate}
    \item \textbf{baseline trivial}: previsão constante de \texttt{nao\_hate}
    na tarefa binária e previsão determinada apenas pela prevalência de cada
    categoria, estimada no treinamento, na tarefa multirrótulo;
    \item \textbf{TF--IDF com regressão logística}: concatenação de n-gramas de
    palavras de ordem 1--2 e n-gramas de caracteres de ordem 3--5, seguida por
    regressão logística. O parâmetro $C$ será selecionado entre 0,1, 1 e 10 no
    mesmo protocolo de cinco folds. Na tarefa multirrótulo será usado
    \texttt{OneVsRestClassifier}, com um classificador logístico independente
    para cada um dos dez tipos;
    \item \textbf{Gemma}: uso das predições originais registradas antes da
    revisão humana, sem novo ajuste. Esse baseline permitirá comparar os BERTs
    com o sistema responsável pelas classificações iniciais.
\end{enumerate}

O baseline TF--IDF será submetido às mesmas partições, agrupamentos e métricas.
Nenhum vocabulário ou transformador TF--IDF será ajustado fora do fold de
treinamento. Como as previsões do Gemma ajudaram a definir estratos e prioridades
de revisão, o desempenho desse baseline não será interpretado como proveniente
de uma amostra independente do próprio sistema. A possível dependência entre a
amostragem e seus erros será explicitada em todas as comparações.

\section{Protocolo de fine-tuning}
\label{sec:fine-tuning}

Todos os encoders receberão cabeçotes de classificação novos, inicializados
aleatoriamente. Os parâmetros do encoder e do cabeçote serão ajustados em
conjunto. A faixa inicial segue a prática de fine-tuning apresentada para BERT,
enquanto decisões adicionais serão verificadas em piloto no próprio domínio
\cite{devlin2019bert}. O protocolo comum será:

\begin{itemize}
    \item otimizador AdamW, com $\beta_1=0{,}9$, $\beta_2=0{,}999$ e
    $\epsilon=10^{-8}$;
    \item tamanho efetivo de lote igual a 32, obtido por acúmulo de gradientes
    quando necessário;
    \item limite superior de oito épocas, entendido como teto operacional e não
    como duração fixa;
    \item avaliação ao final de cada época e parada antecipada após duas épocas
    completas sem redução da perda de validação;
    \item seleção do checkpoint de menor perda de validação, antes da otimização
    de qualquer limiar;
    \item aquecimento linear durante os 10\% primeiros passos;
    \item decaimento linear da taxa de aprendizado após o aquecimento;
    \item norma máxima do gradiente igual a 1,0;
    \item precisão mista FP16 ou BF16 quando suportada pelo hardware;
    \item embaralhamento reprodutível dos exemplos a cada época.
\end{itemize}

A perda de validação foi escolhida para a parada antecipada porque independe de
um limiar de classificação; Macro-F1, F1-\texttt{hate} e AUPRC serão monitorados,
mas não mudarão o checkpoint durante essa etapa. O teto de oito épocas apenas
permite que a regra de paciência opere em tarefas multirrótulo de convergência
mais lenta. O artigo original de BERT usou buscas mais curtas, tipicamente entre
duas e quatro épocas \cite{devlin2019bert}; portanto, será realizado um piloto
com curvas de treinamento. Se as curvas mostrarem convergência sistemática antes
do teto, o custo será reduzido pela parada antecipada. Se permanecerem em melhora
na oitava época, qualquer ampliação exigirá emenda prévia ao protocolo.

Será avaliada a grade cartesiana formada por

\begin{equation}
    \eta \in \{10^{-5}, 2\times10^{-5}, 3\times10^{-5}, 5\times10^{-5}\}
\end{equation}

e

\begin{equation}
    \lambda \in \{0, 0{,}01\},
\end{equation}

em que $\eta$ representa a taxa de aprendizado e $\lambda$ o decaimento de
peso. O valor $10^{-5}$ foi acrescentado às taxas usuais de BERT para contemplar
ajustes mais conservadores, especialmente no modelo Large. A condição de perda
comum ou ponderada será um fator experimental controlado, e não uma decisão
tomada depois de observar o teste.

\subsection{Busca de hiperparâmetros em duas etapas}

Uma grade completa com cinco modelos, duas tarefas, oito combinações, cinco
folds e duas perdas produziria 800 fine-tunings antes das repetições por semente.
Para reduzir custo sem escolher configurações pelo teste, será usada a seguinte
triagem pré-especificada:

\begin{enumerate}
    \item para cada combinação de backbone, tarefa e condição de perda, as oito
    combinações de $\eta$ e $\lambda$ serão executadas por até duas épocas em uma
    divisão interna fixa do conjunto de seleção, com semente 42;
    \item as duas melhores combinações de cada célula serão continuadas até o
    teto de oito épocas, sujeitas à mesma parada antecipada;
    \item a melhor combinação de cada célula seguirá para os cinco folds. Se as
    duas primeiras diferirem menos de 0,005 em Macro-F1, ambas seguirão para
    evitar descarte baseado em ruído da triagem;
    \item a configuração será selecionada pela métrica OOF global, obtida pela
    concatenação das predições dos cinco folds com limiar 0,5. Somente depois
    dessa seleção serão definidos calibração e limiares.
\end{enumerate}

Sem empates, a triagem compreende 160 execuções curtas, 40 continuações e 100
ajustes em folds, em lugar de 800 ajustes completos. Antes da execução será
estimado o orçamento em GPU-horas a partir do piloto, incluindo as repetições
finais e a ablação de arquitetura. Se o orçamento disponível for insuficiente,
o escopo será reduzido por regra registrada antes do teste, priorizando as
perguntas de pesquisa, e não por resultados parciais convenientes.

Na tarefa binária, Macro-F1 será o critério de seleção; F1-\texttt{hate} e AUPRC
serão critérios secundários. Na tarefa multirrótulo, será usado Macro-F1 das
categorias previamente elegíveis. Empates inferiores a 0,005 serão resolvidos,
nesta ordem, por menor variância entre folds, menor perda OOF e menor taxa de
aprendizado. Número de parâmetros e tempo de inferência não serão usados para
desempatar taxas do mesmo backbone, pois permanecem essencialmente constantes.

\subsection{Repetições e definição do modelo final}

Após a seleção, cada condição experimental será treinada no conjunto de seleção
e ajuste com as sementes 13, 42, 97, 123 e 2026. O número de épocas será fixado
na mediana das melhores épocas dos cinco folds, evitando parada antecipada no
teste ou no conjunto de calibração. As probabilidades dos cinco modelos serão
preservadas para medir variabilidade de otimização.

Para a avaliação principal de cada condição, o preditor final será o ensemble
pré-especificado das cinco sementes, formado pela média dos logits antes da
calibração. Nenhuma semente será escolhida pelo desempenho no teste. Caso seja
necessário disponibilizar também um único checkpoint por restrição operacional,
a semente 42 será usada por definição prévia e seus resultados serão reportados
separadamente dos resultados do ensemble.

\section{Limiarização e calibração}
\label{sec:limiarizacao-calibracao}

Na tarefa binária, o ensemble produzirá um logit médio e uma probabilidade
estimada de \texttt{hate}. Depois de fixados backbone, hiperparâmetros, perda,
época e sementes, o limiar $\tau$ será escolhido exclusivamente no conjunto de
calibração. Serão testados valores entre 0,01 e 0,99 em incrementos de 0,01,
selecionando-se o valor que maximizar o Macro-F1. Em empate, será escolhido o
limiar mais próximo de 0,5 e, persistindo o empate, o maior limiar. O valor 0,5
também será reportado como referência.

Na tarefa multirrótulo, cada categoria $\ell$ terá um limiar próprio
$\tau_{\ell}$, escolhido na mesma grade e no mesmo conjunto para maximizar o F1
da categoria. Se houver menos de cinco exemplos positivos de uma categoria na
calibração, seu limiar será mantido em 0,5 e o resultado será tratado como
exploratório. A estabilidade será descrita reestimando os limiares em
reamostragens agrupadas da calibração e, apenas como diagnóstico, nas validações
dos cinco folds; serão reportados mediana, intervalo interquartil e amplitude.

As probabilidades da tarefa binária serão adicionalmente calibradas por
\textit{temperature scaling} ajustado nos logits médios das cinco sementes no
conjunto reservado de calibração. Essa separação evita transferir diretamente
um calibrador estimado sobre logits OOF, produzidos por modelos distintos, para
um modelo final reajustado. Redes neurais modernas podem produzir probabilidades
excessivamente confiantes, e o método oferece uma correção pós-processada simples
\cite{guo2017calibration}.

A temperatura não ponderada estimará calibração para a distribuição da amostra
anotada. Uma temperatura ponderada, obtida pela minimização da log-verossimilhança
negativa com pesos de inclusão, será estimada para a população somente se a
auditoria do desenho amostral validar esses pesos. Caso contrário, não será feita
alegação de que $P(\texttt{hate}\mid\text{amostra})$ corresponde a
$P(\texttt{hate}\mid\text{população})$. O teste receberá apenas transformações e
limiares já congelados. ECE será calculado com 15 intervalos de frequência
aproximadamente igual; também serão apresentados Brier score e diagramas de
confiabilidade antes e depois da calibração.

\section{Métricas de avaliação}
\label{sec:metricas-avaliacao}

\subsection{Tarefa binária}

O Macro-F1 será a métrica primária de comparação geral. F1-\texttt{hate} e AUPRC
serão métricas centrais para as hipóteses sobre a classe positiva, mas não
substituirão retroativamente a métrica definida para cada contraste na
Tabela~\ref{tab:matriz-rq}. Para duas classes $C$, Macro-F1 será calculado
por

\begin{equation}
    \operatorname{MacroF1} = \frac{1}{|C|}\sum_{c\in C}
    \frac{2P_cR_c}{P_c+R_c},
\end{equation}

em que $P_c$ e $R_c$ representam precisão e revocação da classe $c$. O uso de
Macro-F1 impede que o desempenho da classe majoritária domine sozinho a
comparação.

Também serão reportados:

\begin{itemize}
    \item precisão, revocação e F1 da classe \texttt{hate};
    \item área sob a curva Precisão--Revocação (AUPRC);
    \item área sob a curva ROC (AUROC);
    \item coeficiente de correlação de Matthews (MCC);
    \item acurácia balanceada;
    \item matriz de confusão;
    \item Brier score;
    \item erro esperado de calibração (ECE);
    \item curvas de confiabilidade antes e depois da calibração.
\end{itemize}

A AUPRC receberá destaque porque a classe positiva é minoritária. Em cenários
desbalanceados, curvas Precisão--Revocação são mais informativas sobre a
qualidade das predições positivas do que a curva ROC isoladamente
\cite{saito2015precisionrecall}.

\subsection{Tarefa multirrótulo}

A métrica primária será o Macro-F1 calculado sobre as categorias elegíveis para
análise confirmatória. Serão apresentados, adicionalmente:

\begin{itemize}
    \item Macro-F1 incluindo todos os rótulos ativos;
    \item Micro-F1;
    \item precisão, revocação, F1 e PR-AUC por categoria;
    \item \textit{Hamming loss};
    \item \textit{subset accuracy}, que exige correspondência exata do conjunto
    de rótulos;
    \item média da quantidade de rótulos previstos por texto;
    \item matriz de coocorrência entre tipos verdadeiros e previstos.
\end{itemize}

Serão realizadas duas avaliações complementares. A avaliação \textit{oracle}
aplicará o classificador de tipos somente aos registros cujo rótulo binário gold
seja \texttt{hate}, medindo a capacidade específica de reconhecer categorias.
A avaliação ponta a ponta executará primeiro o classificador binário e somente
atribuirá tipos aos textos previstos como positivos. A diferença entre os dois
resultados quantificará a propagação de erros da primeira etapa.

\subsection{Incerteza e comparação entre modelos}

Serão separadas duas fontes de incerteza. A variabilidade de otimização será
descrita pela média, desvio-padrão, mínimo e máximo entre as cinco sementes. A
incerteza amostral será estimada sobre a predição final pré-especificada por
2.000 reamostragens bootstrap do conjunto de teste \cite{efron1993bootstrap}. O
bootstrap será estratificado
pelos estratos de amostragem e terá como unidade o maior agrupamento entre
\texttt{duplicate\_group\_id} e \texttt{context\_group\_id}, preservando a
dependência entre textos relacionados. Na versão ponderada, os pesos serão
mantidos em cada reamostragem.

Diferenças entre modelos serão examinadas por bootstrap pareado, calculando
diretamente, em cada reamostragem, $\Delta M=M_A-M_B$ sobre os mesmos exemplos.
Não será inferida diferença pela simples sobreposição de intervalos individuais.
Os quatro contrastes primários da Tabela~\ref{tab:matriz-rq} serão apresentados
sem expansão posterior do conjunto de hipóteses. Para famílias de comparações
secundárias entre modelos, plataformas ou categorias, os valores de $p$ serão
ajustados pelo procedimento de Holm \cite{holm1979procedure}. Magnitude,
intervalo de confiança e relação
com $\delta$ serão priorizados sobre uma decisão binária de significância.

Além dos resultados globais, as métricas serão desagregadas por plataforma,
tipo de conteúdo e categoria de ódio. Subgrupos com menos de 30 exemplos
positivos não receberão interpretação comparativa; seus valores serão exibidos
apenas como evidência exploratória, acompanhados do tamanho da amostra e do
intervalo de confiança.

Para avaliar generalização entre plataformas, será realizado um experimento
\textit{leave-one-platform-out}: em cada rodada, uma plataforma será mantida
inteiramente fora do ajuste e usada como domínio de teste. A rodada somente será
confirmatória quando a plataforma excluída possuir ao menos 100 registros e 20
positivos; nos demais casos será exploratória. Os resultados serão reportados
para YouTube, WhatsApp, Telegram, Twitter/X, Facebook, Reddit e TikTok, sem
combinar artificialmente plataformas pequenas para obter significância.

Também serão registrados parâmetros, GPU-horas, memória máxima, duração do
treinamento e tempo de inferência por mil textos. A interpretação final comparará
ganho de desempenho e custo, evitando concluir que um modelo maior é preferível
quando a diferença estiver abaixo de $\delta$ e o custo for substancialmente
superior.

\section{Testes funcionais e análise de erros}
\label{sec:testes-funcionais}

Métricas agregadas em um conjunto de teste podem ocultar falhas sistemáticas de
modelos de detecção de discurso de ódio. O HateCheck demonstra a utilidade de
testes funcionais construídos para fenômenos linguísticos específicos
\cite{rottger2021hatecheck}. Com base nessa abordagem, será produzido um pequeno
conjunto diagnóstico em português brasileiro, separado de treinamento,
validação e teste estatístico.

O conjunto funcional terá, como meta inicial, pelo menos 20 exemplos por
fenômeno, equilibrando casos positivos e negativos quando isso for semanticamente
aplicável. Os exemplos poderão ser adaptados de taxonomias funcionais publicadas,
construídos de forma controlada ou extraídos do domínio, desde que origem e
licença sejam registradas. O conjunto deverá contemplar, no mínimo:

\begin{itemize}
    \item negação de uma afirmação preconceituosa;
    \item citação de discurso de ódio para denúncia ou análise;
    \item contra-discurso e defesa do grupo atacado;
    \item ironia e sarcasmo;
    \item menção neutra a grupos sociais;
    \item insulto individual sem ataque generalizante a um grupo;
    \item grafias alteradas, separação de caracteres e uso de símbolos;
    \item ameaça, incitação, desumanização e criminalização coletiva;
    \item ataques explícitos e implícitos;
    \item variações de gênero, número e referência pronominal.
\end{itemize}

Os casos serão classificados por dois avaliadores independentes e divergências
serão adjudicadas pelo protocolo da Seção~\ref{sec:qualidade-anotacao}. Antes do
uso, cada caso será comparado com treino, calibração e teste por hash, Jaccard e
busca aproximada; casos equivalentes serão removidos ou reescritos e novamente
validados. O conjunto funcional não será usado para ajustar hiperparâmetros. O
desempenho será relatado por fenômeno, sem misturá-lo à estimativa principal do
teste da v0.1.

Após a avaliação quantitativa, será conduzida análise qualitativa de erros. Para
cada modelo serão examinados, no mínimo, 50 falsos positivos e 50 falsos
negativos da tarefa binária, sorteados com semente 2026 e estratificação por
plataforma e faixa de confiança. Na tarefa multirrótulo serão sorteados erros das
categorias mais frequentes e incluídos os erros disponíveis das categorias raras,
respeitando as restrições de privacidade.

Antes de examinar as predições será congelada a seguinte taxonomia inicial:
\texttt{E01} contexto ausente; \texttt{E02} ironia; \texttt{E03} sarcasmo;
\texttt{E04} negação; \texttt{E05} alvo implícito; \texttt{E06} ambiguidade
lexical; \texttt{E07} linguagem codificada; \texttt{E08} truncamento;
\texttt{E09} possível erro de anotação; \texttt{E10} conflito entre categorias;
\texttt{E11} conhecimento externo necessário; e \texttt{E12} dependência
multimodal. Novos códigos poderão emergir, mas serão identificados como
pós-especificados. Uma parcela mínima de 20\% dos erros será codificada por dois
avaliadores, com concordância reportada.

Cada item indicará se depende de imagem, vídeo, mensagem anterior, thread,
evento externo ou conhecimento sociopolítico. Será reportada a proporção de
erros potencialmente insolúveis pelo texto isolado, distinguindo limitação do
modelo de limitação da modalidade de entrada.

\section{Reprodutibilidade, documentação e uso responsável}
\label{sec:reprodutibilidade}

Cada execução deverá gerar um manifesto contendo:

\begin{itemize}
    \item versão e hash SHA-256 do dataset;
    \item hash do código e da consulta de extração;
    \item commit do repositório, estado de alterações locais e comando executado;
    \item arquivos imutáveis de \texttt{train\_ids}, \texttt{validation\_ids},
    \texttt{calibration\_ids}, \texttt{test\_ids},
    \texttt{duplicate\_groups} e \texttt{context\_groups};
    \item checkpoint e revisão do tokenizer;
    \item versões de Python, PyTorch, Transformers e demais dependências;
    \item modelo, tarefa, hiperparâmetros e semente;
    \item número de épocas, melhor época nos folds e curvas de perda;
    \item hardware, precisão numérica, memória máxima, GPU-horas e tempo de
    treinamento e inferência;
    \item limiares, temperatura de calibração e pesos de classe;
    \item métricas globais, ponderadas e por subgrupo;
    \item caminho e hash dos checkpoints resultantes.
\end{itemize}

O ambiente de software será congelado por arquivo de dependências com versões
exatas e, preferencialmente, por imagem de contêiner. Não se dependerá apenas da
semente para reconstruir partições. As configurações de determinismo do PyTorch,
CUDA e cuDNN serão registradas, incluindo flags ativadas e operações conhecidas
como não determinísticas. Quando determinismo estrito reduzir de modo relevante
o desempenho computacional ou não estiver disponível, a exceção será documentada.

A documentação do dataset descreverá origem, composição, processo de amostragem,
população representada, critérios de anotação, limitações e usos não recomendados.
Esse procedimento é coerente com a proposta de \textit{data statements}, que
busca tornar explícitas as condições sob as quais resultados de PLN podem ou não
ser generalizados \cite{bender2018datastatements}. Cada modelo terá também uma
ficha com finalidade, desempenho, grupos avaliados, limitações, riscos e uso
pretendido, seguindo a proposta de \textit{model cards}
\cite{mitchell2019modelcards}.

Os classificadores não deverão ser descritos como substitutos da avaliação
humana. Seu uso previsto é priorizar, apoiar ou acelerar a triagem de conteúdos.
Decisões de alto impacto, estudos de prevalência e análises sobre grupos pouco
representados deverão considerar a incerteza, o desenho amostral e a revisão de
especialistas.

Exemplos divulgados em artigos, dissertação ou fichas de modelo passarão por
revisão específica de risco de reidentificação. Conforme o risco, serão usados
mascaramento adicional, paráfrase fiel, exemplo sintético controlado ou omissão.
Identificadores de origem e textos brutos não acompanharão os artefatos públicos.

\section{Continuidade do processo}
\label{sec:continuidade-processo}

O teste da v0.1 permanecerá congelado nas versões posteriores, permitindo medir
ganhos ou regressões no mesmo conjunto histórico. Contudo, após ter sido
consultado em sucessivos ciclos, ele será interpretado como benchmark histórico
ou teste de regressão, e não indefinidamente como única avaliação confirmatória.
Novos registros avaliados serão incorporados ao desenvolvimento das versões
v0.2, v0.3 e seguintes. Cada versão reservará um teste temporal formado por
conteúdos estritamente posteriores ao período de treinamento; o teste temporal
mais recente assumirá progressivamente o papel confirmatório.

Para cada atualização serão repetidos os controles de elegibilidade,
duplicidade, estratificação, treinamento e avaliação. A data de corte temporal
será definida antes do treinamento e os grupos de contexto não poderão atravessá-la.
Um modelo novo somente substituirá o anterior se não apresentar regressão
relevante no teste histórico e demonstrar desempenho aceitável no teste temporal
mais recente.

O pré-treinamento adaptativo em dados não rotulados do domínio poderá ser
investigado em uma etapa futura. Estudos anteriores mostram que a continuação do
pré-treinamento em textos do domínio ou da tarefa pode produzir ganhos mesmo em
cenários com poucos rótulos \cite{gururangan2020dontstop}. Contudo, esse
procedimento não integrará o benchmark principal da v0.1, que comparará os
checkpoints públicos originais sob o mesmo protocolo supervisionado.

\section{Estudos de ablação}
\label{sec:ablacoes}

As decisões de engenharia que possam alterar as conclusões serão transformadas
em contrastes controlados. As ablações pré-especificadas são:

\begin{itemize}
    \item perda comum versus perda ponderada;
    \item arquitetura hierárquica versus conjunta;
    \item encoder monolíngue versus multilíngue;
    \item probabilidades brutas versus calibradas;
    \item truncamento apenas no início versus início e fim e, se necessário,
    \textit{chunking};
    \item texto anonimizado com marcadores especiais versus anonimização sem os
    marcadores explícitos.
\end{itemize}

Cada ablação alterará um componente por vez, manterá partições e sementes e será
analisada de forma pareada. Perda, idioma, calibração e propagação hierárquica
correspondem às perguntas primárias; truncamento e mascaramento serão análises
secundárias no backbone vencedor, salvo se o piloto mostrar efeito suficiente
para justificar expansão previamente documentada.

\section{Fluxo experimental consolidado}
\label{sec:fluxo-experimental}

O procedimento completo seguirá a ordem abaixo. A posição do teste no fluxo
representa isolamento lógico: seus rótulos somente serão acessados depois que os
artefatos anteriores estiverem congelados.

\begin{enumerate}
    \item definir universo, estratos, probabilidades de inclusão e lotes;
    \item concluir a anotação simples, auditar a dupla anotação histórica e
    documentar concordância, divergências e exclusões;
    \item congelar dataset, taxonomia e manual, com hashes e manifests;
    \item executar anonimização, normalização e testes de consistência;
    \item detectar duplicatas, formar componentes conexos e grupos de contexto;
    \item criar e versionar as partições de seleção, calibração e teste;
    \item perfilar comprimentos e executar apenas as ablações pré-especificadas;
    \item realizar triagem de hiperparâmetros e validação cruzada;
    \item produzir e concatenar predições OOF para seleção de configuração;
    \item fixar época e treinar as cinco sementes finais;
    \item formar o ensemble, ajustar calibração e limiares no conjunto próprio;
    \item congelar checkpoint, tokenizer, calibrador, limiares e código;
    \item abrir uma única vez o teste para a avaliação principal;
    \item executar bootstrap pareado, análises por plataforma e eficiência;
    \item executar avaliações externa, temporal e funcional quando disponíveis;
    \item conduzir análise qualitativa reprodutível e publicar a documentação.
\end{enumerate}

\section{Ameaças à validade}
\label{sec:ameacas-validade}

As principais ameaças à interpretação dos resultados são:

\begin{itemize}
    \item \textbf{viés de seleção}: a amostra foi balanceada com base em sinais
    do Gemma e pode não reproduzir diretamente a prevalência do universo;
    \item \textbf{revisão simples}: a predominância de decisões
    \texttt{single} impede estimar concordância entre avaliadores na maior parte
    da v0.1;
    \item \textbf{classes raras}: algumas categorias possuem poucas dezenas de
    exemplos e não sustentam conclusões confirmatórias;
    \item \textbf{mudança temporal}: linguagem ofensiva, códigos e plataformas
    podem mudar depois do congelamento;
    \item \textbf{dependência de contexto}: o texto isolado pode não conter
    informações presentes na conversa, imagem, vídeo ou evento relacionado;
    \item \textbf{propagação de erros}: no pipeline hierárquico, um falso
    negativo binário impede a identificação dos tipos de ódio;
    \item \textbf{generalização limitada}: os resultados serão válidos para a
    composição observada de plataformas, períodos e categorias, não para todo
    conteúdo em português brasileiro.
\end{itemize}

Essas limitações deverão acompanhar qualquer divulgação de métricas ou uso do
modelo. A expansão contínua do dataset será orientada também pelos erros e pelas
lacunas de cobertura detectadas na v0.1.

\section{Tabelas para registro dos resultados}
\label{sec:tabelas-resultados}

As tabelas seguintes definem o formato de apresentação, mas não contêm valores
experimentais. Os campos marcados com ``---'' deverão ser preenchidos somente
após a execução do protocolo.

\begin{table}[htbp]
    \centering
    \caption{Modelo de tabela para a tarefa binária no teste fixo.}
    \label{tab:modelo-resultados-binarios}
    \resizebox{\textwidth}{!}{%
    \begin{tabular}{lrrrrrrr}
        \toprule
        \textbf{Modelo} & \textbf{Macro-F1} & \textbf{F1 hate} &
        \textbf{Precisão hate} & \textbf{Revocação hate} & \textbf{AUPRC} &
        \textbf{MCC} & \textbf{ECE} \\
        \midrule
        TF--IDF + regressão logística & --- & --- & --- & --- & --- & --- & --- \\
        Gemma & --- & --- & --- & --- & --- & --- & --- \\
        BERTimbau Base & --- & --- & --- & --- & --- & --- & --- \\
        BERTimbau Large & --- & --- & --- & --- & --- & --- & --- \\
        BERTabaporu Base & --- & --- & --- & --- & --- & --- & --- \\
        mBERT & --- & --- & --- & --- & --- & --- & --- \\
        XLM-R Base & --- & --- & --- & --- & --- & --- & --- \\
        \bottomrule
    \end{tabular}}
    \par\smallskip
    {\footnotesize Fonte: elaboração própria. Valores a serem preenchidos após
    os experimentos.\par}
\end{table}

\begin{table}[htbp]
    \centering
    \caption{Modelo de tabela para a tarefa multirrótulo.}
    \label{tab:modelo-resultados-multirrotulo}
    \resizebox{\textwidth}{!}{%
    \begin{tabular}{lrrrrrr}
        \toprule
        \textbf{Modelo} & \textbf{Macro-F1 elegível} &
        \textbf{Macro-F1 total} & \textbf{Micro-F1} & \textbf{PR-AUC} &
        \textbf{Hamming loss} & \textbf{Subset accuracy} \\
        \midrule
        TF--IDF + regressão logística & --- & --- & --- & --- & --- & --- \\
        BERTimbau Base & --- & --- & --- & --- & --- & --- \\
        BERTimbau Large & --- & --- & --- & --- & --- & --- \\
        BERTabaporu Base & --- & --- & --- & --- & --- & --- \\
        mBERT & --- & --- & --- & --- & --- & --- \\
        XLM-R Base & --- & --- & --- & --- & --- & --- \\
        \bottomrule
    \end{tabular}}
    \par\smallskip
    {\footnotesize Fonte: elaboração própria. Valores a serem preenchidos após
    os experimentos.\par}
\end{table}

\begin{table}[htbp]
    \centering
    \caption{Modelo de tabela para resultados por categoria de ódio.}
    \label{tab:modelo-resultados-categoria}
    \resizebox{\textwidth}{!}{%
    \begin{tabular}{lrrrrrr}
        \toprule
        \textbf{Categoria} & \textbf{$n$ teste} & \textbf{Precisão} &
        \textbf{Revocação} & \textbf{F1} & \textbf{IC95\% F1} & \textbf{PR-AUC} \\
        \midrule
        Ageísmo & --- & --- & --- & --- & --- & --- \\
        Classismo & --- & --- & --- & --- & --- & --- \\
        Gordofobia & --- & --- & --- & --- & --- & --- \\
        Capacitismo & --- & --- & --- & --- & --- & --- \\
        LGBTfobia & --- & --- & --- & --- & --- & --- \\
        Ódio político & --- & --- & --- & --- & --- & --- \\
        Racismo & --- & --- & --- & --- & --- & --- \\
        Intolerância religiosa & --- & --- & --- & --- & --- & --- \\
        Sexismo & --- & --- & --- & --- & --- & --- \\
        Xenofobia & --- & --- & --- & --- & --- & --- \\
        \bottomrule
    \end{tabular}}
    \par\smallskip
    {\footnotesize Fonte: elaboração própria. Categorias abaixo do limiar
    confirmatório deverão ser identificadas como exploratórias.\par}
\end{table}

\begin{table}[htbp]
    \centering
    \caption{Modelo de tabela para desempenho e eficiência computacional.}
    \label{tab:modelo-eficiencia}
    \resizebox{\textwidth}{!}{%
    \begin{tabular}{lrrrrr}
        \toprule
        \textbf{Modelo} & \textbf{Macro-F1} & \textbf{Parâmetros} &
        \textbf{GPU-horas} & \textbf{Memória máxima} &
        \textbf{Tempo/1.000 textos} \\
        \midrule
        BERTimbau Base & --- & --- & --- & --- & --- \\
        BERTimbau Large & --- & --- & --- & --- & --- \\
        BERTabaporu Base & --- & --- & --- & --- & --- \\
        mBERT & --- & --- & --- & --- & --- \\
        XLM-R Base & --- & --- & --- & --- & --- \\
        \bottomrule
    \end{tabular}}
    \par\smallskip
    {\footnotesize Fonte: elaboração própria. Valores a serem preenchidos após
    os experimentos.\par}
\end{table}

\section{Critérios de conclusão da etapa v0.1}
\label{sec:criterios-conclusao}

A etapa de treinamento da v0.1 será considerada metodologicamente concluída
quando:

\begin{enumerate}
    \item o primeiro lote completo que alcance ao menos 21.248 decisões válidas
    tiver sido congelado e versionado;
    \item a suficiência amostral tiver sido examinada e as análises confirmatórias
    e exploratórias estiverem separadas;
    \item a dupla anotação histórica, as medidas de concordância e as exclusões
    por divergência tiverem sido documentadas, e a ausência de dupla anotação
    contemporânea estiver explicitamente declarada como limitação;
    \item as partições tiverem sido verificadas quanto a duplicatas e vazamento;
    \item todos os baselines e cinco encoders tiverem sido executados nas mesmas
    divisões;
    \item os hiperparâmetros tiverem sido definidos por OOF e os limiares e a
    calibração pelo conjunto próprio, sem uso do teste;
    \item as cinco sementes tiverem sido executadas para cada configuração
    selecionada;
    \item o ensemble final tiver sido fixado sem seleção posterior de semente;
    \item as métricas globais, ponderadas, por plataforma e por categoria
    tiverem sido calculadas com intervalos de confiança e comparações pareadas;
    \item a análise funcional e a análise qualitativa de erros tiverem sido
    documentadas;
    \item os manifests, checkpoints, hashes, tabelas e fichas de documentação
    tiverem sido arquivados de forma reprodutível.
\end{enumerate}

% -----------------------------------------------------------------------------
% Referências incorporadas ao próprio capítulo
% -----------------------------------------------------------------------------

\begin{thebibliography}{99}

\bibitem[Bender e Friedman(2018)]{bender2018datastatements}
BENDER, Emily M.; FRIEDMAN, Batya.
Data statements for natural language processing: toward mitigating system bias
and enabling better science.
\textit{Transactions of the Association for Computational Linguistics}, v. 6,
p. 587--604, 2018.
DOI: \url{https://doi.org/10.1162/tacl_a_00041}.

\bibitem[Cohen(1960)]{cohen1960kappa}
COHEN, Jacob.
A coefficient of agreement for nominal scales.
\textit{Educational and Psychological Measurement}, v. 20, n. 1, p. 37--46,
1960.
DOI: \url{https://doi.org/10.1177/001316446002000104}.

\bibitem[Conneau et al.(2020)]{conneau2020xlmr}
CONNEAU, Alexis et al.
Unsupervised cross-lingual representation learning at scale.
In: \textit{Proceedings of the 58th Annual Meeting of the Association for
Computational Linguistics}. Online: Association for Computational Linguistics,
2020. p. 8440--8451.
Disponível em: \url{https://arxiv.org/abs/1911.02116}.

\bibitem[da Costa et al.(2023)]{dacosta2023bertabaporu}
DA COSTA, Pablo Botton; PAVAN, Matheus Camasmie; DOS SANTOS, Wesley Ramos;
DA SILVA, Samuel Caetano; PARABONI, Ivandre.
BERTabaporu: assessing a genre-specific language model for Portuguese NLP.
In: \textit{Proceedings of the 14th International Conference on Recent Advances
in Natural Language Processing}. Varna: INCOMA, 2023. p. 217--223.
DOI: \url{https://doi.org/10.26615/978-954-452-092-2_024}.

\bibitem[da Costa(2026)]{dacosta2026bertabaporu}
DA COSTA, Pablo Botton.
\textit{BERTabaporu Base Uncased: model card}.
Hugging Face, 2026. Disponível em:
\url{https://huggingface.co/pablocosta/bertabaporu-base-uncased}.
Acesso em: 11 set. 2026.

\bibitem[Devlin et al.(2019)]{devlin2019bert}
DEVLIN, Jacob; CHANG, Ming-Wei; LEE, Kenton; TOUTANOVA, Kristina.
BERT: pre-training of deep bidirectional transformers for language understanding.
In: \textit{Proceedings of the 2019 Conference of the North American Chapter of
the Association for Computational Linguistics: Human Language Technologies}.
Minneapolis: Association for Computational Linguistics, 2019. p. 4171--4186.
Disponível em: \url{https://arxiv.org/abs/1810.04805}.

\bibitem[Dodge et al.(2020)]{dodge2020finetuning}
DODGE, Jesse et al.
Fine-tuning pretrained language models: weight initializations, data orders,
and early stopping.
\textit{arXiv preprint arXiv:2002.06305}, 2020.
Disponível em: \url{https://arxiv.org/abs/2002.06305}.

\bibitem[Efron e Tibshirani(1993)]{efron1993bootstrap}
EFRON, Bradley; TIBSHIRANI, Robert J.
\textit{An introduction to the bootstrap}.
New York: Chapman \& Hall, 1993.
DOI: \url{https://doi.org/10.1007/978-1-4899-4541-9}.

\bibitem[Google BERT(2026)]{google2026mbert}
GOOGLE BERT.
\textit{BERT multilingual base model (cased): model card}.
Hugging Face, 2026.
Disponível em:
\url{https://huggingface.co/google-bert/bert-base-multilingual-cased}.
Acesso em: 1 set. 2026.

\bibitem[Guo et al.(2017)]{guo2017calibration}
GUO, Chuan; PLEISS, Geoff; SUN, Yu; WEINBERGER, Kilian Q.
On calibration of modern neural networks.
In: \textit{Proceedings of the 34th International Conference on Machine Learning}.
Sydney: PMLR, 2017. v. 70, p. 1321--1330.
Disponível em: \url{https://arxiv.org/abs/1706.04599}.

\bibitem[Holm(1979)]{holm1979procedure}
HOLM, Sture.
A simple sequentially rejective multiple test procedure.
\textit{Scandinavian Journal of Statistics}, v. 6, n. 2, p. 65--70, 1979.

\bibitem[Gururangan et al.(2020)]{gururangan2020dontstop}
GURURANGAN, Suchin et al.
Don't stop pretraining: adapt language models to domains and tasks.
In: \textit{Proceedings of the 58th Annual Meeting of the Association for
Computational Linguistics}. Online: Association for Computational Linguistics,
2020. p. 8342--8360.
DOI: \url{https://doi.org/10.18653/v1/2020.acl-main.740}.

\bibitem[Krippendorff(2018)]{krippendorff2018content}
KRIPPENDORFF, Klaus.
\textit{Content analysis: an introduction to its methodology}.
4. ed. Thousand Oaks: SAGE, 2018.

\bibitem[Meta AI(2026)]{meta2026xlmr}
META AI.
\textit{XLM-RoBERTa base: model card}.
Hugging Face, 2026.
Disponível em: \url{https://huggingface.co/FacebookAI/xlm-roberta-base}.
Acesso em: 1 set. 2026.

\bibitem[Mitchell et al.(2019)]{mitchell2019modelcards}
MITCHELL, Margaret et al.
Model cards for model reporting.
In: \textit{Proceedings of the Conference on Fairness, Accountability, and
Transparency}. New York: ACM, 2019. p. 220--229.
DOI: \url{https://doi.org/10.1145/3287560.3287596}.

\bibitem[NeuralMind(2026)]{neuralmind2026bertimbau}
NEURALMIND.
\textit{BERTimbau Base and Large: model cards}.
Hugging Face, 2026.
Disponível em:
\url{https://huggingface.co/neuralmind/bert-base-portuguese-cased} e
\url{https://huggingface.co/neuralmind/bert-large-portuguese-cased}.
Acesso em: 1 set. 2026.

\bibitem[Röttger et al.(2021)]{rottger2021hatecheck}
RÖTTGER, Paul; VIDGEN, Bertie; NGUYEN, Dong; WASEEM, Zeerak; MARGETTS, Helen;
PIERREHUMBERT, Janet.
HateCheck: functional tests for hate speech detection models.
In: \textit{Proceedings of the 59th Annual Meeting of the Association for
Computational Linguistics and the 11th International Joint Conference on
Natural Language Processing}. Online: Association for Computational Linguistics,
2021. p. 41--58.
DOI: \url{https://doi.org/10.18653/v1/2021.acl-long.4}.

\bibitem[Saito e Rehmsmeier(2015)]{saito2015precisionrecall}
SAITO, Takaya; REHMSMEIER, Marc.
The precision-recall plot is more informative than the ROC plot when evaluating
binary classifiers on imbalanced datasets.
\textit{PLOS ONE}, v. 10, n. 3, e0118432, 2015.
DOI: \url{https://doi.org/10.1371/journal.pone.0118432}.

\bibitem[Sechidis, Tsoumakas e Vlahavas(2011)]{sechidis2011stratification}
SECHIDIS, Konstantinos; TSOUMAKAS, Grigorios; VLAHAVAS, Ioannis.
On the stratification of multi-label data.
In: \textit{Machine Learning and Knowledge Discovery in Databases}.
Berlin: Springer, 2011. p. 145--158.
DOI: \url{https://doi.org/10.1007/978-3-642-23808-6_10}.

\bibitem[Souza, Nogueira e Lotufo(2020)]{souza2020bertimbau}
SOUZA, Fábio; NOGUEIRA, Rodrigo; LOTUFO, Roberto.
BERTimbau: pretrained BERT models for Brazilian Portuguese.
In: \textit{Intelligent Systems: 9th Brazilian Conference, BRACIS 2020}.
Cham: Springer, 2020. p. 403--417.
DOI: \url{https://doi.org/10.1007/978-3-030-61377-8_28}.

\bibitem[Wolf et al.(2020)]{wolf2020transformers}
WOLF, Thomas et al.
Transformers: state-of-the-art natural language processing.
In: \textit{Proceedings of the 2020 Conference on Empirical Methods in Natural
Language Processing: System Demonstrations}. Online: Association for
Computational Linguistics, 2020. p. 38--45.
DOI: \url{https://doi.org/10.18653/v1/2020.emnlp-demos.6}.

\end{thebibliography}

%
% Pacotes recomendados no preâmbulo do documento principal:
% \usepackage[T1]{fontenc}
% \usepackage[utf8]{inputenc} % dispensável em distribuições LaTeX recentes
% \usepackage[brazil]{babel}
% \usepackage{amsmath,amssymb}
% \usepackage{booktabs,longtable,array}
% \usepackage{graphicx}
% \usepackage{url}
% \usepackage[hidelinks]{hyperref}
% \usepackage[alf]{abntex2cite}
%
% As referências foram incorporadas ao final deste arquivo para que o capítulo
% não dependa de um arquivo .bib externo. Caso o documento principal já possua
% bibliografia própria, os itens podem ser transferidos para ela sem alterar as
% chaves usadas nas citações.

\chapter{Metodologia de treinamento dos modelos BERT}
\label{cap:metodologia-bert-v01}

\section{Visão geral e objetivos}
\label{sec:visao-geral-bert}

Este capítulo estabelece o protocolo para construção, treinamento, seleção e
avaliação dos primeiros classificadores supervisionados do projeto Radar Ódio
a partir do \textit{Dataset} v0.1. Essa versão será formada pelas avaliações
humanas resolvidas que correspondam a, no mínimo, 10\% do universo elegível. O
objetivo não é tratar a v0.1 como versão definitiva do conjunto de dados, mas
como o primeiro marco reprodutível de um processo contínuo de aprendizagem e
avaliação, a ser ampliado nas versões posteriores.

A estratégia emprega modelos pré-treinados da família BERT e derivados. O uso
de representações contextuais pré-treinadas permite adaptar um modelo de língua
a uma tarefa supervisionada mediante a adição de uma camada de classificação e
o ajuste conjunto dos parâmetros, procedimento conhecido como \textit{fine-tuning}
\cite{devlin2019bert}. Para o português brasileiro, os modelos BERTimbau
constituem referências monolíngues relevantes, tendo sido pré-treinados em um
grande corpus nacional e avaliados em diferentes tarefas de processamento de
linguagem natural \cite{souza2020bertimbau}. Modelos multilíngues serão incluídos
como controles comparativos, permitindo verificar se a especialização no idioma
produz vantagem no domínio de textos oriundos de redes sociais.

O protocolo contempla duas tarefas relacionadas, porém treinadas e avaliadas
separadamente:

\begin{enumerate}
    \item \textbf{detecção binária de discurso de ódio}: classificação de cada
    texto como \texttt{hate} ou \texttt{nao\_hate};
    \item \textbf{classificação multirrótulo dos tipos de ódio}: atribuição de
    zero ou mais categorias temáticas aos registros cujo rótulo gold seja
    \texttt{hate}.
\end{enumerate}

Essa separação define a arquitetura hierárquica candidata: o primeiro modelo
decide se o conteúdo contém discurso de ódio; quando a decisão é positiva, o
segundo identifica os tipos aplicáveis. Sua adoção não será tratada como uma
vantagem presumida. Ela será comparada a uma arquitetura conjunta e terá a
propagação de erros quantificada, conforme a Seção~\ref{sec:arquiteturas}.

\section{Perguntas de pesquisa, hipóteses e contribuição}
\label{sec:perguntas-hipoteses}

Os modelos são empregados como instrumentos para investigar generalização,
desbalanceamento, calibração e propagação de erros, e não apenas como alternativas
de implementação. As perguntas de pesquisa são:

\begin{description}
    \item[RQ1:] um encoder monolíngue pré-treinado em português brasileiro
    apresenta desempenho superior a encoders multilíngues comparáveis na
    detecção de discurso de ódio?
    \item[RQ2:] a ponderação da função de perda melhora a identificação da
    classe minoritária \texttt{hate} e dos tipos menos frequentes?
    \item[RQ3:] a calibração pós-processada reduz o erro das probabilidades
    estimadas sem deteriorar a discriminação?
    \item[RQ4:] o desempenho varia entre plataformas e se mantém quando uma
    plataforma inteira é excluída do treinamento?
    \item[RQ5:] qual é a perda de desempenho causada pela propagação de erros
    entre a detecção binária e a classificação dos tipos?
    \item[RQ6:] o desempenho se mantém em dados coletados depois do período
    representado pela v0.1?
    \item[RQ7:] o pré-treinamento específico em textos de redes sociais produz
    vantagem sobre o pré-treinamento geral da web, especialmente no X/Twitter?
\end{description}

As hipóteses e os contrastes primários serão registrados antes da abertura do
teste. Para RQ1, o contraste principal será BERTimbau Base contra mBERT, pois
ambos possuem porte semelhante; BERTimbau Large e XLM-R serão comparações
secundárias. Para RQ7, o contraste secundário pré-especificado será BERTabaporu
Base contra BERTimbau Base, mantendo as mesmas partições e o mesmo protocolo de ajuste. Seja
$M$ a métrica previamente indicada para cada pergunta. Serão
testadas as seguintes hipóteses direcionais:

\begin{align}
    H_{0,1}&: M_{\mathrm{BERTimbau\ Base}}-M_{\mathrm{mBERT}} \leq 0,
    &H_{1,1}&: M_{\mathrm{BERTimbau\ Base}}-M_{\mathrm{mBERT}} > 0;\\
    H_{0,2}&: M_{\mathrm{ponderada}}-M_{\mathrm{comum}} \leq 0,
    &H_{1,2}&: M_{\mathrm{ponderada}}-M_{\mathrm{comum}} > 0;\\
    H_{0,3}&: E_{\mathrm{calibrado}}-E_{\mathrm{bruto}} \geq 0,
    &H_{1,3}&: E_{\mathrm{calibrado}}-E_{\mathrm{bruto}} < 0;\\
    H_{0,5}&: F1_{\mathrm{oracle}}-F1_{\mathrm{ponta\ a\ ponta}} \leq 0,
    &H_{1,5}&: F1_{\mathrm{oracle}}-F1_{\mathrm{ponta\ a\ ponta}} > 0.
\end{align}

Em RQ3, $E$ designa o Brier score, com o ECE como medida complementar. RQ4 e
RQ6 serão avaliadas como hipóteses de heterogeneidade e degradação de domínio,
respectivamente. RQ6 é prospectiva e somente poderá receber resposta
confirmatória quando houver um teste temporal posterior à v0.1.

Além de significância estatística, será adotada como referência uma diferença
mínima cientificamente relevante de $\delta=0{,}02$ em Macro-F1 ou F1 da classe
\texttt{hate}. Esse valor constitui uma decisão operacional prévia, e não uma
constante universal; análises de sensibilidade usarão também $\delta=0{,}01$ e
$\delta=0{,}03$. A conclusão distinguirá: ausência de evidência, diferença
estatisticamente compatível com zero e diferença superior ao limiar de relevância
prática.

\begin{longtable}{p{0.06\textwidth}p{0.25\textwidth}p{0.30\textwidth}p{0.18\textwidth}}
    \caption{Matriz pré-especificada entre pergunta, hipótese, experimento e métrica.}
    \label{tab:matriz-rq} \\
    \toprule
    \textbf{RQ} & \textbf{Contraste} & \textbf{Experimento} & \textbf{Métrica principal} \\
    \midrule
    \endfirsthead
    \toprule
    \textbf{RQ} & \textbf{Contraste} & \textbf{Experimento} & \textbf{Métrica principal} \\
    \midrule
    \endhead
    RQ1 & monolíngue versus multilíngue & BERTimbau Base versus mBERT nas mesmas partições & Macro-F1 \\
    RQ2 & perda ponderada versus comum & ablação pareada por backbone e semente & F1-\texttt{hate}; AUPRC \\
    RQ3 & probabilidade calibrada versus bruta & \textit{temperature scaling} em conjunto próprio & Brier score; ECE \\
    RQ4 & dentro versus fora do domínio & resultados por plataforma e \textit{leave-one-platform-out} & Macro-F1 \\
    RQ5 & \textit{oracle} versus ponta a ponta & execução das duas formas sobre os mesmos exemplos & $\Delta$Macro-F1 \\
    RQ6 & período histórico versus posterior & teste temporal nas versões posteriores & Macro-F1; AUPRC \\
    RQ7 & pré-treinamento social versus web & BERTabaporu Base versus BERTimbau Base, globalmente e por plataforma & Macro-F1 \\
    \bottomrule
\end{longtable}

A contribuição científica esperada é, portanto, a caracterização controlada de
como idioma de pré-treinamento, desenho amostral, desbalanceamento, calibração,
plataforma e organização hierárquica afetam classificadores de discurso de ódio
em português brasileiro. A construção do dataset e a implementação dos modelos
são meios para essa investigação. Resultados confirmatórios serão separados de
análises exploratórias, especialmente para categorias raras e plataformas com
baixa cobertura.

\section{Definição e congelamento do Dataset v0.1}
\label{sec:dataset-v01}

\subsection{Retrato da base em 10 de setembro de 2026}

Na data de referência deste protocolo, o universo continha 212.476 análises com
estado de processamento bem-sucedido. A base gold reunia 25.090 textos canônicos
com decisão humana resolvida, correspondentes a 11,81\% do universo. Desse total,
3.995 registros haviam sido classificados como \texttt{hate} e 21.095 como
\texttt{nao\_hate}. A distribuição evidencia desbalanceamento entre as classes:
aproximadamente 15,92\% dos registros gold eram positivos e 84,08\% eram
negativos.

\begin{table}[htbp]
    \centering
    \caption{Retrato da base gold em 10 de setembro de 2026.}
    \label{tab:retrato-base-gold}
    \begin{tabular}{lrr}
        \toprule
        \textbf{Indicador} & \textbf{Quantidade} & \textbf{Percentual} \\
        \midrule
        Universo elegível & 212.476 & 100,00\% \\
        Decisões gold resolvidas & 25.090 & 11,81\% \\
        Rótulo \texttt{hate} & 3.995 & 15,92\% do gold \\
        Rótulo \texttt{nao\_hate} & 21.095 & 84,08\% do gold \\
        \bottomrule
    \end{tabular}
    \par\smallskip
    {\footnotesize Fonte: elaboração própria com base nas tabelas de avaliação
    do projeto.\par}
\end{table}

As decisões gold estavam distribuídas entre sete plataformas, conforme a
Tabela~\ref{tab:distribuicao-plataforma}. Essa distribuição deverá ser recalculada
no momento do congelamento, pois o acréscimo de avaliações pode alterar tanto as
contagens quanto as proporções.

\begin{table}[htbp]
    \centering
    \caption{Distribuição preliminar das decisões gold por plataforma.}
    \label{tab:distribuicao-plataforma}
    \begin{tabular}{lr}
        \toprule
        \textbf{Plataforma} & \textbf{Quantidade} \\
        \midrule
        YouTube & 7.735 \\
        Facebook & 4.891 \\
        Reddit & 3.101 \\
        WhatsApp & 3.037 \\
        Twitter & 2.895 \\
        Telegram & 2.660 \\
        TikTok & 771 \\
        \midrule
        \textbf{Total} & \textbf{25.090} \\
        \bottomrule
    \end{tabular}
    \par\smallskip
    {\footnotesize Fonte: elaboração própria com base nas tabelas de avaliação
    do projeto.\par}
\end{table}

Entre os textos positivos, uma mesma publicação pode receber mais de um tipo de
ódio. Por isso, a soma das frequências da Tabela~\ref{tab:tipos-odio-preliminar}
é superior ao número total de textos classificados como \texttt{hate}.

\begin{table}[htbp]
    \centering
    \caption{Frequência preliminar dos tipos de ódio na base gold.}
    \label{tab:tipos-odio-preliminar}
    \begin{tabular}{llr}
        \toprule
        \textbf{Código} & \textbf{Categoria} & \textbf{Ocorrências} \\
        \midrule
        \texttt{political} & Ódio político & 1.841 \\
        \texttt{sexism} & Sexismo & 1.020 \\
        \texttt{lgbtphobia} & LGBTfobia & 725 \\
        \texttt{aporophobia} & Classismo & 316 \\
        \texttt{religious\_intolerance} & Intolerância religiosa & 255 \\
        \texttt{racism} & Racismo & 241 \\
        \texttt{xenophobia} & Xenofobia & 141 \\
        \texttt{body\_shame} & Gordofobia & 123 \\
        \texttt{ageism} & Ageísmo & 52 \\
        \texttt{capacitism} & Capacitismo & 39 \\
        \texttt{other} & Categoria histórica inativa & 1 \\
        \bottomrule
    \end{tabular}
    \par\smallskip
    {\footnotesize Fonte: elaboração própria. As frequências são multirrótulo e,
    portanto, não devem ser somadas para estimar o número de textos positivos.\par}
\end{table}

As categorias menos frequentes ainda não oferecem quantidade suficiente para
uma avaliação confirmatória estável. Esse problema não será resolvido por meio
da eliminação automática de categorias ou pela criação de exemplos artificiais
antes da separação dos dados. Em vez disso, todos os rótulos ativos serão
preservados no treinamento e os resultados das classes raras serão explicitamente
marcados como exploratórios e acompanhados de intervalos de confiança.

\subsection{Justificativa do tamanho da v0.1}

O percentual de 10\% constitui um marco operacional para iniciar o ciclo de
treinamento, e não uma justificativa estatística isolada. O retrato efetivamente
alcançado reúne 25.090 registros e 3.995 positivos, com prevalência observada de
15,92\%. Sob amostragem aleatória simples, a margem de erro aproximada de 95\%
para essa prevalência seria de 0,45 ponto percentual. Esse cálculo não será
interpretado como precisão efetiva porque a base foi estratificada, possui pesos
desiguais e contém grupos correlacionados.

Antes do congelamento será produzido um relatório de suficiência com: contagem
esperada e observada de positivos; contagem por categoria, plataforma e tipo de
conteúdo; tamanho efetivo da amostra; largura dos intervalos de confiança
bootstrap; e quantidade de positivos em cada partição. Para pesos $w_i$, o
tamanho efetivo descritivo será calculado por

\begin{equation}
    n_{\mathrm{efetivo}} =
    \frac{\left(\sum_i w_i\right)^2}{\sum_i w_i^2}.
\end{equation}

Uma análise de poder por simulação, baseada apenas na parcela de desenvolvimento
disponível, estimará a probabilidade de detectar uma diferença pareada de
$\delta=0{,}02$ nas comparações primárias, com nível de 5\% e poder-alvo de 80\%.
A simulação preservará prevalência, estratos e grupos e apresentará também a
menor diferença detectável sob o tamanho alcançado. Seus parâmetros e código
serão arquivados; nenhuma observação do teste será usada para ajustar a análise.

O estudo binário será considerado apto à análise confirmatória apenas se o teste
contiver quantidade suficiente de positivos para produzir estimativas estáveis e
se os intervalos alcançarem precisão compatível com as comparações primárias.
Para os tipos de ódio, a classificação confirmatória ou exploratória obedecerá
aos critérios da Seção~\ref{sec:tarefa-multirrotulo}. Assim, ultrapassar o limiar
de 21.248 registros autoriza o início do benchmark, mas não transforma automaticamente
todas as análises por subgrupo em confirmatórias.

\subsection{Regra de congelamento}

Dez por cento de 212.476 registros correspondem a 21.247,6 observações. Assim,
o limiar operacional da v0.1 será de 21.248 decisões gold válidas. Como as
avaliações são organizadas em lotes, o congelamento ocorrerá ao término do
primeiro lote completo que fizer a base atingir ou ultrapassar esse limiar. Não
será realizado corte dentro de um lote já iniciado, evitando uma seleção
oportunista de casos.

O congelamento deverá produzir um instantâneo imutável, acompanhado de:

\begin{itemize}
    \item data e hora da extração;
    \item consulta SQL ou código de exportação versionado;
    \item número total de registros e distribuição dos rótulos;
    \item identificador do último lote incorporado;
    \item hash SHA-256 do arquivo exportado;
    \item versão da taxonomia e do protocolo de avaliação;
    \item histórico de exclusões e correções aplicadas.
\end{itemize}

Em 10 de setembro de 2026 foi produzida uma exportação candidata com 25.090
decisões válidas. O volume já supera o limiar mínimo da v0.1. Entretanto, o lote
24 ainda constava administrativamente como \texttt{active}, embora todos os seus
953 itens possuíssem decisão resolvida. Por isso, o snapshot pode iniciar a
preparação experimental, mas seu congelamento definitivo deverá registrar o
encerramento administrativo do lote e confirmar novamente as contagens e hashes.

Uma vez congelada, a v0.1 não deverá receber alterações silenciosas. Correções
posteriores de rótulos deverão gerar uma nova versão do conjunto de dados, com
registro das diferenças em relação à versão anterior.

O manifesto de congelamento deverá identificar separadamente as versões do
dataset, da taxonomia e do manual de anotação, por exemplo:

\begin{quote}
\texttt{dataset = radarodio-v0.1}\par
\texttt{taxonomy = taxonomy-v1.2}\par
\texttt{annotation-guidelines = annotation-guidelines-v1.2}
\end{quote}

Mudanças posteriores no manual ou na taxonomia não modificarão retroativamente
o significado de decisões produzidas sob uma versão anterior.

\subsection{Qualidade da anotação e formação do gold standard}
\label{sec:qualidade-anotacao}

Não será realizada nova dupla anotação sistemática da v0.1, pois esse procedimento
é operacionalmente inviável no cronograma do projeto. O controle de confiabilidade
utilizará os arquivos históricos de dupla anotação do Reddit como estudo auxiliar,
mantido separado do teste principal. A auditoria disponível identificou 1.644
pares binários válidos, dos quais 1.512 apresentaram consenso e 132 divergência,
correspondendo a 91,97\% de concordância bruta e $\kappa$ de Cohen de aproximadamente
0,39. Os 1.512 casos consensuais somente poderão ser incorporados depois de
vinculação ao corpus, deduplicação e harmonização da taxonomia; divergências sem
adjudicação permanecerão fora do gold.

Antes da anotação definitiva, o projeto deverá arquivar uma ficha metodológica
com o número e o perfil funcional dos avaliadores, treinamento recebido, versão
do manual, exemplos de treinamento, exercício de qualificação, critérios de
aprovação e procedimento para dúvidas. Informações pessoais dos avaliadores não
entrarão no arquivo destinado ao treinamento dos modelos.

Na tarefa binária, a concordância será descrita pelo percentual bruto, pelo
kappa de Cohen e pelo alfa de Krippendorff. Na tarefa multirrótulo serão
calculados o Jaccard e o F1 entre conjuntos, além da concordância e do alfa de
Krippendorff por categoria \cite{cohen1960kappa,krippendorff2018content}. Todos
os coeficientes serão acompanhados do número de
itens, distribuição marginal e intervalo de confiança. A concordância não será
usada isoladamente como prova de validade; desacordos serão analisados por
categoria, plataforma e fenômeno linguístico.

Haverá divergência quando os avaliadores históricos discordarem sobre o rótulo
binário ou sobre ao menos um tipo ativo. Sem nova adjudicação rastreável, esses
casos não serão transformados em rótulos de treinamento. Se uma adjudicação
pontual vier a ocorrer, deverá registrar decisão, justificativa, versão das
diretrizes, data e categorias em disputa, preservando as duas respostas originais.

Na data de referência, todas as 25.090 decisões resolvidas no sistema tinham
estado \texttt{single}. A v0.1 poderá ser usada para desenvolvimento e avaliação
interna, mas a ausência de dupla anotação contemporânea e multiplataforma limita
a força das conclusões sobre confiabilidade do gold. O estudo histórico do Reddit
será reportado como evidência auxiliar, sem generalização automática para as
demais plataformas.

\subsection{Critérios de elegibilidade}

Serão incluídos somente textos canônicos cujo consenso possua um dos estados
\texttt{single}, \texttt{agreed} ou \texttt{adjudicated}. Serão excluídos:

\begin{itemize}
    \item avaliações puladas ou ainda atribuídas;
    \item conflitos que não tenham sido adjudicados;
    \item registros sem texto ou com falha de codificação não recuperável;
    \item registros sem rótulo final consistente;
    \item cópias não canônicas de textos duplicados;
    \item observações que violem as regras vigentes da taxonomia.
\end{itemize}

Na data de referência, todas as 25.090 decisões resolvidas tinham estado
\texttt{single}. Consequentemente, não é possível estimar concordância entre
avaliadores para essa parcela da base. Essa condição será registrada como uma
limitação da v0.1. Os casos históricos de revisão dupla serão mantidos em um
conjunto auxiliar com proveniência própria e não alterarão silenciosamente o
estado das decisões produzidas pelo sistema.

O código histórico \texttt{other}, atualmente inativo, não integrará o vetor de
saída da tarefa multirrótulo. O registro que ainda contém esse código poderá ser
mantido na tarefa binária, mas deverá ser excluído da tarefa de tipos até que
seja revisto e recodificado em uma categoria ativa.

Antes da exportação serão aplicadas as seguintes invariantes:

\begin{align}
    \texttt{nao\_hate} &\Rightarrow
    \texttt{gold\_hate\_types}=\varnothing,\\
    \texttt{hate} &\Rightarrow
    \left|\texttt{gold\_hate\_types}\right|\geq 1,
\end{align}

em que todos os tipos devem pertencer à versão ativa da taxonomia. A condição
\texttt{hate} sem tipo não será silenciosamente convertida em negativa: o caso
será devolvido para adjudicação ou ficará fora da tarefa multirrótulo até ser
resolvido. Também serão verificadas unicidade de \texttt{analysis\_id}, presença
dos campos obrigatórios, reprodução do \texttt{text\_hash}, ausência de tipos
inativos e inexistência de rótulos conflitantes dentro de um grupo de duplicatas.
O relatório dessas validações será armazenado com o snapshot.

\subsection{Contrato de exportação}

A extração para treinamento deverá seguir o contrato apresentado na
Tabela~\ref{tab:contrato-exportacao}. Campos administrativos e identidades dos
avaliadores não serão incorporados ao conjunto utilizado pelos modelos.

\begin{longtable}{p{0.22\textwidth}p{0.16\textwidth}p{0.48\textwidth}}
    \caption{Contrato mínimo de exportação do Dataset v0.1.}
    \label{tab:contrato-exportacao} \\
    \toprule
    \textbf{Campo} & \textbf{Tipo} & \textbf{Finalidade} \\
    \midrule
    \endfirsthead
    \toprule
    \textbf{Campo} & \textbf{Tipo} & \textbf{Finalidade} \\
    \midrule
    \endhead
    \texttt{analysis\_id} & inteiro & Identificador interno e rastreabilidade da análise. \\
    \texttt{text\_hash} & texto & Identificação de duplicatas exatas e definição dos grupos de partição. \\
    \texttt{text} & texto & Conteúdo anonimizado que será fornecido ao tokenizer. \\
    \texttt{platform} & categoria & Auditoria e avaliação estratificada por plataforma. \\
    \texttt{content\_kind} & categoria & Identificação de post ou comentário. \\
    \texttt{gold\_label} & categoria & Alvo binário: \texttt{hate} ou \texttt{nao\_hate}. \\
    \texttt{gold\_hate\_types} & lista & Conjunto dos tipos de ódio aplicáveis ao texto. \\
    \texttt{decision\_status} & categoria & Indicação de decisão simples, concordante ou adjudicada. \\
    \texttt{batch\_id} & inteiro & Identificação do lote de amostragem. \\
    \texttt{stratum} & texto & Estrato usado na seleção e no cálculo de pesos amostrais. \\
    \texttt{inclusion\_probability} & numérico & Probabilidade de inclusão, quando conhecida e auditável. \\
    \texttt{duplicate\_group\_id} & texto & Grupo transitivo de duplicatas e quase duplicatas. \\
    \texttt{context\_group\_id} & texto/nulo & Conversa, tópico ou publicação de origem compartilhada. \\
    \texttt{content\_created\_at} & data/hora/nulo & Data do conteúdo para avaliação temporal. \\
    \texttt{reviewed\_at} & data/hora & Controle da versão e avaliações temporais futuras. \\
    \texttt{taxonomy\_version} & texto & Versão da taxonomia aplicada ao rótulo. \\
    \texttt{guidelines\_version} & texto & Versão do manual de anotação. \\
    \texttt{context\_dependency} & categoria/nulo & Indício de dependência de imagem, vídeo, conversa ou evento externo. \\
    \bottomrule
\end{longtable}

Os campos \texttt{gemma\_label}, \texttt{gemma\_hate\_probability} e
\texttt{gemma\_hate\_types} poderão ser exportados em uma tabela auxiliar
restrita à auditoria, ao cálculo dos estratos e à reprodução do baseline. Eles
não poderão ser concatenados ao texto nem empregados como atributos de entrada
dos classificadores BERT. Essa separação impede que os novos modelos apenas
reproduzam os sinais do sistema que orientou a amostragem.

\section{Preparação dos textos}
\label{sec:preparacao-textos}

O pré-processamento será deliberadamente conservador. Modelos BERT utilizam
tokenização em subpalavras e foram concebidos para aprender representações a
partir de sequências próximas da linguagem observada. Remoções agressivas podem
eliminar sinais relevantes para discurso de ódio, especialmente pontuação,
emojis, hashtags, caixa, erros ortográficos e alterações intencionais de
grafia.

Cada texto passará, nesta ordem, pelas seguintes operações:

\begin{enumerate}
    \item conversão para Unicode e normalização na forma NFC;
    \item substituição de sequências inválidas de codificação quando a correção
    for inequívoca;
    \item padronização de quebras de linha e espaços consecutivos;
    \item substituição de URLs pelo marcador \texttt{<URL>};
    \item substituição de nomes de usuário e menções pelo marcador
    \texttt{<USER>};
    \item substituição de endereços de e-mail por \texttt{<EMAIL>};
    \item substituição de números de telefone por \texttt{<PHONE>};
    \item remoção de espaços nas extremidades, sem alterar o restante do texto.
\end{enumerate}

Os marcadores serão adicionados como tokens especiais indivisíveis ao tokenizer
de cada checkpoint, com redimensionamento correspondente da matriz de embeddings.
As novas linhas serão inicializadas pelo procedimento padrão definido na
configuração do checkpoint, treinadas durante o fine-tuning e controladas pelas
mesmas sementes das demais inicializações. A lista de tokens, os IDs resultantes
e o hash do tokenizer alterado serão registrados. Uma ablação no conjunto de
desenvolvimento comparará o mascaramento com a manutenção das formas originais,
após anonimização irreversível dos identificadores, para medir eventual perda de
sinais úteis.

Não serão aplicados \textit{stemming}, lematização, remoção de \textit{stopwords},
tradução, conversão geral para minúsculas ou correção ortográfica. Os modelos
sensíveis à caixa receberão a capitalização original. Expressões regulares,
dicionários e versões de qualquer reconhecedor de informação pessoal serão
versionados. Nomes de grupos sociais necessários à interpretação do alvo não
serão removidos como se fossem identificadores pessoais. O texto bruto ficará
em ambiente de acesso restrito; somente a versão anonimizada poderá integrar os
artefatos de treinamento.

O BERTabaporu é um modelo \textit{uncased}; portanto, sua normalização de caixa
será realizada pelo tokenizer do próprio checkpoint. O arquivo de entrada comum
continuará preservando a capitalização original, evitando que a inclusão desse
modelo modifique os dados recebidos pelos encoders sensíveis à caixa.

\subsection{Duplicatas e prevenção de vazamento}

Primeiramente, o mapa de duplicatas do sistema será usado para manter uma única
linha por \texttt{text\_hash} canônico. Em seguida, será criada uma representação
normalizada em 5-gramas de caracteres. Pares com similaridade de Jaccard maior
ou igual a 0,90 serão tratados como quase duplicatas para fins de partição. A
detecção de quase duplicatas não implicará exclusão automática: os registros
serão reunidos sob um mesmo identificador de grupo, garantindo que permaneçam
na mesma partição e no mesmo fold.

O limiar 0,90 será validado antes da divisão definitiva. Serão sorteados pares
dos intervalos 0,80--0,89, 0,90--0,94 e 0,95--1,00 para julgamento humano cego,
registrando precisão e revocação aproximadas do detector de quase duplicatas.
Se a inspeção indicar separação inadequada, qualquer alteração do limiar será
feita e documentada antes da abertura do teste.

O agrupamento será transitivo. Os textos serão vértices de um grafo; cada par
que satisfizer o limiar formará uma aresta; e cada componente conexo receberá
um único \texttt{duplicate\_group\_id}. Assim, se $A$ for semelhante a $B$ e
$B$ a $C$, os três permanecerão juntos mesmo quando a similaridade direta entre
$A$ e $C$ for inferior a 0,90.

Quando houver identificador confiável de conversa, tópico ou encadeamento, todos
os textos do mesmo contexto também serão mantidos no mesmo grupo. Considera-se
contexto compartilhado a mesma postagem original e seus comentários, a mesma
thread, o mesmo encadeamento de respostas ou uma conversa identificada pelo
mesmo ID de origem. Um \texttt{context\_group\_id} prevalecerá sobre a decisão
individual de partição. Essa regra reduz o risco de um comentário no treinamento
revelar o conteúdo de outro comentário relacionado presente no teste.

\subsection{Comprimento das sequências}

Antes de definir o truncamento, será calculada, separadamente para cada tokenizer,
a distribuição do comprimento dos textos. Serão reportados média, mediana,
percentis 90, 95 e 99, máximo e as proporções $P(L>128)$, $P(L>256)$ e
$P(L>512)$. O limite técnico superior será de 512 tokens, incluindo os tokens
especiais do encoder.

Três políticas serão comparadas no desenvolvimento: (A) primeiros 510 tokens de
conteúdo; (B) 255 tokens iniciais e 255 finais; e (C) janelas de até 510 tokens
com sobreposição, agregadas no nível do documento. Para a política C, tamanho da
janela, sobreposição e regra de agregação serão fixados antes do teste. Se menos
de 5\% dos textos excederem 512 tokens e não houver diferença relevante no
desenvolvimento, a política B será adotada por parcimônia. Caso contrário, a
política será tratada como hiperparâmetro de ablação. A quantidade de textos
truncados, o conteúdo médio descartado e o custo adicional do \textit{chunking}
serão registrados.

\section{Definição das tarefas de aprendizagem}
\label{sec:tarefas-aprendizagem}

\subsection{Tarefa A: detecção binária}

Para um texto $x_i$, o alvo binário será definido por

\begin{equation}
    y_i =
    \begin{cases}
        1, & \text{se o rótulo gold for \texttt{hate}},\\
        0, & \text{se o rótulo gold for \texttt{nao\_hate}}.
    \end{cases}
\end{equation}

O vetor contextual associado ao token de classificação será fornecido a uma
camada linear com duas saídas, seguida por \textit{softmax}. A função de perda
será a entropia cruzada. Serão comparadas duas condições: perda sem ponderação e
perda ponderada por classe. Na condição ponderada, o peso da classe $c$ será
calculado apenas no subconjunto de treinamento de cada fold:

\begin{equation}
    \alpha_c = \frac{N_{\mathrm{treino}}}{2N_c},
\end{equation}

em que $N_c$ é o número de exemplos da classe $c$ no fold de treinamento.

\subsection{Tarefa B: classificação multirrótulo}
\label{sec:tarefa-multirrotulo}

A tarefa multirrótulo utilizará somente os textos com rótulo gold
\texttt{hate}. Para cada texto, será construído um vetor binário
$\mathbf{y}_i \in \{0,1\}^{10}$, na seguinte ordem fixa:

\begin{enumerate}
    \item \texttt{ageism};
    \item \texttt{aporophobia};
    \item \texttt{body\_shame};
    \item \texttt{capacitism};
    \item \texttt{lgbtphobia};
    \item \texttt{political};
    \item \texttt{racism};
    \item \texttt{religious\_intolerance};
    \item \texttt{sexism};
    \item \texttt{xenophobia}.
\end{enumerate}

No banco de origem, o rótulo de sexismo ainda está armazenado sob o código
histórico \texttt{misogyny}. Durante a preparação da tarefa multirrótulo, esse
valor será mapeado deterministicamente para \texttt{sexism}, sem alterar o
arquivo bruto. A tabela de mapeamento e sua versão integrarão o manifesto do
dataset.

O cabeçote terá dez saídas independentes com ativação sigmoide. A perda será a
entropia cruzada binária com logits, implementada por
\texttt{BCEWithLogitsLoss}. Também serão comparadas uma condição sem ponderação
e outra com peso positivo por categoria. Para a categoria $\ell$, o peso será
calculado exclusivamente no fold de treinamento:

\begin{equation}
    p_{\ell} = \frac{N_{\ell}^{-}}{N_{\ell}^{+}},
\end{equation}

em que $N_{\ell}^{+}$ e $N_{\ell}^{-}$ são, respectivamente, as quantidades de
exemplos positivos e negativos para o rótulo $\ell$.

Uma categoria será considerada elegível para inferência confirmatória somente
se possuir pelo menos 100 ocorrências na v0.1 e pelo menos 20 ocorrências no
conjunto de teste. Esses valores são salvaguardas operacionais mínimas para
evitar folds praticamente vazios; não garantem, sozinhos, precisão inferencial.
Por isso, serão apresentados intervalos de confiança e uma análise de
sensibilidade exigindo 20, 30 e 50 positivos no teste. Categorias cujo intervalo
de 95\% para F1 permaneça excessivamente amplo serão descritas como exploratórias,
mesmo que ultrapassem o limite mínimo.

\section{Arquiteturas e propagação de erros}
\label{sec:arquiteturas}

O experimento principal comparará duas organizações sob o mesmo encoder,
partições, pré-processamento e orçamento de ajuste:

\begin{enumerate}
    \item \textbf{hierárquica}: um classificador binário é executado primeiro e
    o classificador multirrótulo somente produz tipos quando o primeiro estágio
    prevê \texttt{hate};
    \item \textbf{conjunta não hierárquica}: um único encoder possui onze saídas
    sigmoides, uma para \texttt{hate} e dez para os tipos. Registros
    \texttt{nao\_hate} recebem zero nos dez tipos, e nenhuma regra de bloqueio é
    aplicada antes da medição. Inconsistências, como prever um tipo sem prever
    \texttt{hate}, serão contabilizadas.
\end{enumerate}

Na arquitetura conjunta, a perda será a soma ponderada da perda binária e da
perda dos tipos,

\begin{equation}
    \mathcal{L}_{\mathrm{conjunta}} =
    \mathcal{L}_{\mathrm{binária}} +
    \gamma\mathcal{L}_{\mathrm{tipos}},
\end{equation}

com $\gamma$ definido exclusivamente no desenvolvimento. Para limitar o custo,
essa ablação será executada inicialmente com o backbone vencedor da tarefa
binária; uma extensão aos demais backbones será identificada como análise
secundária.

O classificador de tipos será avaliado tanto com a entrada binária gold
(condição \textit{oracle}) quanto no pipeline completo. A propagação de erros
será quantificada por

\begin{equation}
    \Delta F1_{\mathrm{propagação}} =
    F1_{\mathrm{oracle}}-F1_{\mathrm{ponta\ a\ ponta}},
\end{equation}

globalmente, por categoria e por plataforma. A comparação hierárquica versus
conjunta utilizará predições pareadas no mesmo teste; qualquer vantagem deverá
ser interpretada junto do custo computacional e da taxa de inconsistências.

\section{Particionamento e desenho de validação}
\label{sec:particionamento}

\subsection{Teste fixo e conjunto de desenvolvimento}

Após a formação dos grupos de duplicatas, 20\% da v0.1 serão reservados como
conjunto de teste fixo. Dos 80\% restantes, 12,5\% serão separados previamente
como conjunto de calibração, equivalentes a 10\% da base completa. Os 70\%
restantes formarão o conjunto de seleção e ajuste. A proporção global será,
portanto, 70\% para seleção/treinamento, 10\% para calibração e limiares e 20\%
para teste. Todas as divisões respeitarão os grupos e a estratificação.

O conjunto de calibração terá três usos previamente limitados: estimar o
parâmetro de calibração, definir os limiares da configuração já selecionada e
medir a estabilidade dessas escolhas. Ele não será usado para escolher backbone,
taxa de aprendizado, perda, época ou política de truncamento. O teste será
armazenado separadamente e não poderá participar de:

\begin{itemize}
    \item ajuste de hiperparâmetros;
    \item seleção de modelo ou de época;
    \item definição de pesos de classe;
    \item escolha de limiares de decisão;
    \item calibração de probabilidades ou escolha de seus métodos;
    \item análise exploratória destinada a alterar o pipeline.
\end{itemize}

A divisão será realizada com estratificação iterativa multirrótulo, abordagem
destinada a conservar distribuições de rótulos em problemas nos quais uma mesma
instância pode pertencer a várias classes \cite{sechidis2011stratification}. O
vetor de estratificação combinará o rótulo binário, os dez tipos de ódio, a
plataforma, o tipo de conteúdo e o estrato de amostragem. A unidade indivisível
será o grupo de duplicatas, e não o registro individual.

\subsection{Validação cruzada}

O conjunto de seleção e ajuste será dividido em cinco folds por meio da mesma
estratégia de agrupamento e estratificação. Em cada rodada, quatro folds serão
usados para treinamento e um para validação. Ao final das cinco rodadas, cada
registro desse conjunto terá exatamente uma predição fora da amostra,
denominada \textit{out-of-fold} (OOF).

As predições OOF serão usadas para comparar configurações e estimar sua
variabilidade sem consultar o conjunto de calibração. A métrica OOF global,
calculada após concatenar as cinco predições de validação, será o critério
principal de seleção. A média e o desvio-padrão das cinco métricas calculadas
separadamente também serão reportados, mas não substituirão o critério global.
Essa escolha é especialmente importante porque o
fine-tuning de modelos pré-treinados pode apresentar variação relevante em
função da inicialização e da ordem dos dados, sobretudo em conjuntos menores
\cite{dodge2020finetuning}.

\subsection{Correção do desenho amostral}

Os lotes de avaliação foram balanceados a partir de informações como plataforma,
rótulo previsto pelo Gemma e tipo previsto de ódio. Logo, a distribuição da
base anotada pode diferir da distribuição do universo completo. O manifesto de
amostragem registrará o universo disponível em cada lote, definição e tamanho
dos estratos, quantidade selecionada, algoritmo, semente, critérios de exclusão,
uso das previsões do Gemma e qualquer priorização manual. Para cada estrato $h$,
será calculado o peso de expansão

\begin{equation}
    w_h = \frac{N_h}{n_h},
\end{equation}

em que $N_h$ é o número de registros elegíveis do estrato no universo e $n_h$
é o número de registros selecionados do mesmo estrato. A expressão corresponde
ao inverso da probabilidade de inclusão somente quando a seleção é conhecida e
equiprovável dentro de cada estrato. Para estabilidade numérica, os pesos serão
normalizados para média 1 dentro de cada partição.

As métricas do teste serão apresentadas em duas versões: não ponderada, que
descreve diretamente o conjunto anotado, e ponderada, que busca estimar o
desempenho sob a composição do universo original. Essa segunda interpretação
depende das hipóteses de inclusão conhecível, positividade e seleção ignorável
dentro dos estratos. Seleção manual, revisão dirigida ou exclusão subjetiva
viola essas hipóteses e não é corrigida automaticamente por $w_h$; nessa situação,
a versão ponderada será qualificada como ajuste descritivo, não como estimativa
representativa da população. Os pesos não serão calculados com base no rótulo
gold do teste e jamais serão derivados de seus erros.

Sempre que operacionalmente possível, será sorteada do universo uma amostra
probabilística independente das previsões do Gemma. Ela será anotada pelo mesmo
protocolo e usada como teste externo da prevalência, dos pesos e da generalização.
Se essa amostra não existir na v0.1, a ausência será registrada como limitação e
sua construção será prioridade da versão seguinte.

\section{Modelos comparados}
\label{sec:modelos-comparados}

Cinco encoders serão submetidos ao mesmo desenho experimental:

\begin{description}
    \item[BERTimbau Base:] \texttt{neuralmind/bert-base-portuguese-cased},
    modelo BERT Base monolíngue, sensível à caixa e pré-treinado para português
    brasileiro \cite{souza2020bertimbau,neuralmind2026bertimbau};
    \item[BERTimbau Large:] \texttt{neuralmind/bert-large-portuguese-cased},
    versão de maior capacidade do mesmo projeto
    \cite{souza2020bertimbau,neuralmind2026bertimbau};
    \item[BERTabaporu Base:]
    \texttt{pablocosta/bertabaporu-base-uncased}, modelo BERT Base monolíngue
    sem distinção de caixa, pré-treinado em aproximadamente 238 milhões de
    tweets em português brasileiro, provenientes de mais de 100 mil usuários e
    totalizando cerca de 2,9 bilhões de tokens
    \cite{dacosta2023bertabaporu,dacosta2026bertabaporu};
    \item[mBERT:] \texttt{google-bert/bert-base-multilingual-cased}, modelo BERT
    multilíngue sensível à caixa, pré-treinado em 104 idiomas
    \cite{devlin2019bert,google2026mbert};
    \item[XLM-R Base:] \texttt{FacebookAI/xlm-roberta-base}, encoder
    multilíngue baseado em RoBERTa, incluído como comparador da família de
    Transformers bidirecionais derivados de BERT
    \cite{conneau2020xlmr,meta2026xlmr}.
\end{description}

O BERTabaporu permitirá testar diretamente o efeito do gênero de pré-treinamento:
BERTimbau Base representa o português geral da web, enquanto BERTabaporu Base
representa a linguagem de redes sociais. Como a v0.1 é multiplataforma, eventual
vantagem no X/Twitter não será generalizada para as demais fontes. Serão
reportados o resultado global, o resultado por plataforma e a interação entre
modelo e plataforma. A versão
\texttt{pablocosta/bertabaporu-large-uncased} ficará como extensão secundária,
condicionada ao orçamento computacional, e não substituirá a execução do Base
no benchmark principal.

Os checkpoints e tokenizers serão carregados por meio da biblioteca
Transformers \cite{wolf2020transformers}. A revisão exata de cada repositório
será fixada por hash de commit no início dos experimentos, evitando que uma
atualização remota altere silenciosamente os resultados.

\subsection{Baselines}

Três famílias de baselines serão mantidas:

\begin{enumerate}
    \item \textbf{baseline trivial}: previsão constante de \texttt{nao\_hate}
    na tarefa binária e previsão determinada apenas pela prevalência de cada
    categoria, estimada no treinamento, na tarefa multirrótulo;
    \item \textbf{TF--IDF com regressão logística}: concatenação de n-gramas de
    palavras de ordem 1--2 e n-gramas de caracteres de ordem 3--5, seguida por
    regressão logística. O parâmetro $C$ será selecionado entre 0,1, 1 e 10 no
    mesmo protocolo de cinco folds. Na tarefa multirrótulo será usado
    \texttt{OneVsRestClassifier}, com um classificador logístico independente
    para cada um dos dez tipos;
    \item \textbf{Gemma}: uso das predições originais registradas antes da
    revisão humana, sem novo ajuste. Esse baseline permitirá comparar os BERTs
    com o sistema responsável pelas classificações iniciais.
\end{enumerate}

O baseline TF--IDF será submetido às mesmas partições, agrupamentos e métricas.
Nenhum vocabulário ou transformador TF--IDF será ajustado fora do fold de
treinamento. Como as previsões do Gemma ajudaram a definir estratos e prioridades
de revisão, o desempenho desse baseline não será interpretado como proveniente
de uma amostra independente do próprio sistema. A possível dependência entre a
amostragem e seus erros será explicitada em todas as comparações.

\section{Protocolo de fine-tuning}
\label{sec:fine-tuning}

Todos os encoders receberão cabeçotes de classificação novos, inicializados
aleatoriamente. Os parâmetros do encoder e do cabeçote serão ajustados em
conjunto. A faixa inicial segue a prática de fine-tuning apresentada para BERT,
enquanto decisões adicionais serão verificadas em piloto no próprio domínio
\cite{devlin2019bert}. O protocolo comum será:

\begin{itemize}
    \item otimizador AdamW, com $\beta_1=0{,}9$, $\beta_2=0{,}999$ e
    $\epsilon=10^{-8}$;
    \item tamanho efetivo de lote igual a 32, obtido por acúmulo de gradientes
    quando necessário;
    \item limite superior de oito épocas, entendido como teto operacional e não
    como duração fixa;
    \item avaliação ao final de cada época e parada antecipada após duas épocas
    completas sem redução da perda de validação;
    \item seleção do checkpoint de menor perda de validação, antes da otimização
    de qualquer limiar;
    \item aquecimento linear durante os 10\% primeiros passos;
    \item decaimento linear da taxa de aprendizado após o aquecimento;
    \item norma máxima do gradiente igual a 1,0;
    \item precisão mista FP16 ou BF16 quando suportada pelo hardware;
    \item embaralhamento reprodutível dos exemplos a cada época.
\end{itemize}

A perda de validação foi escolhida para a parada antecipada porque independe de
um limiar de classificação; Macro-F1, F1-\texttt{hate} e AUPRC serão monitorados,
mas não mudarão o checkpoint durante essa etapa. O teto de oito épocas apenas
permite que a regra de paciência opere em tarefas multirrótulo de convergência
mais lenta. O artigo original de BERT usou buscas mais curtas, tipicamente entre
duas e quatro épocas \cite{devlin2019bert}; portanto, será realizado um piloto
com curvas de treinamento. Se as curvas mostrarem convergência sistemática antes
do teto, o custo será reduzido pela parada antecipada. Se permanecerem em melhora
na oitava época, qualquer ampliação exigirá emenda prévia ao protocolo.

Será avaliada a grade cartesiana formada por

\begin{equation}
    \eta \in \{10^{-5}, 2\times10^{-5}, 3\times10^{-5}, 5\times10^{-5}\}
\end{equation}

e

\begin{equation}
    \lambda \in \{0, 0{,}01\},
\end{equation}

em que $\eta$ representa a taxa de aprendizado e $\lambda$ o decaimento de
peso. O valor $10^{-5}$ foi acrescentado às taxas usuais de BERT para contemplar
ajustes mais conservadores, especialmente no modelo Large. A condição de perda
comum ou ponderada será um fator experimental controlado, e não uma decisão
tomada depois de observar o teste.

\subsection{Busca de hiperparâmetros em duas etapas}

Uma grade completa com cinco modelos, duas tarefas, oito combinações, cinco
folds e duas perdas produziria 800 fine-tunings antes das repetições por semente.
Para reduzir custo sem escolher configurações pelo teste, será usada a seguinte
triagem pré-especificada:

\begin{enumerate}
    \item para cada combinação de backbone, tarefa e condição de perda, as oito
    combinações de $\eta$ e $\lambda$ serão executadas por até duas épocas em uma
    divisão interna fixa do conjunto de seleção, com semente 42;
    \item as duas melhores combinações de cada célula serão continuadas até o
    teto de oito épocas, sujeitas à mesma parada antecipada;
    \item a melhor combinação de cada célula seguirá para os cinco folds. Se as
    duas primeiras diferirem menos de 0,005 em Macro-F1, ambas seguirão para
    evitar descarte baseado em ruído da triagem;
    \item a configuração será selecionada pela métrica OOF global, obtida pela
    concatenação das predições dos cinco folds com limiar 0,5. Somente depois
    dessa seleção serão definidos calibração e limiares.
\end{enumerate}

Sem empates, a triagem compreende 160 execuções curtas, 40 continuações e 100
ajustes em folds, em lugar de 800 ajustes completos. Antes da execução será
estimado o orçamento em GPU-horas a partir do piloto, incluindo as repetições
finais e a ablação de arquitetura. Se o orçamento disponível for insuficiente,
o escopo será reduzido por regra registrada antes do teste, priorizando as
perguntas de pesquisa, e não por resultados parciais convenientes.

Na tarefa binária, Macro-F1 será o critério de seleção; F1-\texttt{hate} e AUPRC
serão critérios secundários. Na tarefa multirrótulo, será usado Macro-F1 das
categorias previamente elegíveis. Empates inferiores a 0,005 serão resolvidos,
nesta ordem, por menor variância entre folds, menor perda OOF e menor taxa de
aprendizado. Número de parâmetros e tempo de inferência não serão usados para
desempatar taxas do mesmo backbone, pois permanecem essencialmente constantes.

\subsection{Repetições e definição do modelo final}

Após a seleção, cada condição experimental será treinada no conjunto de seleção
e ajuste com as sementes 13, 42, 97, 123 e 2026. O número de épocas será fixado
na mediana das melhores épocas dos cinco folds, evitando parada antecipada no
teste ou no conjunto de calibração. As probabilidades dos cinco modelos serão
preservadas para medir variabilidade de otimização.

Para a avaliação principal de cada condição, o preditor final será o ensemble
pré-especificado das cinco sementes, formado pela média dos logits antes da
calibração. Nenhuma semente será escolhida pelo desempenho no teste. Caso seja
necessário disponibilizar também um único checkpoint por restrição operacional,
a semente 42 será usada por definição prévia e seus resultados serão reportados
separadamente dos resultados do ensemble.

\section{Limiarização e calibração}
\label{sec:limiarizacao-calibracao}

Na tarefa binária, o ensemble produzirá um logit médio e uma probabilidade
estimada de \texttt{hate}. Depois de fixados backbone, hiperparâmetros, perda,
época e sementes, o limiar $\tau$ será escolhido exclusivamente no conjunto de
calibração. Serão testados valores entre 0,01 e 0,99 em incrementos de 0,01,
selecionando-se o valor que maximizar o Macro-F1. Em empate, será escolhido o
limiar mais próximo de 0,5 e, persistindo o empate, o maior limiar. O valor 0,5
também será reportado como referência.

Na tarefa multirrótulo, cada categoria $\ell$ terá um limiar próprio
$\tau_{\ell}$, escolhido na mesma grade e no mesmo conjunto para maximizar o F1
da categoria. Se houver menos de cinco exemplos positivos de uma categoria na
calibração, seu limiar será mantido em 0,5 e o resultado será tratado como
exploratório. A estabilidade será descrita reestimando os limiares em
reamostragens agrupadas da calibração e, apenas como diagnóstico, nas validações
dos cinco folds; serão reportados mediana, intervalo interquartil e amplitude.

As probabilidades da tarefa binária serão adicionalmente calibradas por
\textit{temperature scaling} ajustado nos logits médios das cinco sementes no
conjunto reservado de calibração. Essa separação evita transferir diretamente
um calibrador estimado sobre logits OOF, produzidos por modelos distintos, para
um modelo final reajustado. Redes neurais modernas podem produzir probabilidades
excessivamente confiantes, e o método oferece uma correção pós-processada simples
\cite{guo2017calibration}.

A temperatura não ponderada estimará calibração para a distribuição da amostra
anotada. Uma temperatura ponderada, obtida pela minimização da log-verossimilhança
negativa com pesos de inclusão, será estimada para a população somente se a
auditoria do desenho amostral validar esses pesos. Caso contrário, não será feita
alegação de que $P(\texttt{hate}\mid\text{amostra})$ corresponde a
$P(\texttt{hate}\mid\text{população})$. O teste receberá apenas transformações e
limiares já congelados. ECE será calculado com 15 intervalos de frequência
aproximadamente igual; também serão apresentados Brier score e diagramas de
confiabilidade antes e depois da calibração.

\section{Métricas de avaliação}
\label{sec:metricas-avaliacao}

\subsection{Tarefa binária}

O Macro-F1 será a métrica primária de comparação geral. F1-\texttt{hate} e AUPRC
serão métricas centrais para as hipóteses sobre a classe positiva, mas não
substituirão retroativamente a métrica definida para cada contraste na
Tabela~\ref{tab:matriz-rq}. Para duas classes $C$, Macro-F1 será calculado
por

\begin{equation}
    \operatorname{MacroF1} = \frac{1}{|C|}\sum_{c\in C}
    \frac{2P_cR_c}{P_c+R_c},
\end{equation}

em que $P_c$ e $R_c$ representam precisão e revocação da classe $c$. O uso de
Macro-F1 impede que o desempenho da classe majoritária domine sozinho a
comparação.

Também serão reportados:

\begin{itemize}
    \item precisão, revocação e F1 da classe \texttt{hate};
    \item área sob a curva Precisão--Revocação (AUPRC);
    \item área sob a curva ROC (AUROC);
    \item coeficiente de correlação de Matthews (MCC);
    \item acurácia balanceada;
    \item matriz de confusão;
    \item Brier score;
    \item erro esperado de calibração (ECE);
    \item curvas de confiabilidade antes e depois da calibração.
\end{itemize}

A AUPRC receberá destaque porque a classe positiva é minoritária. Em cenários
desbalanceados, curvas Precisão--Revocação são mais informativas sobre a
qualidade das predições positivas do que a curva ROC isoladamente
\cite{saito2015precisionrecall}.

\subsection{Tarefa multirrótulo}

A métrica primária será o Macro-F1 calculado sobre as categorias elegíveis para
análise confirmatória. Serão apresentados, adicionalmente:

\begin{itemize}
    \item Macro-F1 incluindo todos os rótulos ativos;
    \item Micro-F1;
    \item precisão, revocação, F1 e PR-AUC por categoria;
    \item \textit{Hamming loss};
    \item \textit{subset accuracy}, que exige correspondência exata do conjunto
    de rótulos;
    \item média da quantidade de rótulos previstos por texto;
    \item matriz de coocorrência entre tipos verdadeiros e previstos.
\end{itemize}

Serão realizadas duas avaliações complementares. A avaliação \textit{oracle}
aplicará o classificador de tipos somente aos registros cujo rótulo binário gold
seja \texttt{hate}, medindo a capacidade específica de reconhecer categorias.
A avaliação ponta a ponta executará primeiro o classificador binário e somente
atribuirá tipos aos textos previstos como positivos. A diferença entre os dois
resultados quantificará a propagação de erros da primeira etapa.

\subsection{Incerteza e comparação entre modelos}

Serão separadas duas fontes de incerteza. A variabilidade de otimização será
descrita pela média, desvio-padrão, mínimo e máximo entre as cinco sementes. A
incerteza amostral será estimada sobre a predição final pré-especificada por
2.000 reamostragens bootstrap do conjunto de teste \cite{efron1993bootstrap}. O
bootstrap será estratificado
pelos estratos de amostragem e terá como unidade o maior agrupamento entre
\texttt{duplicate\_group\_id} e \texttt{context\_group\_id}, preservando a
dependência entre textos relacionados. Na versão ponderada, os pesos serão
mantidos em cada reamostragem.

Diferenças entre modelos serão examinadas por bootstrap pareado, calculando
diretamente, em cada reamostragem, $\Delta M=M_A-M_B$ sobre os mesmos exemplos.
Não será inferida diferença pela simples sobreposição de intervalos individuais.
Os quatro contrastes primários da Tabela~\ref{tab:matriz-rq} serão apresentados
sem expansão posterior do conjunto de hipóteses. Para famílias de comparações
secundárias entre modelos, plataformas ou categorias, os valores de $p$ serão
ajustados pelo procedimento de Holm \cite{holm1979procedure}. Magnitude,
intervalo de confiança e relação
com $\delta$ serão priorizados sobre uma decisão binária de significância.

Além dos resultados globais, as métricas serão desagregadas por plataforma,
tipo de conteúdo e categoria de ódio. Subgrupos com menos de 30 exemplos
positivos não receberão interpretação comparativa; seus valores serão exibidos
apenas como evidência exploratória, acompanhados do tamanho da amostra e do
intervalo de confiança.

Para avaliar generalização entre plataformas, será realizado um experimento
\textit{leave-one-platform-out}: em cada rodada, uma plataforma será mantida
inteiramente fora do ajuste e usada como domínio de teste. A rodada somente será
confirmatória quando a plataforma excluída possuir ao menos 100 registros e 20
positivos; nos demais casos será exploratória. Os resultados serão reportados
para YouTube, WhatsApp, Telegram, Twitter/X, Facebook, Reddit e TikTok, sem
combinar artificialmente plataformas pequenas para obter significância.

Também serão registrados parâmetros, GPU-horas, memória máxima, duração do
treinamento e tempo de inferência por mil textos. A interpretação final comparará
ganho de desempenho e custo, evitando concluir que um modelo maior é preferível
quando a diferença estiver abaixo de $\delta$ e o custo for substancialmente
superior.

\section{Testes funcionais e análise de erros}
\label{sec:testes-funcionais}

Métricas agregadas em um conjunto de teste podem ocultar falhas sistemáticas de
modelos de detecção de discurso de ódio. O HateCheck demonstra a utilidade de
testes funcionais construídos para fenômenos linguísticos específicos
\cite{rottger2021hatecheck}. Com base nessa abordagem, será produzido um pequeno
conjunto diagnóstico em português brasileiro, separado de treinamento,
validação e teste estatístico.

O conjunto funcional terá, como meta inicial, pelo menos 20 exemplos por
fenômeno, equilibrando casos positivos e negativos quando isso for semanticamente
aplicável. Os exemplos poderão ser adaptados de taxonomias funcionais publicadas,
construídos de forma controlada ou extraídos do domínio, desde que origem e
licença sejam registradas. O conjunto deverá contemplar, no mínimo:

\begin{itemize}
    \item negação de uma afirmação preconceituosa;
    \item citação de discurso de ódio para denúncia ou análise;
    \item contra-discurso e defesa do grupo atacado;
    \item ironia e sarcasmo;
    \item menção neutra a grupos sociais;
    \item insulto individual sem ataque generalizante a um grupo;
    \item grafias alteradas, separação de caracteres e uso de símbolos;
    \item ameaça, incitação, desumanização e criminalização coletiva;
    \item ataques explícitos e implícitos;
    \item variações de gênero, número e referência pronominal.
\end{itemize}

Os casos serão classificados por dois avaliadores independentes e divergências
serão adjudicadas pelo protocolo da Seção~\ref{sec:qualidade-anotacao}. Antes do
uso, cada caso será comparado com treino, calibração e teste por hash, Jaccard e
busca aproximada; casos equivalentes serão removidos ou reescritos e novamente
validados. O conjunto funcional não será usado para ajustar hiperparâmetros. O
desempenho será relatado por fenômeno, sem misturá-lo à estimativa principal do
teste da v0.1.

Após a avaliação quantitativa, será conduzida análise qualitativa de erros. Para
cada modelo serão examinados, no mínimo, 50 falsos positivos e 50 falsos
negativos da tarefa binária, sorteados com semente 2026 e estratificação por
plataforma e faixa de confiança. Na tarefa multirrótulo serão sorteados erros das
categorias mais frequentes e incluídos os erros disponíveis das categorias raras,
respeitando as restrições de privacidade.

Antes de examinar as predições será congelada a seguinte taxonomia inicial:
\texttt{E01} contexto ausente; \texttt{E02} ironia; \texttt{E03} sarcasmo;
\texttt{E04} negação; \texttt{E05} alvo implícito; \texttt{E06} ambiguidade
lexical; \texttt{E07} linguagem codificada; \texttt{E08} truncamento;
\texttt{E09} possível erro de anotação; \texttt{E10} conflito entre categorias;
\texttt{E11} conhecimento externo necessário; e \texttt{E12} dependência
multimodal. Novos códigos poderão emergir, mas serão identificados como
pós-especificados. Uma parcela mínima de 20\% dos erros será codificada por dois
avaliadores, com concordância reportada.

Cada item indicará se depende de imagem, vídeo, mensagem anterior, thread,
evento externo ou conhecimento sociopolítico. Será reportada a proporção de
erros potencialmente insolúveis pelo texto isolado, distinguindo limitação do
modelo de limitação da modalidade de entrada.

\section{Reprodutibilidade, documentação e uso responsável}
\label{sec:reprodutibilidade}

Cada execução deverá gerar um manifesto contendo:

\begin{itemize}
    \item versão e hash SHA-256 do dataset;
    \item hash do código e da consulta de extração;
    \item commit do repositório, estado de alterações locais e comando executado;
    \item arquivos imutáveis de \texttt{train\_ids}, \texttt{validation\_ids},
    \texttt{calibration\_ids}, \texttt{test\_ids},
    \texttt{duplicate\_groups} e \texttt{context\_groups};
    \item checkpoint e revisão do tokenizer;
    \item versões de Python, PyTorch, Transformers e demais dependências;
    \item modelo, tarefa, hiperparâmetros e semente;
    \item número de épocas, melhor época nos folds e curvas de perda;
    \item hardware, precisão numérica, memória máxima, GPU-horas e tempo de
    treinamento e inferência;
    \item limiares, temperatura de calibração e pesos de classe;
    \item métricas globais, ponderadas e por subgrupo;
    \item caminho e hash dos checkpoints resultantes.
\end{itemize}

O ambiente de software será congelado por arquivo de dependências com versões
exatas e, preferencialmente, por imagem de contêiner. Não se dependerá apenas da
semente para reconstruir partições. As configurações de determinismo do PyTorch,
CUDA e cuDNN serão registradas, incluindo flags ativadas e operações conhecidas
como não determinísticas. Quando determinismo estrito reduzir de modo relevante
o desempenho computacional ou não estiver disponível, a exceção será documentada.

A documentação do dataset descreverá origem, composição, processo de amostragem,
população representada, critérios de anotação, limitações e usos não recomendados.
Esse procedimento é coerente com a proposta de \textit{data statements}, que
busca tornar explícitas as condições sob as quais resultados de PLN podem ou não
ser generalizados \cite{bender2018datastatements}. Cada modelo terá também uma
ficha com finalidade, desempenho, grupos avaliados, limitações, riscos e uso
pretendido, seguindo a proposta de \textit{model cards}
\cite{mitchell2019modelcards}.

Os classificadores não deverão ser descritos como substitutos da avaliação
humana. Seu uso previsto é priorizar, apoiar ou acelerar a triagem de conteúdos.
Decisões de alto impacto, estudos de prevalência e análises sobre grupos pouco
representados deverão considerar a incerteza, o desenho amostral e a revisão de
especialistas.

Exemplos divulgados em artigos, dissertação ou fichas de modelo passarão por
revisão específica de risco de reidentificação. Conforme o risco, serão usados
mascaramento adicional, paráfrase fiel, exemplo sintético controlado ou omissão.
Identificadores de origem e textos brutos não acompanharão os artefatos públicos.

\section{Continuidade do processo}
\label{sec:continuidade-processo}

O teste da v0.1 permanecerá congelado nas versões posteriores, permitindo medir
ganhos ou regressões no mesmo conjunto histórico. Contudo, após ter sido
consultado em sucessivos ciclos, ele será interpretado como benchmark histórico
ou teste de regressão, e não indefinidamente como única avaliação confirmatória.
Novos registros avaliados serão incorporados ao desenvolvimento das versões
v0.2, v0.3 e seguintes. Cada versão reservará um teste temporal formado por
conteúdos estritamente posteriores ao período de treinamento; o teste temporal
mais recente assumirá progressivamente o papel confirmatório.

Para cada atualização serão repetidos os controles de elegibilidade,
duplicidade, estratificação, treinamento e avaliação. A data de corte temporal
será definida antes do treinamento e os grupos de contexto não poderão atravessá-la.
Um modelo novo somente substituirá o anterior se não apresentar regressão
relevante no teste histórico e demonstrar desempenho aceitável no teste temporal
mais recente.

O pré-treinamento adaptativo em dados não rotulados do domínio poderá ser
investigado em uma etapa futura. Estudos anteriores mostram que a continuação do
pré-treinamento em textos do domínio ou da tarefa pode produzir ganhos mesmo em
cenários com poucos rótulos \cite{gururangan2020dontstop}. Contudo, esse
procedimento não integrará o benchmark principal da v0.1, que comparará os
checkpoints públicos originais sob o mesmo protocolo supervisionado.

\section{Estudos de ablação}
\label{sec:ablacoes}

As decisões de engenharia que possam alterar as conclusões serão transformadas
em contrastes controlados. As ablações pré-especificadas são:

\begin{itemize}
    \item perda comum versus perda ponderada;
    \item arquitetura hierárquica versus conjunta;
    \item encoder monolíngue versus multilíngue;
    \item probabilidades brutas versus calibradas;
    \item truncamento apenas no início versus início e fim e, se necessário,
    \textit{chunking};
    \item texto anonimizado com marcadores especiais versus anonimização sem os
    marcadores explícitos.
\end{itemize}

Cada ablação alterará um componente por vez, manterá partições e sementes e será
analisada de forma pareada. Perda, idioma, calibração e propagação hierárquica
correspondem às perguntas primárias; truncamento e mascaramento serão análises
secundárias no backbone vencedor, salvo se o piloto mostrar efeito suficiente
para justificar expansão previamente documentada.

\section{Fluxo experimental consolidado}
\label{sec:fluxo-experimental}

O procedimento completo seguirá a ordem abaixo. A posição do teste no fluxo
representa isolamento lógico: seus rótulos somente serão acessados depois que os
artefatos anteriores estiverem congelados.

\begin{enumerate}
    \item definir universo, estratos, probabilidades de inclusão e lotes;
    \item concluir a anotação simples, auditar a dupla anotação histórica e
    documentar concordância, divergências e exclusões;
    \item congelar dataset, taxonomia e manual, com hashes e manifests;
    \item executar anonimização, normalização e testes de consistência;
    \item detectar duplicatas, formar componentes conexos e grupos de contexto;
    \item criar e versionar as partições de seleção, calibração e teste;
    \item perfilar comprimentos e executar apenas as ablações pré-especificadas;
    \item realizar triagem de hiperparâmetros e validação cruzada;
    \item produzir e concatenar predições OOF para seleção de configuração;
    \item fixar época e treinar as cinco sementes finais;
    \item formar o ensemble, ajustar calibração e limiares no conjunto próprio;
    \item congelar checkpoint, tokenizer, calibrador, limiares e código;
    \item abrir uma única vez o teste para a avaliação principal;
    \item executar bootstrap pareado, análises por plataforma e eficiência;
    \item executar avaliações externa, temporal e funcional quando disponíveis;
    \item conduzir análise qualitativa reprodutível e publicar a documentação.
\end{enumerate}

\section{Ameaças à validade}
\label{sec:ameacas-validade}

As principais ameaças à interpretação dos resultados são:

\begin{itemize}
    \item \textbf{viés de seleção}: a amostra foi balanceada com base em sinais
    do Gemma e pode não reproduzir diretamente a prevalência do universo;
    \item \textbf{revisão simples}: a predominância de decisões
    \texttt{single} impede estimar concordância entre avaliadores na maior parte
    da v0.1;
    \item \textbf{classes raras}: algumas categorias possuem poucas dezenas de
    exemplos e não sustentam conclusões confirmatórias;
    \item \textbf{mudança temporal}: linguagem ofensiva, códigos e plataformas
    podem mudar depois do congelamento;
    \item \textbf{dependência de contexto}: o texto isolado pode não conter
    informações presentes na conversa, imagem, vídeo ou evento relacionado;
    \item \textbf{propagação de erros}: no pipeline hierárquico, um falso
    negativo binário impede a identificação dos tipos de ódio;
    \item \textbf{generalização limitada}: os resultados serão válidos para a
    composição observada de plataformas, períodos e categorias, não para todo
    conteúdo em português brasileiro.
\end{itemize}

Essas limitações deverão acompanhar qualquer divulgação de métricas ou uso do
modelo. A expansão contínua do dataset será orientada também pelos erros e pelas
lacunas de cobertura detectadas na v0.1.

\section{Tabelas para registro dos resultados}
\label{sec:tabelas-resultados}

As tabelas seguintes definem o formato de apresentação, mas não contêm valores
experimentais. Os campos marcados com ``---'' deverão ser preenchidos somente
após a execução do protocolo.

\begin{table}[htbp]
    \centering
    \caption{Modelo de tabela para a tarefa binária no teste fixo.}
    \label{tab:modelo-resultados-binarios}
    \resizebox{\textwidth}{!}{%
    \begin{tabular}{lrrrrrrr}
        \toprule
        \textbf{Modelo} & \textbf{Macro-F1} & \textbf{F1 hate} &
        \textbf{Precisão hate} & \textbf{Revocação hate} & \textbf{AUPRC} &
        \textbf{MCC} & \textbf{ECE} \\
        \midrule
        TF--IDF + regressão logística & --- & --- & --- & --- & --- & --- & --- \\
        Gemma & --- & --- & --- & --- & --- & --- & --- \\
        BERTimbau Base & --- & --- & --- & --- & --- & --- & --- \\
        BERTimbau Large & --- & --- & --- & --- & --- & --- & --- \\
        BERTabaporu Base & --- & --- & --- & --- & --- & --- & --- \\
        mBERT & --- & --- & --- & --- & --- & --- & --- \\
        XLM-R Base & --- & --- & --- & --- & --- & --- & --- \\
        \bottomrule
    \end{tabular}}
    \par\smallskip
    {\footnotesize Fonte: elaboração própria. Valores a serem preenchidos após
    os experimentos.\par}
\end{table}

\begin{table}[htbp]
    \centering
    \caption{Modelo de tabela para a tarefa multirrótulo.}
    \label{tab:modelo-resultados-multirrotulo}
    \resizebox{\textwidth}{!}{%
    \begin{tabular}{lrrrrrr}
        \toprule
        \textbf{Modelo} & \textbf{Macro-F1 elegível} &
        \textbf{Macro-F1 total} & \textbf{Micro-F1} & \textbf{PR-AUC} &
        \textbf{Hamming loss} & \textbf{Subset accuracy} \\
        \midrule
        TF--IDF + regressão logística & --- & --- & --- & --- & --- & --- \\
        BERTimbau Base & --- & --- & --- & --- & --- & --- \\
        BERTimbau Large & --- & --- & --- & --- & --- & --- \\
        BERTabaporu Base & --- & --- & --- & --- & --- & --- \\
        mBERT & --- & --- & --- & --- & --- & --- \\
        XLM-R Base & --- & --- & --- & --- & --- & --- \\
        \bottomrule
    \end{tabular}}
    \par\smallskip
    {\footnotesize Fonte: elaboração própria. Valores a serem preenchidos após
    os experimentos.\par}
\end{table}

\begin{table}[htbp]
    \centering
    \caption{Modelo de tabela para resultados por categoria de ódio.}
    \label{tab:modelo-resultados-categoria}
    \resizebox{\textwidth}{!}{%
    \begin{tabular}{lrrrrrr}
        \toprule
        \textbf{Categoria} & \textbf{$n$ teste} & \textbf{Precisão} &
        \textbf{Revocação} & \textbf{F1} & \textbf{IC95\% F1} & \textbf{PR-AUC} \\
        \midrule
        Ageísmo & --- & --- & --- & --- & --- & --- \\
        Classismo & --- & --- & --- & --- & --- & --- \\
        Gordofobia & --- & --- & --- & --- & --- & --- \\
        Capacitismo & --- & --- & --- & --- & --- & --- \\
        LGBTfobia & --- & --- & --- & --- & --- & --- \\
        Ódio político & --- & --- & --- & --- & --- & --- \\
        Racismo & --- & --- & --- & --- & --- & --- \\
        Intolerância religiosa & --- & --- & --- & --- & --- & --- \\
        Sexismo & --- & --- & --- & --- & --- & --- \\
        Xenofobia & --- & --- & --- & --- & --- & --- \\
        \bottomrule
    \end{tabular}}
    \par\smallskip
    {\footnotesize Fonte: elaboração própria. Categorias abaixo do limiar
    confirmatório deverão ser identificadas como exploratórias.\par}
\end{table}

\begin{table}[htbp]
    \centering
    \caption{Modelo de tabela para desempenho e eficiência computacional.}
    \label{tab:modelo-eficiencia}
    \resizebox{\textwidth}{!}{%
    \begin{tabular}{lrrrrr}
        \toprule
        \textbf{Modelo} & \textbf{Macro-F1} & \textbf{Parâmetros} &
        \textbf{GPU-horas} & \textbf{Memória máxima} &
        \textbf{Tempo/1.000 textos} \\
        \midrule
        BERTimbau Base & --- & --- & --- & --- & --- \\
        BERTimbau Large & --- & --- & --- & --- & --- \\
        BERTabaporu Base & --- & --- & --- & --- & --- \\
        mBERT & --- & --- & --- & --- & --- \\
        XLM-R Base & --- & --- & --- & --- & --- \\
        \bottomrule
    \end{tabular}}
    \par\smallskip
    {\footnotesize Fonte: elaboração própria. Valores a serem preenchidos após
    os experimentos.\par}
\end{table}

\section{Critérios de conclusão da etapa v0.1}
\label{sec:criterios-conclusao}

A etapa de treinamento da v0.1 será considerada metodologicamente concluída
quando:

\begin{enumerate}
    \item o primeiro lote completo que alcance ao menos 21.248 decisões válidas
    tiver sido congelado e versionado;
    \item a suficiência amostral tiver sido examinada e as análises confirmatórias
    e exploratórias estiverem separadas;
    \item a dupla anotação histórica, as medidas de concordância e as exclusões
    por divergência tiverem sido documentadas, e a ausência de dupla anotação
    contemporânea estiver explicitamente declarada como limitação;
    \item as partições tiverem sido verificadas quanto a duplicatas e vazamento;
    \item todos os baselines e cinco encoders tiverem sido executados nas mesmas
    divisões;
    \item os hiperparâmetros tiverem sido definidos por OOF e os limiares e a
    calibração pelo conjunto próprio, sem uso do teste;
    \item as cinco sementes tiverem sido executadas para cada configuração
    selecionada;
    \item o ensemble final tiver sido fixado sem seleção posterior de semente;
    \item as métricas globais, ponderadas, por plataforma e por categoria
    tiverem sido calculadas com intervalos de confiança e comparações pareadas;
    \item a análise funcional e a análise qualitativa de erros tiverem sido
    documentadas;
    \item os manifests, checkpoints, hashes, tabelas e fichas de documentação
    tiverem sido arquivados de forma reprodutível.
\end{enumerate}

% -----------------------------------------------------------------------------
% Referências incorporadas ao próprio capítulo
% -----------------------------------------------------------------------------

\begin{thebibliography}{99}

\bibitem[Bender e Friedman(2018)]{bender2018datastatements}
BENDER, Emily M.; FRIEDMAN, Batya.
Data statements for natural language processing: toward mitigating system bias
and enabling better science.
\textit{Transactions of the Association for Computational Linguistics}, v. 6,
p. 587--604, 2018.
DOI: \url{https://doi.org/10.1162/tacl_a_00041}.

\bibitem[Cohen(1960)]{cohen1960kappa}
COHEN, Jacob.
A coefficient of agreement for nominal scales.
\textit{Educational and Psychological Measurement}, v. 20, n. 1, p. 37--46,
1960.
DOI: \url{https://doi.org/10.1177/001316446002000104}.

\bibitem[Conneau et al.(2020)]{conneau2020xlmr}
CONNEAU, Alexis et al.
Unsupervised cross-lingual representation learning at scale.
In: \textit{Proceedings of the 58th Annual Meeting of the Association for
Computational Linguistics}. Online: Association for Computational Linguistics,
2020. p. 8440--8451.
Disponível em: \url{https://arxiv.org/abs/1911.02116}.

\bibitem[da Costa et al.(2023)]{dacosta2023bertabaporu}
DA COSTA, Pablo Botton; PAVAN, Matheus Camasmie; DOS SANTOS, Wesley Ramos;
DA SILVA, Samuel Caetano; PARABONI, Ivandre.
BERTabaporu: assessing a genre-specific language model for Portuguese NLP.
In: \textit{Proceedings of the 14th International Conference on Recent Advances
in Natural Language Processing}. Varna: INCOMA, 2023. p. 217--223.
DOI: \url{https://doi.org/10.26615/978-954-452-092-2_024}.

\bibitem[da Costa(2026)]{dacosta2026bertabaporu}
DA COSTA, Pablo Botton.
\textit{BERTabaporu Base Uncased: model card}.
Hugging Face, 2026. Disponível em:
\url{https://huggingface.co/pablocosta/bertabaporu-base-uncased}.
Acesso em: 11 set. 2026.

\bibitem[Devlin et al.(2019)]{devlin2019bert}
DEVLIN, Jacob; CHANG, Ming-Wei; LEE, Kenton; TOUTANOVA, Kristina.
BERT: pre-training of deep bidirectional transformers for language understanding.
In: \textit{Proceedings of the 2019 Conference of the North American Chapter of
the Association for Computational Linguistics: Human Language Technologies}.
Minneapolis: Association for Computational Linguistics, 2019. p. 4171--4186.
Disponível em: \url{https://arxiv.org/abs/1810.04805}.

\bibitem[Dodge et al.(2020)]{dodge2020finetuning}
DODGE, Jesse et al.
Fine-tuning pretrained language models: weight initializations, data orders,
and early stopping.
\textit{arXiv preprint arXiv:2002.06305}, 2020.
Disponível em: \url{https://arxiv.org/abs/2002.06305}.

\bibitem[Efron e Tibshirani(1993)]{efron1993bootstrap}
EFRON, Bradley; TIBSHIRANI, Robert J.
\textit{An introduction to the bootstrap}.
New York: Chapman \& Hall, 1993.
DOI: \url{https://doi.org/10.1007/978-1-4899-4541-9}.

\bibitem[Google BERT(2026)]{google2026mbert}
GOOGLE BERT.
\textit{BERT multilingual base model (cased): model card}.
Hugging Face, 2026.
Disponível em:
\url{https://huggingface.co/google-bert/bert-base-multilingual-cased}.
Acesso em: 1 set. 2026.

\bibitem[Guo et al.(2017)]{guo2017calibration}
GUO, Chuan; PLEISS, Geoff; SUN, Yu; WEINBERGER, Kilian Q.
On calibration of modern neural networks.
In: \textit{Proceedings of the 34th International Conference on Machine Learning}.
Sydney: PMLR, 2017. v. 70, p. 1321--1330.
Disponível em: \url{https://arxiv.org/abs/1706.04599}.

\bibitem[Holm(1979)]{holm1979procedure}
HOLM, Sture.
A simple sequentially rejective multiple test procedure.
\textit{Scandinavian Journal of Statistics}, v. 6, n. 2, p. 65--70, 1979.

\bibitem[Gururangan et al.(2020)]{gururangan2020dontstop}
GURURANGAN, Suchin et al.
Don't stop pretraining: adapt language models to domains and tasks.
In: \textit{Proceedings of the 58th Annual Meeting of the Association for
Computational Linguistics}. Online: Association for Computational Linguistics,
2020. p. 8342--8360.
DOI: \url{https://doi.org/10.18653/v1/2020.acl-main.740}.

\bibitem[Krippendorff(2018)]{krippendorff2018content}
KRIPPENDORFF, Klaus.
\textit{Content analysis: an introduction to its methodology}.
4. ed. Thousand Oaks: SAGE, 2018.

\bibitem[Meta AI(2026)]{meta2026xlmr}
META AI.
\textit{XLM-RoBERTa base: model card}.
Hugging Face, 2026.
Disponível em: \url{https://huggingface.co/FacebookAI/xlm-roberta-base}.
Acesso em: 1 set. 2026.

\bibitem[Mitchell et al.(2019)]{mitchell2019modelcards}
MITCHELL, Margaret et al.
Model cards for model reporting.
In: \textit{Proceedings of the Conference on Fairness, Accountability, and
Transparency}. New York: ACM, 2019. p. 220--229.
DOI: \url{https://doi.org/10.1145/3287560.3287596}.

\bibitem[NeuralMind(2026)]{neuralmind2026bertimbau}
NEURALMIND.
\textit{BERTimbau Base and Large: model cards}.
Hugging Face, 2026.
Disponível em:
\url{https://huggingface.co/neuralmind/bert-base-portuguese-cased} e
\url{https://huggingface.co/neuralmind/bert-large-portuguese-cased}.
Acesso em: 1 set. 2026.

\bibitem[Röttger et al.(2021)]{rottger2021hatecheck}
RÖTTGER, Paul; VIDGEN, Bertie; NGUYEN, Dong; WASEEM, Zeerak; MARGETTS, Helen;
PIERREHUMBERT, Janet.
HateCheck: functional tests for hate speech detection models.
In: \textit{Proceedings of the 59th Annual Meeting of the Association for
Computational Linguistics and the 11th International Joint Conference on
Natural Language Processing}. Online: Association for Computational Linguistics,
2021. p. 41--58.
DOI: \url{https://doi.org/10.18653/v1/2021.acl-long.4}.

\bibitem[Saito e Rehmsmeier(2015)]{saito2015precisionrecall}
SAITO, Takaya; REHMSMEIER, Marc.
The precision-recall plot is more informative than the ROC plot when evaluating
binary classifiers on imbalanced datasets.
\textit{PLOS ONE}, v. 10, n. 3, e0118432, 2015.
DOI: \url{https://doi.org/10.1371/journal.pone.0118432}.

\bibitem[Sechidis, Tsoumakas e Vlahavas(2011)]{sechidis2011stratification}
SECHIDIS, Konstantinos; TSOUMAKAS, Grigorios; VLAHAVAS, Ioannis.
On the stratification of multi-label data.
In: \textit{Machine Learning and Knowledge Discovery in Databases}.
Berlin: Springer, 2011. p. 145--158.
DOI: \url{https://doi.org/10.1007/978-3-642-23808-6_10}.

\bibitem[Souza, Nogueira e Lotufo(2020)]{souza2020bertimbau}
SOUZA, Fábio; NOGUEIRA, Rodrigo; LOTUFO, Roberto.
BERTimbau: pretrained BERT models for Brazilian Portuguese.
In: \textit{Intelligent Systems: 9th Brazilian Conference, BRACIS 2020}.
Cham: Springer, 2020. p. 403--417.
DOI: \url{https://doi.org/10.1007/978-3-030-61377-8_28}.

\bibitem[Wolf et al.(2020)]{wolf2020transformers}
WOLF, Thomas et al.
Transformers: state-of-the-art natural language processing.
In: \textit{Proceedings of the 2020 Conference on Empirical Methods in Natural
Language Processing: System Demonstrations}. Online: Association for
Computational Linguistics, 2020. p. 38--45.
DOI: \url{https://doi.org/10.18653/v1/2020.emnlp-demos.6}.

\end{thebibliography}

%
% Pacotes recomendados no preâmbulo do documento principal:
% \usepackage[T1]{fontenc}
% \usepackage[utf8]{inputenc} % dispensável em distribuições LaTeX recentes
% \usepackage[brazil]{babel}
% \usepackage{amsmath,amssymb}
% \usepackage{booktabs,longtable,array}
% \usepackage{graphicx}
% \usepackage{url}
% \usepackage[hidelinks]{hyperref}
% \usepackage[alf]{abntex2cite}
%
% As referências foram incorporadas ao final deste arquivo para que o capítulo
% não dependa de um arquivo .bib externo. Caso o documento principal já possua
% bibliografia própria, os itens podem ser transferidos para ela sem alterar as
% chaves usadas nas citações.

\chapter{Metodologia de treinamento dos modelos BERT}
\label{cap:metodologia-bert-v01}

\section{Visão geral e objetivos}
\label{sec:visao-geral-bert}

Este capítulo estabelece o protocolo para construção, treinamento, seleção e
avaliação dos primeiros classificadores supervisionados do projeto Radar Ódio
a partir do \textit{Dataset} v0.1. Essa versão será formada pelas avaliações
humanas resolvidas que correspondam a, no mínimo, 10\% do universo elegível. O
objetivo não é tratar a v0.1 como versão definitiva do conjunto de dados, mas
como o primeiro marco reprodutível de um processo contínuo de aprendizagem e
avaliação, a ser ampliado nas versões posteriores.

A estratégia emprega modelos pré-treinados da família BERT e derivados. O uso
de representações contextuais pré-treinadas permite adaptar um modelo de língua
a uma tarefa supervisionada mediante a adição de uma camada de classificação e
o ajuste conjunto dos parâmetros, procedimento conhecido como \textit{fine-tuning}
\cite{devlin2019bert}. Para o português brasileiro, os modelos BERTimbau
constituem referências monolíngues relevantes, tendo sido pré-treinados em um
grande corpus nacional e avaliados em diferentes tarefas de processamento de
linguagem natural \cite{souza2020bertimbau}. Modelos multilíngues serão incluídos
como controles comparativos, permitindo verificar se a especialização no idioma
produz vantagem no domínio de textos oriundos de redes sociais.

O protocolo contempla duas tarefas relacionadas, porém treinadas e avaliadas
separadamente:

\begin{enumerate}
    \item \textbf{detecção binária de discurso de ódio}: classificação de cada
    texto como \texttt{hate} ou \texttt{nao\_hate};
    \item \textbf{classificação multirrótulo dos tipos de ódio}: atribuição de
    zero ou mais categorias temáticas aos registros cujo rótulo gold seja
    \texttt{hate}.
\end{enumerate}

Essa separação define a arquitetura hierárquica candidata: o primeiro modelo
decide se o conteúdo contém discurso de ódio; quando a decisão é positiva, o
segundo identifica os tipos aplicáveis. Sua adoção não será tratada como uma
vantagem presumida. Ela será comparada a uma arquitetura conjunta e terá a
propagação de erros quantificada, conforme a Seção~\ref{sec:arquiteturas}.

\section{Perguntas de pesquisa, hipóteses e contribuição}
\label{sec:perguntas-hipoteses}

Os modelos são empregados como instrumentos para investigar generalização,
desbalanceamento, calibração e propagação de erros, e não apenas como alternativas
de implementação. As perguntas de pesquisa são:

\begin{description}
    \item[RQ1:] um encoder monolíngue pré-treinado em português brasileiro
    apresenta desempenho superior a encoders multilíngues comparáveis na
    detecção de discurso de ódio?
    \item[RQ2:] a ponderação da função de perda melhora a identificação da
    classe minoritária \texttt{hate} e dos tipos menos frequentes?
    \item[RQ3:] a calibração pós-processada reduz o erro das probabilidades
    estimadas sem deteriorar a discriminação?
    \item[RQ4:] o desempenho varia entre plataformas e se mantém quando uma
    plataforma inteira é excluída do treinamento?
    \item[RQ5:] qual é a perda de desempenho causada pela propagação de erros
    entre a detecção binária e a classificação dos tipos?
    \item[RQ6:] o desempenho se mantém em dados coletados depois do período
    representado pela v0.1?
    \item[RQ7:] o pré-treinamento específico em textos de redes sociais produz
    vantagem sobre o pré-treinamento geral da web, especialmente no X/Twitter?
\end{description}

As hipóteses e os contrastes primários serão registrados antes da abertura do
teste. Para RQ1, o contraste principal será BERTimbau Base contra mBERT, pois
ambos possuem porte semelhante; BERTimbau Large e XLM-R serão comparações
secundárias. Para RQ7, o contraste secundário pré-especificado será BERTabaporu
Base contra BERTimbau Base, mantendo as mesmas partições e o mesmo protocolo de ajuste. Seja
$M$ a métrica previamente indicada para cada pergunta. Serão
testadas as seguintes hipóteses direcionais:

\begin{align}
    H_{0,1}&: M_{\mathrm{BERTimbau\ Base}}-M_{\mathrm{mBERT}} \leq 0,
    &H_{1,1}&: M_{\mathrm{BERTimbau\ Base}}-M_{\mathrm{mBERT}} > 0;\\
    H_{0,2}&: M_{\mathrm{ponderada}}-M_{\mathrm{comum}} \leq 0,
    &H_{1,2}&: M_{\mathrm{ponderada}}-M_{\mathrm{comum}} > 0;\\
    H_{0,3}&: E_{\mathrm{calibrado}}-E_{\mathrm{bruto}} \geq 0,
    &H_{1,3}&: E_{\mathrm{calibrado}}-E_{\mathrm{bruto}} < 0;\\
    H_{0,5}&: F1_{\mathrm{oracle}}-F1_{\mathrm{ponta\ a\ ponta}} \leq 0,
    &H_{1,5}&: F1_{\mathrm{oracle}}-F1_{\mathrm{ponta\ a\ ponta}} > 0.
\end{align}

Em RQ3, $E$ designa o Brier score, com o ECE como medida complementar. RQ4 e
RQ6 serão avaliadas como hipóteses de heterogeneidade e degradação de domínio,
respectivamente. RQ6 é prospectiva e somente poderá receber resposta
confirmatória quando houver um teste temporal posterior à v0.1.

Além de significância estatística, será adotada como referência uma diferença
mínima cientificamente relevante de $\delta=0{,}02$ em Macro-F1 ou F1 da classe
\texttt{hate}. Esse valor constitui uma decisão operacional prévia, e não uma
constante universal; análises de sensibilidade usarão também $\delta=0{,}01$ e
$\delta=0{,}03$. A conclusão distinguirá: ausência de evidência, diferença
estatisticamente compatível com zero e diferença superior ao limiar de relevância
prática.

\begin{longtable}{p{0.06\textwidth}p{0.25\textwidth}p{0.30\textwidth}p{0.18\textwidth}}
    \caption{Matriz pré-especificada entre pergunta, hipótese, experimento e métrica.}
    \label{tab:matriz-rq} \\
    \toprule
    \textbf{RQ} & \textbf{Contraste} & \textbf{Experimento} & \textbf{Métrica principal} \\
    \midrule
    \endfirsthead
    \toprule
    \textbf{RQ} & \textbf{Contraste} & \textbf{Experimento} & \textbf{Métrica principal} \\
    \midrule
    \endhead
    RQ1 & monolíngue versus multilíngue & BERTimbau Base versus mBERT nas mesmas partições & Macro-F1 \\
    RQ2 & perda ponderada versus comum & ablação pareada por backbone e semente & F1-\texttt{hate}; AUPRC \\
    RQ3 & probabilidade calibrada versus bruta & \textit{temperature scaling} em conjunto próprio & Brier score; ECE \\
    RQ4 & dentro versus fora do domínio & resultados por plataforma e \textit{leave-one-platform-out} & Macro-F1 \\
    RQ5 & \textit{oracle} versus ponta a ponta & execução das duas formas sobre os mesmos exemplos & $\Delta$Macro-F1 \\
    RQ6 & período histórico versus posterior & teste temporal nas versões posteriores & Macro-F1; AUPRC \\
    RQ7 & pré-treinamento social versus web & BERTabaporu Base versus BERTimbau Base, globalmente e por plataforma & Macro-F1 \\
    \bottomrule
\end{longtable}

A contribuição científica esperada é, portanto, a caracterização controlada de
como idioma de pré-treinamento, desenho amostral, desbalanceamento, calibração,
plataforma e organização hierárquica afetam classificadores de discurso de ódio
em português brasileiro. A construção do dataset e a implementação dos modelos
são meios para essa investigação. Resultados confirmatórios serão separados de
análises exploratórias, especialmente para categorias raras e plataformas com
baixa cobertura.

\section{Definição e congelamento do Dataset v0.1}
\label{sec:dataset-v01}

\subsection{Retrato da base em 10 de setembro de 2026}

Na data de referência deste protocolo, o universo continha 212.476 análises com
estado de processamento bem-sucedido. A base gold reunia 25.090 textos canônicos
com decisão humana resolvida, correspondentes a 11,81\% do universo. Desse total,
3.995 registros haviam sido classificados como \texttt{hate} e 21.095 como
\texttt{nao\_hate}. A distribuição evidencia desbalanceamento entre as classes:
aproximadamente 15,92\% dos registros gold eram positivos e 84,08\% eram
negativos.

\begin{table}[htbp]
    \centering
    \caption{Retrato da base gold em 10 de setembro de 2026.}
    \label{tab:retrato-base-gold}
    \begin{tabular}{lrr}
        \toprule
        \textbf{Indicador} & \textbf{Quantidade} & \textbf{Percentual} \\
        \midrule
        Universo elegível & 212.476 & 100,00\% \\
        Decisões gold resolvidas & 25.090 & 11,81\% \\
        Rótulo \texttt{hate} & 3.995 & 15,92\% do gold \\
        Rótulo \texttt{nao\_hate} & 21.095 & 84,08\% do gold \\
        \bottomrule
    \end{tabular}
    \par\smallskip
    {\footnotesize Fonte: elaboração própria com base nas tabelas de avaliação
    do projeto.\par}
\end{table}

As decisões gold estavam distribuídas entre sete plataformas, conforme a
Tabela~\ref{tab:distribuicao-plataforma}. Essa distribuição deverá ser recalculada
no momento do congelamento, pois o acréscimo de avaliações pode alterar tanto as
contagens quanto as proporções.

\begin{table}[htbp]
    \centering
    \caption{Distribuição preliminar das decisões gold por plataforma.}
    \label{tab:distribuicao-plataforma}
    \begin{tabular}{lr}
        \toprule
        \textbf{Plataforma} & \textbf{Quantidade} \\
        \midrule
        YouTube & 7.735 \\
        Facebook & 4.891 \\
        Reddit & 3.101 \\
        WhatsApp & 3.037 \\
        Twitter & 2.895 \\
        Telegram & 2.660 \\
        TikTok & 771 \\
        \midrule
        \textbf{Total} & \textbf{25.090} \\
        \bottomrule
    \end{tabular}
    \par\smallskip
    {\footnotesize Fonte: elaboração própria com base nas tabelas de avaliação
    do projeto.\par}
\end{table}

Entre os textos positivos, uma mesma publicação pode receber mais de um tipo de
ódio. Por isso, a soma das frequências da Tabela~\ref{tab:tipos-odio-preliminar}
é superior ao número total de textos classificados como \texttt{hate}.

\begin{table}[htbp]
    \centering
    \caption{Frequência preliminar dos tipos de ódio na base gold.}
    \label{tab:tipos-odio-preliminar}
    \begin{tabular}{llr}
        \toprule
        \textbf{Código} & \textbf{Categoria} & \textbf{Ocorrências} \\
        \midrule
        \texttt{political} & Ódio político & 1.841 \\
        \texttt{sexism} & Sexismo & 1.020 \\
        \texttt{lgbtphobia} & LGBTfobia & 725 \\
        \texttt{aporophobia} & Classismo & 316 \\
        \texttt{religious\_intolerance} & Intolerância religiosa & 255 \\
        \texttt{racism} & Racismo & 241 \\
        \texttt{xenophobia} & Xenofobia & 141 \\
        \texttt{body\_shame} & Gordofobia & 123 \\
        \texttt{ageism} & Ageísmo & 52 \\
        \texttt{capacitism} & Capacitismo & 39 \\
        \texttt{other} & Categoria histórica inativa & 1 \\
        \bottomrule
    \end{tabular}
    \par\smallskip
    {\footnotesize Fonte: elaboração própria. As frequências são multirrótulo e,
    portanto, não devem ser somadas para estimar o número de textos positivos.\par}
\end{table}

As categorias menos frequentes ainda não oferecem quantidade suficiente para
uma avaliação confirmatória estável. Esse problema não será resolvido por meio
da eliminação automática de categorias ou pela criação de exemplos artificiais
antes da separação dos dados. Em vez disso, todos os rótulos ativos serão
preservados no treinamento e os resultados das classes raras serão explicitamente
marcados como exploratórios e acompanhados de intervalos de confiança.

\subsection{Justificativa do tamanho da v0.1}

O percentual de 10\% constitui um marco operacional para iniciar o ciclo de
treinamento, e não uma justificativa estatística isolada. O retrato efetivamente
alcançado reúne 25.090 registros e 3.995 positivos, com prevalência observada de
15,92\%. Sob amostragem aleatória simples, a margem de erro aproximada de 95\%
para essa prevalência seria de 0,45 ponto percentual. Esse cálculo não será
interpretado como precisão efetiva porque a base foi estratificada, possui pesos
desiguais e contém grupos correlacionados.

Antes do congelamento será produzido um relatório de suficiência com: contagem
esperada e observada de positivos; contagem por categoria, plataforma e tipo de
conteúdo; tamanho efetivo da amostra; largura dos intervalos de confiança
bootstrap; e quantidade de positivos em cada partição. Para pesos $w_i$, o
tamanho efetivo descritivo será calculado por

\begin{equation}
    n_{\mathrm{efetivo}} =
    \frac{\left(\sum_i w_i\right)^2}{\sum_i w_i^2}.
\end{equation}

Uma análise de poder por simulação, baseada apenas na parcela de desenvolvimento
disponível, estimará a probabilidade de detectar uma diferença pareada de
$\delta=0{,}02$ nas comparações primárias, com nível de 5\% e poder-alvo de 80\%.
A simulação preservará prevalência, estratos e grupos e apresentará também a
menor diferença detectável sob o tamanho alcançado. Seus parâmetros e código
serão arquivados; nenhuma observação do teste será usada para ajustar a análise.

O estudo binário será considerado apto à análise confirmatória apenas se o teste
contiver quantidade suficiente de positivos para produzir estimativas estáveis e
se os intervalos alcançarem precisão compatível com as comparações primárias.
Para os tipos de ódio, a classificação confirmatória ou exploratória obedecerá
aos critérios da Seção~\ref{sec:tarefa-multirrotulo}. Assim, ultrapassar o limiar
de 21.248 registros autoriza o início do benchmark, mas não transforma automaticamente
todas as análises por subgrupo em confirmatórias.

\subsection{Regra de congelamento}

Dez por cento de 212.476 registros correspondem a 21.247,6 observações. Assim,
o limiar operacional da v0.1 será de 21.248 decisões gold válidas. Como as
avaliações são organizadas em lotes, o congelamento ocorrerá ao término do
primeiro lote completo que fizer a base atingir ou ultrapassar esse limiar. Não
será realizado corte dentro de um lote já iniciado, evitando uma seleção
oportunista de casos.

O congelamento deverá produzir um instantâneo imutável, acompanhado de:

\begin{itemize}
    \item data e hora da extração;
    \item consulta SQL ou código de exportação versionado;
    \item número total de registros e distribuição dos rótulos;
    \item identificador do último lote incorporado;
    \item hash SHA-256 do arquivo exportado;
    \item versão da taxonomia e do protocolo de avaliação;
    \item histórico de exclusões e correções aplicadas.
\end{itemize}

Em 10 de setembro de 2026 foi produzida uma exportação candidata com 25.090
decisões válidas. O volume já supera o limiar mínimo da v0.1. Entretanto, o lote
24 ainda constava administrativamente como \texttt{active}, embora todos os seus
953 itens possuíssem decisão resolvida. Por isso, o snapshot pode iniciar a
preparação experimental, mas seu congelamento definitivo deverá registrar o
encerramento administrativo do lote e confirmar novamente as contagens e hashes.

Uma vez congelada, a v0.1 não deverá receber alterações silenciosas. Correções
posteriores de rótulos deverão gerar uma nova versão do conjunto de dados, com
registro das diferenças em relação à versão anterior.

O manifesto de congelamento deverá identificar separadamente as versões do
dataset, da taxonomia e do manual de anotação, por exemplo:

\begin{quote}
\texttt{dataset = radarodio-v0.1}\par
\texttt{taxonomy = taxonomy-v1.2}\par
\texttt{annotation-guidelines = annotation-guidelines-v1.2}
\end{quote}

Mudanças posteriores no manual ou na taxonomia não modificarão retroativamente
o significado de decisões produzidas sob uma versão anterior.

\subsection{Qualidade da anotação e formação do gold standard}
\label{sec:qualidade-anotacao}

Não será realizada nova dupla anotação sistemática da v0.1, pois esse procedimento
é operacionalmente inviável no cronograma do projeto. O controle de confiabilidade
utilizará os arquivos históricos de dupla anotação do Reddit como estudo auxiliar,
mantido separado do teste principal. A auditoria disponível identificou 1.644
pares binários válidos, dos quais 1.512 apresentaram consenso e 132 divergência,
correspondendo a 91,97\% de concordância bruta e $\kappa$ de Cohen de aproximadamente
0,39. Os 1.512 casos consensuais somente poderão ser incorporados depois de
vinculação ao corpus, deduplicação e harmonização da taxonomia; divergências sem
adjudicação permanecerão fora do gold.

Antes da anotação definitiva, o projeto deverá arquivar uma ficha metodológica
com o número e o perfil funcional dos avaliadores, treinamento recebido, versão
do manual, exemplos de treinamento, exercício de qualificação, critérios de
aprovação e procedimento para dúvidas. Informações pessoais dos avaliadores não
entrarão no arquivo destinado ao treinamento dos modelos.

Na tarefa binária, a concordância será descrita pelo percentual bruto, pelo
kappa de Cohen e pelo alfa de Krippendorff. Na tarefa multirrótulo serão
calculados o Jaccard e o F1 entre conjuntos, além da concordância e do alfa de
Krippendorff por categoria \cite{cohen1960kappa,krippendorff2018content}. Todos
os coeficientes serão acompanhados do número de
itens, distribuição marginal e intervalo de confiança. A concordância não será
usada isoladamente como prova de validade; desacordos serão analisados por
categoria, plataforma e fenômeno linguístico.

Haverá divergência quando os avaliadores históricos discordarem sobre o rótulo
binário ou sobre ao menos um tipo ativo. Sem nova adjudicação rastreável, esses
casos não serão transformados em rótulos de treinamento. Se uma adjudicação
pontual vier a ocorrer, deverá registrar decisão, justificativa, versão das
diretrizes, data e categorias em disputa, preservando as duas respostas originais.

Na data de referência, todas as 25.090 decisões resolvidas no sistema tinham
estado \texttt{single}. A v0.1 poderá ser usada para desenvolvimento e avaliação
interna, mas a ausência de dupla anotação contemporânea e multiplataforma limita
a força das conclusões sobre confiabilidade do gold. O estudo histórico do Reddit
será reportado como evidência auxiliar, sem generalização automática para as
demais plataformas.

\subsection{Critérios de elegibilidade}

Serão incluídos somente textos canônicos cujo consenso possua um dos estados
\texttt{single}, \texttt{agreed} ou \texttt{adjudicated}. Serão excluídos:

\begin{itemize}
    \item avaliações puladas ou ainda atribuídas;
    \item conflitos que não tenham sido adjudicados;
    \item registros sem texto ou com falha de codificação não recuperável;
    \item registros sem rótulo final consistente;
    \item cópias não canônicas de textos duplicados;
    \item observações que violem as regras vigentes da taxonomia.
\end{itemize}

Na data de referência, todas as 25.090 decisões resolvidas tinham estado
\texttt{single}. Consequentemente, não é possível estimar concordância entre
avaliadores para essa parcela da base. Essa condição será registrada como uma
limitação da v0.1. Os casos históricos de revisão dupla serão mantidos em um
conjunto auxiliar com proveniência própria e não alterarão silenciosamente o
estado das decisões produzidas pelo sistema.

O código histórico \texttt{other}, atualmente inativo, não integrará o vetor de
saída da tarefa multirrótulo. O registro que ainda contém esse código poderá ser
mantido na tarefa binária, mas deverá ser excluído da tarefa de tipos até que
seja revisto e recodificado em uma categoria ativa.

Antes da exportação serão aplicadas as seguintes invariantes:

\begin{align}
    \texttt{nao\_hate} &\Rightarrow
    \texttt{gold\_hate\_types}=\varnothing,\\
    \texttt{hate} &\Rightarrow
    \left|\texttt{gold\_hate\_types}\right|\geq 1,
\end{align}

em que todos os tipos devem pertencer à versão ativa da taxonomia. A condição
\texttt{hate} sem tipo não será silenciosamente convertida em negativa: o caso
será devolvido para adjudicação ou ficará fora da tarefa multirrótulo até ser
resolvido. Também serão verificadas unicidade de \texttt{analysis\_id}, presença
dos campos obrigatórios, reprodução do \texttt{text\_hash}, ausência de tipos
inativos e inexistência de rótulos conflitantes dentro de um grupo de duplicatas.
O relatório dessas validações será armazenado com o snapshot.

\subsection{Contrato de exportação}

A extração para treinamento deverá seguir o contrato apresentado na
Tabela~\ref{tab:contrato-exportacao}. Campos administrativos e identidades dos
avaliadores não serão incorporados ao conjunto utilizado pelos modelos.

\begin{longtable}{p{0.22\textwidth}p{0.16\textwidth}p{0.48\textwidth}}
    \caption{Contrato mínimo de exportação do Dataset v0.1.}
    \label{tab:contrato-exportacao} \\
    \toprule
    \textbf{Campo} & \textbf{Tipo} & \textbf{Finalidade} \\
    \midrule
    \endfirsthead
    \toprule
    \textbf{Campo} & \textbf{Tipo} & \textbf{Finalidade} \\
    \midrule
    \endhead
    \texttt{analysis\_id} & inteiro & Identificador interno e rastreabilidade da análise. \\
    \texttt{text\_hash} & texto & Identificação de duplicatas exatas e definição dos grupos de partição. \\
    \texttt{text} & texto & Conteúdo anonimizado que será fornecido ao tokenizer. \\
    \texttt{platform} & categoria & Auditoria e avaliação estratificada por plataforma. \\
    \texttt{content\_kind} & categoria & Identificação de post ou comentário. \\
    \texttt{gold\_label} & categoria & Alvo binário: \texttt{hate} ou \texttt{nao\_hate}. \\
    \texttt{gold\_hate\_types} & lista & Conjunto dos tipos de ódio aplicáveis ao texto. \\
    \texttt{decision\_status} & categoria & Indicação de decisão simples, concordante ou adjudicada. \\
    \texttt{batch\_id} & inteiro & Identificação do lote de amostragem. \\
    \texttt{stratum} & texto & Estrato usado na seleção e no cálculo de pesos amostrais. \\
    \texttt{inclusion\_probability} & numérico & Probabilidade de inclusão, quando conhecida e auditável. \\
    \texttt{duplicate\_group\_id} & texto & Grupo transitivo de duplicatas e quase duplicatas. \\
    \texttt{context\_group\_id} & texto/nulo & Conversa, tópico ou publicação de origem compartilhada. \\
    \texttt{content\_created\_at} & data/hora/nulo & Data do conteúdo para avaliação temporal. \\
    \texttt{reviewed\_at} & data/hora & Controle da versão e avaliações temporais futuras. \\
    \texttt{taxonomy\_version} & texto & Versão da taxonomia aplicada ao rótulo. \\
    \texttt{guidelines\_version} & texto & Versão do manual de anotação. \\
    \texttt{context\_dependency} & categoria/nulo & Indício de dependência de imagem, vídeo, conversa ou evento externo. \\
    \bottomrule
\end{longtable}

Os campos \texttt{gemma\_label}, \texttt{gemma\_hate\_probability} e
\texttt{gemma\_hate\_types} poderão ser exportados em uma tabela auxiliar
restrita à auditoria, ao cálculo dos estratos e à reprodução do baseline. Eles
não poderão ser concatenados ao texto nem empregados como atributos de entrada
dos classificadores BERT. Essa separação impede que os novos modelos apenas
reproduzam os sinais do sistema que orientou a amostragem.

\section{Preparação dos textos}
\label{sec:preparacao-textos}

O pré-processamento será deliberadamente conservador. Modelos BERT utilizam
tokenização em subpalavras e foram concebidos para aprender representações a
partir de sequências próximas da linguagem observada. Remoções agressivas podem
eliminar sinais relevantes para discurso de ódio, especialmente pontuação,
emojis, hashtags, caixa, erros ortográficos e alterações intencionais de
grafia.

Cada texto passará, nesta ordem, pelas seguintes operações:

\begin{enumerate}
    \item conversão para Unicode e normalização na forma NFC;
    \item substituição de sequências inválidas de codificação quando a correção
    for inequívoca;
    \item padronização de quebras de linha e espaços consecutivos;
    \item substituição de URLs pelo marcador \texttt{<URL>};
    \item substituição de nomes de usuário e menções pelo marcador
    \texttt{<USER>};
    \item substituição de endereços de e-mail por \texttt{<EMAIL>};
    \item substituição de números de telefone por \texttt{<PHONE>};
    \item remoção de espaços nas extremidades, sem alterar o restante do texto.
\end{enumerate}

Os marcadores serão adicionados como tokens especiais indivisíveis ao tokenizer
de cada checkpoint, com redimensionamento correspondente da matriz de embeddings.
As novas linhas serão inicializadas pelo procedimento padrão definido na
configuração do checkpoint, treinadas durante o fine-tuning e controladas pelas
mesmas sementes das demais inicializações. A lista de tokens, os IDs resultantes
e o hash do tokenizer alterado serão registrados. Uma ablação no conjunto de
desenvolvimento comparará o mascaramento com a manutenção das formas originais,
após anonimização irreversível dos identificadores, para medir eventual perda de
sinais úteis.

Não serão aplicados \textit{stemming}, lematização, remoção de \textit{stopwords},
tradução, conversão geral para minúsculas ou correção ortográfica. Os modelos
sensíveis à caixa receberão a capitalização original. Expressões regulares,
dicionários e versões de qualquer reconhecedor de informação pessoal serão
versionados. Nomes de grupos sociais necessários à interpretação do alvo não
serão removidos como se fossem identificadores pessoais. O texto bruto ficará
em ambiente de acesso restrito; somente a versão anonimizada poderá integrar os
artefatos de treinamento.

O BERTabaporu é um modelo \textit{uncased}; portanto, sua normalização de caixa
será realizada pelo tokenizer do próprio checkpoint. O arquivo de entrada comum
continuará preservando a capitalização original, evitando que a inclusão desse
modelo modifique os dados recebidos pelos encoders sensíveis à caixa.

\subsection{Duplicatas e prevenção de vazamento}

Primeiramente, o mapa de duplicatas do sistema será usado para manter uma única
linha por \texttt{text\_hash} canônico. Em seguida, será criada uma representação
normalizada em 5-gramas de caracteres. Pares com similaridade de Jaccard maior
ou igual a 0,90 serão tratados como quase duplicatas para fins de partição. A
detecção de quase duplicatas não implicará exclusão automática: os registros
serão reunidos sob um mesmo identificador de grupo, garantindo que permaneçam
na mesma partição e no mesmo fold.

O limiar 0,90 será validado antes da divisão definitiva. Serão sorteados pares
dos intervalos 0,80--0,89, 0,90--0,94 e 0,95--1,00 para julgamento humano cego,
registrando precisão e revocação aproximadas do detector de quase duplicatas.
Se a inspeção indicar separação inadequada, qualquer alteração do limiar será
feita e documentada antes da abertura do teste.

O agrupamento será transitivo. Os textos serão vértices de um grafo; cada par
que satisfizer o limiar formará uma aresta; e cada componente conexo receberá
um único \texttt{duplicate\_group\_id}. Assim, se $A$ for semelhante a $B$ e
$B$ a $C$, os três permanecerão juntos mesmo quando a similaridade direta entre
$A$ e $C$ for inferior a 0,90.

Quando houver identificador confiável de conversa, tópico ou encadeamento, todos
os textos do mesmo contexto também serão mantidos no mesmo grupo. Considera-se
contexto compartilhado a mesma postagem original e seus comentários, a mesma
thread, o mesmo encadeamento de respostas ou uma conversa identificada pelo
mesmo ID de origem. Um \texttt{context\_group\_id} prevalecerá sobre a decisão
individual de partição. Essa regra reduz o risco de um comentário no treinamento
revelar o conteúdo de outro comentário relacionado presente no teste.

\subsection{Comprimento das sequências}

Antes de definir o truncamento, será calculada, separadamente para cada tokenizer,
a distribuição do comprimento dos textos. Serão reportados média, mediana,
percentis 90, 95 e 99, máximo e as proporções $P(L>128)$, $P(L>256)$ e
$P(L>512)$. O limite técnico superior será de 512 tokens, incluindo os tokens
especiais do encoder.

Três políticas serão comparadas no desenvolvimento: (A) primeiros 510 tokens de
conteúdo; (B) 255 tokens iniciais e 255 finais; e (C) janelas de até 510 tokens
com sobreposição, agregadas no nível do documento. Para a política C, tamanho da
janela, sobreposição e regra de agregação serão fixados antes do teste. Se menos
de 5\% dos textos excederem 512 tokens e não houver diferença relevante no
desenvolvimento, a política B será adotada por parcimônia. Caso contrário, a
política será tratada como hiperparâmetro de ablação. A quantidade de textos
truncados, o conteúdo médio descartado e o custo adicional do \textit{chunking}
serão registrados.

\section{Definição das tarefas de aprendizagem}
\label{sec:tarefas-aprendizagem}

\subsection{Tarefa A: detecção binária}

Para um texto $x_i$, o alvo binário será definido por

\begin{equation}
    y_i =
    \begin{cases}
        1, & \text{se o rótulo gold for \texttt{hate}},\\
        0, & \text{se o rótulo gold for \texttt{nao\_hate}}.
    \end{cases}
\end{equation}

O vetor contextual associado ao token de classificação será fornecido a uma
camada linear com duas saídas, seguida por \textit{softmax}. A função de perda
será a entropia cruzada. Serão comparadas duas condições: perda sem ponderação e
perda ponderada por classe. Na condição ponderada, o peso da classe $c$ será
calculado apenas no subconjunto de treinamento de cada fold:

\begin{equation}
    \alpha_c = \frac{N_{\mathrm{treino}}}{2N_c},
\end{equation}

em que $N_c$ é o número de exemplos da classe $c$ no fold de treinamento.

\subsection{Tarefa B: classificação multirrótulo}
\label{sec:tarefa-multirrotulo}

A tarefa multirrótulo utilizará somente os textos com rótulo gold
\texttt{hate}. Para cada texto, será construído um vetor binário
$\mathbf{y}_i \in \{0,1\}^{10}$, na seguinte ordem fixa:

\begin{enumerate}
    \item \texttt{ageism};
    \item \texttt{aporophobia};
    \item \texttt{body\_shame};
    \item \texttt{capacitism};
    \item \texttt{lgbtphobia};
    \item \texttt{political};
    \item \texttt{racism};
    \item \texttt{religious\_intolerance};
    \item \texttt{sexism};
    \item \texttt{xenophobia}.
\end{enumerate}

No banco de origem, o rótulo de sexismo ainda está armazenado sob o código
histórico \texttt{misogyny}. Durante a preparação da tarefa multirrótulo, esse
valor será mapeado deterministicamente para \texttt{sexism}, sem alterar o
arquivo bruto. A tabela de mapeamento e sua versão integrarão o manifesto do
dataset.

O cabeçote terá dez saídas independentes com ativação sigmoide. A perda será a
entropia cruzada binária com logits, implementada por
\texttt{BCEWithLogitsLoss}. Também serão comparadas uma condição sem ponderação
e outra com peso positivo por categoria. Para a categoria $\ell$, o peso será
calculado exclusivamente no fold de treinamento:

\begin{equation}
    p_{\ell} = \frac{N_{\ell}^{-}}{N_{\ell}^{+}},
\end{equation}

em que $N_{\ell}^{+}$ e $N_{\ell}^{-}$ são, respectivamente, as quantidades de
exemplos positivos e negativos para o rótulo $\ell$.

Uma categoria será considerada elegível para inferência confirmatória somente
se possuir pelo menos 100 ocorrências na v0.1 e pelo menos 20 ocorrências no
conjunto de teste. Esses valores são salvaguardas operacionais mínimas para
evitar folds praticamente vazios; não garantem, sozinhos, precisão inferencial.
Por isso, serão apresentados intervalos de confiança e uma análise de
sensibilidade exigindo 20, 30 e 50 positivos no teste. Categorias cujo intervalo
de 95\% para F1 permaneça excessivamente amplo serão descritas como exploratórias,
mesmo que ultrapassem o limite mínimo.

\section{Arquiteturas e propagação de erros}
\label{sec:arquiteturas}

O experimento principal comparará duas organizações sob o mesmo encoder,
partições, pré-processamento e orçamento de ajuste:

\begin{enumerate}
    \item \textbf{hierárquica}: um classificador binário é executado primeiro e
    o classificador multirrótulo somente produz tipos quando o primeiro estágio
    prevê \texttt{hate};
    \item \textbf{conjunta não hierárquica}: um único encoder possui onze saídas
    sigmoides, uma para \texttt{hate} e dez para os tipos. Registros
    \texttt{nao\_hate} recebem zero nos dez tipos, e nenhuma regra de bloqueio é
    aplicada antes da medição. Inconsistências, como prever um tipo sem prever
    \texttt{hate}, serão contabilizadas.
\end{enumerate}

Na arquitetura conjunta, a perda será a soma ponderada da perda binária e da
perda dos tipos,

\begin{equation}
    \mathcal{L}_{\mathrm{conjunta}} =
    \mathcal{L}_{\mathrm{binária}} +
    \gamma\mathcal{L}_{\mathrm{tipos}},
\end{equation}

com $\gamma$ definido exclusivamente no desenvolvimento. Para limitar o custo,
essa ablação será executada inicialmente com o backbone vencedor da tarefa
binária; uma extensão aos demais backbones será identificada como análise
secundária.

O classificador de tipos será avaliado tanto com a entrada binária gold
(condição \textit{oracle}) quanto no pipeline completo. A propagação de erros
será quantificada por

\begin{equation}
    \Delta F1_{\mathrm{propagação}} =
    F1_{\mathrm{oracle}}-F1_{\mathrm{ponta\ a\ ponta}},
\end{equation}

globalmente, por categoria e por plataforma. A comparação hierárquica versus
conjunta utilizará predições pareadas no mesmo teste; qualquer vantagem deverá
ser interpretada junto do custo computacional e da taxa de inconsistências.

\section{Particionamento e desenho de validação}
\label{sec:particionamento}

\subsection{Teste fixo e conjunto de desenvolvimento}

Após a formação dos grupos de duplicatas, 20\% da v0.1 serão reservados como
conjunto de teste fixo. Dos 80\% restantes, 12,5\% serão separados previamente
como conjunto de calibração, equivalentes a 10\% da base completa. Os 70\%
restantes formarão o conjunto de seleção e ajuste. A proporção global será,
portanto, 70\% para seleção/treinamento, 10\% para calibração e limiares e 20\%
para teste. Todas as divisões respeitarão os grupos e a estratificação.

O conjunto de calibração terá três usos previamente limitados: estimar o
parâmetro de calibração, definir os limiares da configuração já selecionada e
medir a estabilidade dessas escolhas. Ele não será usado para escolher backbone,
taxa de aprendizado, perda, época ou política de truncamento. O teste será
armazenado separadamente e não poderá participar de:

\begin{itemize}
    \item ajuste de hiperparâmetros;
    \item seleção de modelo ou de época;
    \item definição de pesos de classe;
    \item escolha de limiares de decisão;
    \item calibração de probabilidades ou escolha de seus métodos;
    \item análise exploratória destinada a alterar o pipeline.
\end{itemize}

A divisão será realizada com estratificação iterativa multirrótulo, abordagem
destinada a conservar distribuições de rótulos em problemas nos quais uma mesma
instância pode pertencer a várias classes \cite{sechidis2011stratification}. O
vetor de estratificação combinará o rótulo binário, os dez tipos de ódio, a
plataforma, o tipo de conteúdo e o estrato de amostragem. A unidade indivisível
será o grupo de duplicatas, e não o registro individual.

\subsection{Validação cruzada}

O conjunto de seleção e ajuste será dividido em cinco folds por meio da mesma
estratégia de agrupamento e estratificação. Em cada rodada, quatro folds serão
usados para treinamento e um para validação. Ao final das cinco rodadas, cada
registro desse conjunto terá exatamente uma predição fora da amostra,
denominada \textit{out-of-fold} (OOF).

As predições OOF serão usadas para comparar configurações e estimar sua
variabilidade sem consultar o conjunto de calibração. A métrica OOF global,
calculada após concatenar as cinco predições de validação, será o critério
principal de seleção. A média e o desvio-padrão das cinco métricas calculadas
separadamente também serão reportados, mas não substituirão o critério global.
Essa escolha é especialmente importante porque o
fine-tuning de modelos pré-treinados pode apresentar variação relevante em
função da inicialização e da ordem dos dados, sobretudo em conjuntos menores
\cite{dodge2020finetuning}.

\subsection{Correção do desenho amostral}

Os lotes de avaliação foram balanceados a partir de informações como plataforma,
rótulo previsto pelo Gemma e tipo previsto de ódio. Logo, a distribuição da
base anotada pode diferir da distribuição do universo completo. O manifesto de
amostragem registrará o universo disponível em cada lote, definição e tamanho
dos estratos, quantidade selecionada, algoritmo, semente, critérios de exclusão,
uso das previsões do Gemma e qualquer priorização manual. Para cada estrato $h$,
será calculado o peso de expansão

\begin{equation}
    w_h = \frac{N_h}{n_h},
\end{equation}

em que $N_h$ é o número de registros elegíveis do estrato no universo e $n_h$
é o número de registros selecionados do mesmo estrato. A expressão corresponde
ao inverso da probabilidade de inclusão somente quando a seleção é conhecida e
equiprovável dentro de cada estrato. Para estabilidade numérica, os pesos serão
normalizados para média 1 dentro de cada partição.

As métricas do teste serão apresentadas em duas versões: não ponderada, que
descreve diretamente o conjunto anotado, e ponderada, que busca estimar o
desempenho sob a composição do universo original. Essa segunda interpretação
depende das hipóteses de inclusão conhecível, positividade e seleção ignorável
dentro dos estratos. Seleção manual, revisão dirigida ou exclusão subjetiva
viola essas hipóteses e não é corrigida automaticamente por $w_h$; nessa situação,
a versão ponderada será qualificada como ajuste descritivo, não como estimativa
representativa da população. Os pesos não serão calculados com base no rótulo
gold do teste e jamais serão derivados de seus erros.

Sempre que operacionalmente possível, será sorteada do universo uma amostra
probabilística independente das previsões do Gemma. Ela será anotada pelo mesmo
protocolo e usada como teste externo da prevalência, dos pesos e da generalização.
Se essa amostra não existir na v0.1, a ausência será registrada como limitação e
sua construção será prioridade da versão seguinte.

\section{Modelos comparados}
\label{sec:modelos-comparados}

Cinco encoders serão submetidos ao mesmo desenho experimental:

\begin{description}
    \item[BERTimbau Base:] \texttt{neuralmind/bert-base-portuguese-cased},
    modelo BERT Base monolíngue, sensível à caixa e pré-treinado para português
    brasileiro \cite{souza2020bertimbau,neuralmind2026bertimbau};
    \item[BERTimbau Large:] \texttt{neuralmind/bert-large-portuguese-cased},
    versão de maior capacidade do mesmo projeto
    \cite{souza2020bertimbau,neuralmind2026bertimbau};
    \item[BERTabaporu Base:]
    \texttt{pablocosta/bertabaporu-base-uncased}, modelo BERT Base monolíngue
    sem distinção de caixa, pré-treinado em aproximadamente 238 milhões de
    tweets em português brasileiro, provenientes de mais de 100 mil usuários e
    totalizando cerca de 2,9 bilhões de tokens
    \cite{dacosta2023bertabaporu,dacosta2026bertabaporu};
    \item[mBERT:] \texttt{google-bert/bert-base-multilingual-cased}, modelo BERT
    multilíngue sensível à caixa, pré-treinado em 104 idiomas
    \cite{devlin2019bert,google2026mbert};
    \item[XLM-R Base:] \texttt{FacebookAI/xlm-roberta-base}, encoder
    multilíngue baseado em RoBERTa, incluído como comparador da família de
    Transformers bidirecionais derivados de BERT
    \cite{conneau2020xlmr,meta2026xlmr}.
\end{description}

O BERTabaporu permitirá testar diretamente o efeito do gênero de pré-treinamento:
BERTimbau Base representa o português geral da web, enquanto BERTabaporu Base
representa a linguagem de redes sociais. Como a v0.1 é multiplataforma, eventual
vantagem no X/Twitter não será generalizada para as demais fontes. Serão
reportados o resultado global, o resultado por plataforma e a interação entre
modelo e plataforma. A versão
\texttt{pablocosta/bertabaporu-large-uncased} ficará como extensão secundária,
condicionada ao orçamento computacional, e não substituirá a execução do Base
no benchmark principal.

Os checkpoints e tokenizers serão carregados por meio da biblioteca
Transformers \cite{wolf2020transformers}. A revisão exata de cada repositório
será fixada por hash de commit no início dos experimentos, evitando que uma
atualização remota altere silenciosamente os resultados.

\subsection{Baselines}

Três famílias de baselines serão mantidas:

\begin{enumerate}
    \item \textbf{baseline trivial}: previsão constante de \texttt{nao\_hate}
    na tarefa binária e previsão determinada apenas pela prevalência de cada
    categoria, estimada no treinamento, na tarefa multirrótulo;
    \item \textbf{TF--IDF com regressão logística}: concatenação de n-gramas de
    palavras de ordem 1--2 e n-gramas de caracteres de ordem 3--5, seguida por
    regressão logística. O parâmetro $C$ será selecionado entre 0,1, 1 e 10 no
    mesmo protocolo de cinco folds. Na tarefa multirrótulo será usado
    \texttt{OneVsRestClassifier}, com um classificador logístico independente
    para cada um dos dez tipos;
    \item \textbf{Gemma}: uso das predições originais registradas antes da
    revisão humana, sem novo ajuste. Esse baseline permitirá comparar os BERTs
    com o sistema responsável pelas classificações iniciais.
\end{enumerate}

O baseline TF--IDF será submetido às mesmas partições, agrupamentos e métricas.
Nenhum vocabulário ou transformador TF--IDF será ajustado fora do fold de
treinamento. Como as previsões do Gemma ajudaram a definir estratos e prioridades
de revisão, o desempenho desse baseline não será interpretado como proveniente
de uma amostra independente do próprio sistema. A possível dependência entre a
amostragem e seus erros será explicitada em todas as comparações.

\section{Protocolo de fine-tuning}
\label{sec:fine-tuning}

Todos os encoders receberão cabeçotes de classificação novos, inicializados
aleatoriamente. Os parâmetros do encoder e do cabeçote serão ajustados em
conjunto. A faixa inicial segue a prática de fine-tuning apresentada para BERT,
enquanto decisões adicionais serão verificadas em piloto no próprio domínio
\cite{devlin2019bert}. O protocolo comum será:

\begin{itemize}
    \item otimizador AdamW, com $\beta_1=0{,}9$, $\beta_2=0{,}999$ e
    $\epsilon=10^{-8}$;
    \item tamanho efetivo de lote igual a 32, obtido por acúmulo de gradientes
    quando necessário;
    \item limite superior de oito épocas, entendido como teto operacional e não
    como duração fixa;
    \item avaliação ao final de cada época e parada antecipada após duas épocas
    completas sem redução da perda de validação;
    \item seleção do checkpoint de menor perda de validação, antes da otimização
    de qualquer limiar;
    \item aquecimento linear durante os 10\% primeiros passos;
    \item decaimento linear da taxa de aprendizado após o aquecimento;
    \item norma máxima do gradiente igual a 1,0;
    \item precisão mista FP16 ou BF16 quando suportada pelo hardware;
    \item embaralhamento reprodutível dos exemplos a cada época.
\end{itemize}

A perda de validação foi escolhida para a parada antecipada porque independe de
um limiar de classificação; Macro-F1, F1-\texttt{hate} e AUPRC serão monitorados,
mas não mudarão o checkpoint durante essa etapa. O teto de oito épocas apenas
permite que a regra de paciência opere em tarefas multirrótulo de convergência
mais lenta. O artigo original de BERT usou buscas mais curtas, tipicamente entre
duas e quatro épocas \cite{devlin2019bert}; portanto, será realizado um piloto
com curvas de treinamento. Se as curvas mostrarem convergência sistemática antes
do teto, o custo será reduzido pela parada antecipada. Se permanecerem em melhora
na oitava época, qualquer ampliação exigirá emenda prévia ao protocolo.

Será avaliada a grade cartesiana formada por

\begin{equation}
    \eta \in \{10^{-5}, 2\times10^{-5}, 3\times10^{-5}, 5\times10^{-5}\}
\end{equation}

e

\begin{equation}
    \lambda \in \{0, 0{,}01\},
\end{equation}

em que $\eta$ representa a taxa de aprendizado e $\lambda$ o decaimento de
peso. O valor $10^{-5}$ foi acrescentado às taxas usuais de BERT para contemplar
ajustes mais conservadores, especialmente no modelo Large. A condição de perda
comum ou ponderada será um fator experimental controlado, e não uma decisão
tomada depois de observar o teste.

\subsection{Busca de hiperparâmetros em duas etapas}

Uma grade completa com cinco modelos, duas tarefas, oito combinações, cinco
folds e duas perdas produziria 800 fine-tunings antes das repetições por semente.
Para reduzir custo sem escolher configurações pelo teste, será usada a seguinte
triagem pré-especificada:

\begin{enumerate}
    \item para cada combinação de backbone, tarefa e condição de perda, as oito
    combinações de $\eta$ e $\lambda$ serão executadas por até duas épocas em uma
    divisão interna fixa do conjunto de seleção, com semente 42;
    \item as duas melhores combinações de cada célula serão continuadas até o
    teto de oito épocas, sujeitas à mesma parada antecipada;
    \item a melhor combinação de cada célula seguirá para os cinco folds. Se as
    duas primeiras diferirem menos de 0,005 em Macro-F1, ambas seguirão para
    evitar descarte baseado em ruído da triagem;
    \item a configuração será selecionada pela métrica OOF global, obtida pela
    concatenação das predições dos cinco folds com limiar 0,5. Somente depois
    dessa seleção serão definidos calibração e limiares.
\end{enumerate}

Sem empates, a triagem compreende 160 execuções curtas, 40 continuações e 100
ajustes em folds, em lugar de 800 ajustes completos. Antes da execução será
estimado o orçamento em GPU-horas a partir do piloto, incluindo as repetições
finais e a ablação de arquitetura. Se o orçamento disponível for insuficiente,
o escopo será reduzido por regra registrada antes do teste, priorizando as
perguntas de pesquisa, e não por resultados parciais convenientes.

Na tarefa binária, Macro-F1 será o critério de seleção; F1-\texttt{hate} e AUPRC
serão critérios secundários. Na tarefa multirrótulo, será usado Macro-F1 das
categorias previamente elegíveis. Empates inferiores a 0,005 serão resolvidos,
nesta ordem, por menor variância entre folds, menor perda OOF e menor taxa de
aprendizado. Número de parâmetros e tempo de inferência não serão usados para
desempatar taxas do mesmo backbone, pois permanecem essencialmente constantes.

\subsection{Repetições e definição do modelo final}

Após a seleção, cada condição experimental será treinada no conjunto de seleção
e ajuste com as sementes 13, 42, 97, 123 e 2026. O número de épocas será fixado
na mediana das melhores épocas dos cinco folds, evitando parada antecipada no
teste ou no conjunto de calibração. As probabilidades dos cinco modelos serão
preservadas para medir variabilidade de otimização.

Para a avaliação principal de cada condição, o preditor final será o ensemble
pré-especificado das cinco sementes, formado pela média dos logits antes da
calibração. Nenhuma semente será escolhida pelo desempenho no teste. Caso seja
necessário disponibilizar também um único checkpoint por restrição operacional,
a semente 42 será usada por definição prévia e seus resultados serão reportados
separadamente dos resultados do ensemble.

\section{Limiarização e calibração}
\label{sec:limiarizacao-calibracao}

Na tarefa binária, o ensemble produzirá um logit médio e uma probabilidade
estimada de \texttt{hate}. Depois de fixados backbone, hiperparâmetros, perda,
época e sementes, o limiar $\tau$ será escolhido exclusivamente no conjunto de
calibração. Serão testados valores entre 0,01 e 0,99 em incrementos de 0,01,
selecionando-se o valor que maximizar o Macro-F1. Em empate, será escolhido o
limiar mais próximo de 0,5 e, persistindo o empate, o maior limiar. O valor 0,5
também será reportado como referência.

Na tarefa multirrótulo, cada categoria $\ell$ terá um limiar próprio
$\tau_{\ell}$, escolhido na mesma grade e no mesmo conjunto para maximizar o F1
da categoria. Se houver menos de cinco exemplos positivos de uma categoria na
calibração, seu limiar será mantido em 0,5 e o resultado será tratado como
exploratório. A estabilidade será descrita reestimando os limiares em
reamostragens agrupadas da calibração e, apenas como diagnóstico, nas validações
dos cinco folds; serão reportados mediana, intervalo interquartil e amplitude.

As probabilidades da tarefa binária serão adicionalmente calibradas por
\textit{temperature scaling} ajustado nos logits médios das cinco sementes no
conjunto reservado de calibração. Essa separação evita transferir diretamente
um calibrador estimado sobre logits OOF, produzidos por modelos distintos, para
um modelo final reajustado. Redes neurais modernas podem produzir probabilidades
excessivamente confiantes, e o método oferece uma correção pós-processada simples
\cite{guo2017calibration}.

A temperatura não ponderada estimará calibração para a distribuição da amostra
anotada. Uma temperatura ponderada, obtida pela minimização da log-verossimilhança
negativa com pesos de inclusão, será estimada para a população somente se a
auditoria do desenho amostral validar esses pesos. Caso contrário, não será feita
alegação de que $P(\texttt{hate}\mid\text{amostra})$ corresponde a
$P(\texttt{hate}\mid\text{população})$. O teste receberá apenas transformações e
limiares já congelados. ECE será calculado com 15 intervalos de frequência
aproximadamente igual; também serão apresentados Brier score e diagramas de
confiabilidade antes e depois da calibração.

\section{Métricas de avaliação}
\label{sec:metricas-avaliacao}

\subsection{Tarefa binária}

O Macro-F1 será a métrica primária de comparação geral. F1-\texttt{hate} e AUPRC
serão métricas centrais para as hipóteses sobre a classe positiva, mas não
substituirão retroativamente a métrica definida para cada contraste na
Tabela~\ref{tab:matriz-rq}. Para duas classes $C$, Macro-F1 será calculado
por

\begin{equation}
    \operatorname{MacroF1} = \frac{1}{|C|}\sum_{c\in C}
    \frac{2P_cR_c}{P_c+R_c},
\end{equation}

em que $P_c$ e $R_c$ representam precisão e revocação da classe $c$. O uso de
Macro-F1 impede que o desempenho da classe majoritária domine sozinho a
comparação.

Também serão reportados:

\begin{itemize}
    \item precisão, revocação e F1 da classe \texttt{hate};
    \item área sob a curva Precisão--Revocação (AUPRC);
    \item área sob a curva ROC (AUROC);
    \item coeficiente de correlação de Matthews (MCC);
    \item acurácia balanceada;
    \item matriz de confusão;
    \item Brier score;
    \item erro esperado de calibração (ECE);
    \item curvas de confiabilidade antes e depois da calibração.
\end{itemize}

A AUPRC receberá destaque porque a classe positiva é minoritária. Em cenários
desbalanceados, curvas Precisão--Revocação são mais informativas sobre a
qualidade das predições positivas do que a curva ROC isoladamente
\cite{saito2015precisionrecall}.

\subsection{Tarefa multirrótulo}

A métrica primária será o Macro-F1 calculado sobre as categorias elegíveis para
análise confirmatória. Serão apresentados, adicionalmente:

\begin{itemize}
    \item Macro-F1 incluindo todos os rótulos ativos;
    \item Micro-F1;
    \item precisão, revocação, F1 e PR-AUC por categoria;
    \item \textit{Hamming loss};
    \item \textit{subset accuracy}, que exige correspondência exata do conjunto
    de rótulos;
    \item média da quantidade de rótulos previstos por texto;
    \item matriz de coocorrência entre tipos verdadeiros e previstos.
\end{itemize}

Serão realizadas duas avaliações complementares. A avaliação \textit{oracle}
aplicará o classificador de tipos somente aos registros cujo rótulo binário gold
seja \texttt{hate}, medindo a capacidade específica de reconhecer categorias.
A avaliação ponta a ponta executará primeiro o classificador binário e somente
atribuirá tipos aos textos previstos como positivos. A diferença entre os dois
resultados quantificará a propagação de erros da primeira etapa.

\subsection{Incerteza e comparação entre modelos}

Serão separadas duas fontes de incerteza. A variabilidade de otimização será
descrita pela média, desvio-padrão, mínimo e máximo entre as cinco sementes. A
incerteza amostral será estimada sobre a predição final pré-especificada por
2.000 reamostragens bootstrap do conjunto de teste \cite{efron1993bootstrap}. O
bootstrap será estratificado
pelos estratos de amostragem e terá como unidade o maior agrupamento entre
\texttt{duplicate\_group\_id} e \texttt{context\_group\_id}, preservando a
dependência entre textos relacionados. Na versão ponderada, os pesos serão
mantidos em cada reamostragem.

Diferenças entre modelos serão examinadas por bootstrap pareado, calculando
diretamente, em cada reamostragem, $\Delta M=M_A-M_B$ sobre os mesmos exemplos.
Não será inferida diferença pela simples sobreposição de intervalos individuais.
Os quatro contrastes primários da Tabela~\ref{tab:matriz-rq} serão apresentados
sem expansão posterior do conjunto de hipóteses. Para famílias de comparações
secundárias entre modelos, plataformas ou categorias, os valores de $p$ serão
ajustados pelo procedimento de Holm \cite{holm1979procedure}. Magnitude,
intervalo de confiança e relação
com $\delta$ serão priorizados sobre uma decisão binária de significância.

Além dos resultados globais, as métricas serão desagregadas por plataforma,
tipo de conteúdo e categoria de ódio. Subgrupos com menos de 30 exemplos
positivos não receberão interpretação comparativa; seus valores serão exibidos
apenas como evidência exploratória, acompanhados do tamanho da amostra e do
intervalo de confiança.

Para avaliar generalização entre plataformas, será realizado um experimento
\textit{leave-one-platform-out}: em cada rodada, uma plataforma será mantida
inteiramente fora do ajuste e usada como domínio de teste. A rodada somente será
confirmatória quando a plataforma excluída possuir ao menos 100 registros e 20
positivos; nos demais casos será exploratória. Os resultados serão reportados
para YouTube, WhatsApp, Telegram, Twitter/X, Facebook, Reddit e TikTok, sem
combinar artificialmente plataformas pequenas para obter significância.

Também serão registrados parâmetros, GPU-horas, memória máxima, duração do
treinamento e tempo de inferência por mil textos. A interpretação final comparará
ganho de desempenho e custo, evitando concluir que um modelo maior é preferível
quando a diferença estiver abaixo de $\delta$ e o custo for substancialmente
superior.

\section{Testes funcionais e análise de erros}
\label{sec:testes-funcionais}

Métricas agregadas em um conjunto de teste podem ocultar falhas sistemáticas de
modelos de detecção de discurso de ódio. O HateCheck demonstra a utilidade de
testes funcionais construídos para fenômenos linguísticos específicos
\cite{rottger2021hatecheck}. Com base nessa abordagem, será produzido um pequeno
conjunto diagnóstico em português brasileiro, separado de treinamento,
validação e teste estatístico.

O conjunto funcional terá, como meta inicial, pelo menos 20 exemplos por
fenômeno, equilibrando casos positivos e negativos quando isso for semanticamente
aplicável. Os exemplos poderão ser adaptados de taxonomias funcionais publicadas,
construídos de forma controlada ou extraídos do domínio, desde que origem e
licença sejam registradas. O conjunto deverá contemplar, no mínimo:

\begin{itemize}
    \item negação de uma afirmação preconceituosa;
    \item citação de discurso de ódio para denúncia ou análise;
    \item contra-discurso e defesa do grupo atacado;
    \item ironia e sarcasmo;
    \item menção neutra a grupos sociais;
    \item insulto individual sem ataque generalizante a um grupo;
    \item grafias alteradas, separação de caracteres e uso de símbolos;
    \item ameaça, incitação, desumanização e criminalização coletiva;
    \item ataques explícitos e implícitos;
    \item variações de gênero, número e referência pronominal.
\end{itemize}

Os casos serão classificados por dois avaliadores independentes e divergências
serão adjudicadas pelo protocolo da Seção~\ref{sec:qualidade-anotacao}. Antes do
uso, cada caso será comparado com treino, calibração e teste por hash, Jaccard e
busca aproximada; casos equivalentes serão removidos ou reescritos e novamente
validados. O conjunto funcional não será usado para ajustar hiperparâmetros. O
desempenho será relatado por fenômeno, sem misturá-lo à estimativa principal do
teste da v0.1.

Após a avaliação quantitativa, será conduzida análise qualitativa de erros. Para
cada modelo serão examinados, no mínimo, 50 falsos positivos e 50 falsos
negativos da tarefa binária, sorteados com semente 2026 e estratificação por
plataforma e faixa de confiança. Na tarefa multirrótulo serão sorteados erros das
categorias mais frequentes e incluídos os erros disponíveis das categorias raras,
respeitando as restrições de privacidade.

Antes de examinar as predições será congelada a seguinte taxonomia inicial:
\texttt{E01} contexto ausente; \texttt{E02} ironia; \texttt{E03} sarcasmo;
\texttt{E04} negação; \texttt{E05} alvo implícito; \texttt{E06} ambiguidade
lexical; \texttt{E07} linguagem codificada; \texttt{E08} truncamento;
\texttt{E09} possível erro de anotação; \texttt{E10} conflito entre categorias;
\texttt{E11} conhecimento externo necessário; e \texttt{E12} dependência
multimodal. Novos códigos poderão emergir, mas serão identificados como
pós-especificados. Uma parcela mínima de 20\% dos erros será codificada por dois
avaliadores, com concordância reportada.

Cada item indicará se depende de imagem, vídeo, mensagem anterior, thread,
evento externo ou conhecimento sociopolítico. Será reportada a proporção de
erros potencialmente insolúveis pelo texto isolado, distinguindo limitação do
modelo de limitação da modalidade de entrada.

\section{Reprodutibilidade, documentação e uso responsável}
\label{sec:reprodutibilidade}

Cada execução deverá gerar um manifesto contendo:

\begin{itemize}
    \item versão e hash SHA-256 do dataset;
    \item hash do código e da consulta de extração;
    \item commit do repositório, estado de alterações locais e comando executado;
    \item arquivos imutáveis de \texttt{train\_ids}, \texttt{validation\_ids},
    \texttt{calibration\_ids}, \texttt{test\_ids},
    \texttt{duplicate\_groups} e \texttt{context\_groups};
    \item checkpoint e revisão do tokenizer;
    \item versões de Python, PyTorch, Transformers e demais dependências;
    \item modelo, tarefa, hiperparâmetros e semente;
    \item número de épocas, melhor época nos folds e curvas de perda;
    \item hardware, precisão numérica, memória máxima, GPU-horas e tempo de
    treinamento e inferência;
    \item limiares, temperatura de calibração e pesos de classe;
    \item métricas globais, ponderadas e por subgrupo;
    \item caminho e hash dos checkpoints resultantes.
\end{itemize}

O ambiente de software será congelado por arquivo de dependências com versões
exatas e, preferencialmente, por imagem de contêiner. Não se dependerá apenas da
semente para reconstruir partições. As configurações de determinismo do PyTorch,
CUDA e cuDNN serão registradas, incluindo flags ativadas e operações conhecidas
como não determinísticas. Quando determinismo estrito reduzir de modo relevante
o desempenho computacional ou não estiver disponível, a exceção será documentada.

A documentação do dataset descreverá origem, composição, processo de amostragem,
população representada, critérios de anotação, limitações e usos não recomendados.
Esse procedimento é coerente com a proposta de \textit{data statements}, que
busca tornar explícitas as condições sob as quais resultados de PLN podem ou não
ser generalizados \cite{bender2018datastatements}. Cada modelo terá também uma
ficha com finalidade, desempenho, grupos avaliados, limitações, riscos e uso
pretendido, seguindo a proposta de \textit{model cards}
\cite{mitchell2019modelcards}.

Os classificadores não deverão ser descritos como substitutos da avaliação
humana. Seu uso previsto é priorizar, apoiar ou acelerar a triagem de conteúdos.
Decisões de alto impacto, estudos de prevalência e análises sobre grupos pouco
representados deverão considerar a incerteza, o desenho amostral e a revisão de
especialistas.

Exemplos divulgados em artigos, dissertação ou fichas de modelo passarão por
revisão específica de risco de reidentificação. Conforme o risco, serão usados
mascaramento adicional, paráfrase fiel, exemplo sintético controlado ou omissão.
Identificadores de origem e textos brutos não acompanharão os artefatos públicos.

\section{Continuidade do processo}
\label{sec:continuidade-processo}

O teste da v0.1 permanecerá congelado nas versões posteriores, permitindo medir
ganhos ou regressões no mesmo conjunto histórico. Contudo, após ter sido
consultado em sucessivos ciclos, ele será interpretado como benchmark histórico
ou teste de regressão, e não indefinidamente como única avaliação confirmatória.
Novos registros avaliados serão incorporados ao desenvolvimento das versões
v0.2, v0.3 e seguintes. Cada versão reservará um teste temporal formado por
conteúdos estritamente posteriores ao período de treinamento; o teste temporal
mais recente assumirá progressivamente o papel confirmatório.

Para cada atualização serão repetidos os controles de elegibilidade,
duplicidade, estratificação, treinamento e avaliação. A data de corte temporal
será definida antes do treinamento e os grupos de contexto não poderão atravessá-la.
Um modelo novo somente substituirá o anterior se não apresentar regressão
relevante no teste histórico e demonstrar desempenho aceitável no teste temporal
mais recente.

O pré-treinamento adaptativo em dados não rotulados do domínio poderá ser
investigado em uma etapa futura. Estudos anteriores mostram que a continuação do
pré-treinamento em textos do domínio ou da tarefa pode produzir ganhos mesmo em
cenários com poucos rótulos \cite{gururangan2020dontstop}. Contudo, esse
procedimento não integrará o benchmark principal da v0.1, que comparará os
checkpoints públicos originais sob o mesmo protocolo supervisionado.

\section{Estudos de ablação}
\label{sec:ablacoes}

As decisões de engenharia que possam alterar as conclusões serão transformadas
em contrastes controlados. As ablações pré-especificadas são:

\begin{itemize}
    \item perda comum versus perda ponderada;
    \item arquitetura hierárquica versus conjunta;
    \item encoder monolíngue versus multilíngue;
    \item probabilidades brutas versus calibradas;
    \item truncamento apenas no início versus início e fim e, se necessário,
    \textit{chunking};
    \item texto anonimizado com marcadores especiais versus anonimização sem os
    marcadores explícitos.
\end{itemize}

Cada ablação alterará um componente por vez, manterá partições e sementes e será
analisada de forma pareada. Perda, idioma, calibração e propagação hierárquica
correspondem às perguntas primárias; truncamento e mascaramento serão análises
secundárias no backbone vencedor, salvo se o piloto mostrar efeito suficiente
para justificar expansão previamente documentada.

\section{Fluxo experimental consolidado}
\label{sec:fluxo-experimental}

O procedimento completo seguirá a ordem abaixo. A posição do teste no fluxo
representa isolamento lógico: seus rótulos somente serão acessados depois que os
artefatos anteriores estiverem congelados.

\begin{enumerate}
    \item definir universo, estratos, probabilidades de inclusão e lotes;
    \item concluir a anotação simples, auditar a dupla anotação histórica e
    documentar concordância, divergências e exclusões;
    \item congelar dataset, taxonomia e manual, com hashes e manifests;
    \item executar anonimização, normalização e testes de consistência;
    \item detectar duplicatas, formar componentes conexos e grupos de contexto;
    \item criar e versionar as partições de seleção, calibração e teste;
    \item perfilar comprimentos e executar apenas as ablações pré-especificadas;
    \item realizar triagem de hiperparâmetros e validação cruzada;
    \item produzir e concatenar predições OOF para seleção de configuração;
    \item fixar época e treinar as cinco sementes finais;
    \item formar o ensemble, ajustar calibração e limiares no conjunto próprio;
    \item congelar checkpoint, tokenizer, calibrador, limiares e código;
    \item abrir uma única vez o teste para a avaliação principal;
    \item executar bootstrap pareado, análises por plataforma e eficiência;
    \item executar avaliações externa, temporal e funcional quando disponíveis;
    \item conduzir análise qualitativa reprodutível e publicar a documentação.
\end{enumerate}

\section{Ameaças à validade}
\label{sec:ameacas-validade}

As principais ameaças à interpretação dos resultados são:

\begin{itemize}
    \item \textbf{viés de seleção}: a amostra foi balanceada com base em sinais
    do Gemma e pode não reproduzir diretamente a prevalência do universo;
    \item \textbf{revisão simples}: a predominância de decisões
    \texttt{single} impede estimar concordância entre avaliadores na maior parte
    da v0.1;
    \item \textbf{classes raras}: algumas categorias possuem poucas dezenas de
    exemplos e não sustentam conclusões confirmatórias;
    \item \textbf{mudança temporal}: linguagem ofensiva, códigos e plataformas
    podem mudar depois do congelamento;
    \item \textbf{dependência de contexto}: o texto isolado pode não conter
    informações presentes na conversa, imagem, vídeo ou evento relacionado;
    \item \textbf{propagação de erros}: no pipeline hierárquico, um falso
    negativo binário impede a identificação dos tipos de ódio;
    \item \textbf{generalização limitada}: os resultados serão válidos para a
    composição observada de plataformas, períodos e categorias, não para todo
    conteúdo em português brasileiro.
\end{itemize}

Essas limitações deverão acompanhar qualquer divulgação de métricas ou uso do
modelo. A expansão contínua do dataset será orientada também pelos erros e pelas
lacunas de cobertura detectadas na v0.1.

\section{Tabelas para registro dos resultados}
\label{sec:tabelas-resultados}

As tabelas seguintes definem o formato de apresentação, mas não contêm valores
experimentais. Os campos marcados com ``---'' deverão ser preenchidos somente
após a execução do protocolo.

\begin{table}[htbp]
    \centering
    \caption{Modelo de tabela para a tarefa binária no teste fixo.}
    \label{tab:modelo-resultados-binarios}
    \resizebox{\textwidth}{!}{%
    \begin{tabular}{lrrrrrrr}
        \toprule
        \textbf{Modelo} & \textbf{Macro-F1} & \textbf{F1 hate} &
        \textbf{Precisão hate} & \textbf{Revocação hate} & \textbf{AUPRC} &
        \textbf{MCC} & \textbf{ECE} \\
        \midrule
        TF--IDF + regressão logística & --- & --- & --- & --- & --- & --- & --- \\
        Gemma & --- & --- & --- & --- & --- & --- & --- \\
        BERTimbau Base & --- & --- & --- & --- & --- & --- & --- \\
        BERTimbau Large & --- & --- & --- & --- & --- & --- & --- \\
        BERTabaporu Base & --- & --- & --- & --- & --- & --- & --- \\
        mBERT & --- & --- & --- & --- & --- & --- & --- \\
        XLM-R Base & --- & --- & --- & --- & --- & --- & --- \\
        \bottomrule
    \end{tabular}}
    \par\smallskip
    {\footnotesize Fonte: elaboração própria. Valores a serem preenchidos após
    os experimentos.\par}
\end{table}

\begin{table}[htbp]
    \centering
    \caption{Modelo de tabela para a tarefa multirrótulo.}
    \label{tab:modelo-resultados-multirrotulo}
    \resizebox{\textwidth}{!}{%
    \begin{tabular}{lrrrrrr}
        \toprule
        \textbf{Modelo} & \textbf{Macro-F1 elegível} &
        \textbf{Macro-F1 total} & \textbf{Micro-F1} & \textbf{PR-AUC} &
        \textbf{Hamming loss} & \textbf{Subset accuracy} \\
        \midrule
        TF--IDF + regressão logística & --- & --- & --- & --- & --- & --- \\
        BERTimbau Base & --- & --- & --- & --- & --- & --- \\
        BERTimbau Large & --- & --- & --- & --- & --- & --- \\
        BERTabaporu Base & --- & --- & --- & --- & --- & --- \\
        mBERT & --- & --- & --- & --- & --- & --- \\
        XLM-R Base & --- & --- & --- & --- & --- & --- \\
        \bottomrule
    \end{tabular}}
    \par\smallskip
    {\footnotesize Fonte: elaboração própria. Valores a serem preenchidos após
    os experimentos.\par}
\end{table}

\begin{table}[htbp]
    \centering
    \caption{Modelo de tabela para resultados por categoria de ódio.}
    \label{tab:modelo-resultados-categoria}
    \resizebox{\textwidth}{!}{%
    \begin{tabular}{lrrrrrr}
        \toprule
        \textbf{Categoria} & \textbf{$n$ teste} & \textbf{Precisão} &
        \textbf{Revocação} & \textbf{F1} & \textbf{IC95\% F1} & \textbf{PR-AUC} \\
        \midrule
        Ageísmo & --- & --- & --- & --- & --- & --- \\
        Classismo & --- & --- & --- & --- & --- & --- \\
        Gordofobia & --- & --- & --- & --- & --- & --- \\
        Capacitismo & --- & --- & --- & --- & --- & --- \\
        LGBTfobia & --- & --- & --- & --- & --- & --- \\
        Ódio político & --- & --- & --- & --- & --- & --- \\
        Racismo & --- & --- & --- & --- & --- & --- \\
        Intolerância religiosa & --- & --- & --- & --- & --- & --- \\
        Sexismo & --- & --- & --- & --- & --- & --- \\
        Xenofobia & --- & --- & --- & --- & --- & --- \\
        \bottomrule
    \end{tabular}}
    \par\smallskip
    {\footnotesize Fonte: elaboração própria. Categorias abaixo do limiar
    confirmatório deverão ser identificadas como exploratórias.\par}
\end{table}

\begin{table}[htbp]
    \centering
    \caption{Modelo de tabela para desempenho e eficiência computacional.}
    \label{tab:modelo-eficiencia}
    \resizebox{\textwidth}{!}{%
    \begin{tabular}{lrrrrr}
        \toprule
        \textbf{Modelo} & \textbf{Macro-F1} & \textbf{Parâmetros} &
        \textbf{GPU-horas} & \textbf{Memória máxima} &
        \textbf{Tempo/1.000 textos} \\
        \midrule
        BERTimbau Base & --- & --- & --- & --- & --- \\
        BERTimbau Large & --- & --- & --- & --- & --- \\
        BERTabaporu Base & --- & --- & --- & --- & --- \\
        mBERT & --- & --- & --- & --- & --- \\
        XLM-R Base & --- & --- & --- & --- & --- \\
        \bottomrule
    \end{tabular}}
    \par\smallskip
    {\footnotesize Fonte: elaboração própria. Valores a serem preenchidos após
    os experimentos.\par}
\end{table}

\section{Critérios de conclusão da etapa v0.1}
\label{sec:criterios-conclusao}

A etapa de treinamento da v0.1 será considerada metodologicamente concluída
quando:

\begin{enumerate}
    \item o primeiro lote completo que alcance ao menos 21.248 decisões válidas
    tiver sido congelado e versionado;
    \item a suficiência amostral tiver sido examinada e as análises confirmatórias
    e exploratórias estiverem separadas;
    \item a dupla anotação histórica, as medidas de concordância e as exclusões
    por divergência tiverem sido documentadas, e a ausência de dupla anotação
    contemporânea estiver explicitamente declarada como limitação;
    \item as partições tiverem sido verificadas quanto a duplicatas e vazamento;
    \item todos os baselines e cinco encoders tiverem sido executados nas mesmas
    divisões;
    \item os hiperparâmetros tiverem sido definidos por OOF e os limiares e a
    calibração pelo conjunto próprio, sem uso do teste;
    \item as cinco sementes tiverem sido executadas para cada configuração
    selecionada;
    \item o ensemble final tiver sido fixado sem seleção posterior de semente;
    \item as métricas globais, ponderadas, por plataforma e por categoria
    tiverem sido calculadas com intervalos de confiança e comparações pareadas;
    \item a análise funcional e a análise qualitativa de erros tiverem sido
    documentadas;
    \item os manifests, checkpoints, hashes, tabelas e fichas de documentação
    tiverem sido arquivados de forma reprodutível.
\end{enumerate}

% -----------------------------------------------------------------------------
% Referências incorporadas ao próprio capítulo
% -----------------------------------------------------------------------------

\begin{thebibliography}{99}

\bibitem[Bender e Friedman(2018)]{bender2018datastatements}
BENDER, Emily M.; FRIEDMAN, Batya.
Data statements for natural language processing: toward mitigating system bias
and enabling better science.
\textit{Transactions of the Association for Computational Linguistics}, v. 6,
p. 587--604, 2018.
DOI: \url{https://doi.org/10.1162/tacl_a_00041}.

\bibitem[Cohen(1960)]{cohen1960kappa}
COHEN, Jacob.
A coefficient of agreement for nominal scales.
\textit{Educational and Psychological Measurement}, v. 20, n. 1, p. 37--46,
1960.
DOI: \url{https://doi.org/10.1177/001316446002000104}.

\bibitem[Conneau et al.(2020)]{conneau2020xlmr}
CONNEAU, Alexis et al.
Unsupervised cross-lingual representation learning at scale.
In: \textit{Proceedings of the 58th Annual Meeting of the Association for
Computational Linguistics}. Online: Association for Computational Linguistics,
2020. p. 8440--8451.
Disponível em: \url{https://arxiv.org/abs/1911.02116}.

\bibitem[da Costa et al.(2023)]{dacosta2023bertabaporu}
DA COSTA, Pablo Botton; PAVAN, Matheus Camasmie; DOS SANTOS, Wesley Ramos;
DA SILVA, Samuel Caetano; PARABONI, Ivandre.
BERTabaporu: assessing a genre-specific language model for Portuguese NLP.
In: \textit{Proceedings of the 14th International Conference on Recent Advances
in Natural Language Processing}. Varna: INCOMA, 2023. p. 217--223.
DOI: \url{https://doi.org/10.26615/978-954-452-092-2_024}.

\bibitem[da Costa(2026)]{dacosta2026bertabaporu}
DA COSTA, Pablo Botton.
\textit{BERTabaporu Base Uncased: model card}.
Hugging Face, 2026. Disponível em:
\url{https://huggingface.co/pablocosta/bertabaporu-base-uncased}.
Acesso em: 11 set. 2026.

\bibitem[Devlin et al.(2019)]{devlin2019bert}
DEVLIN, Jacob; CHANG, Ming-Wei; LEE, Kenton; TOUTANOVA, Kristina.
BERT: pre-training of deep bidirectional transformers for language understanding.
In: \textit{Proceedings of the 2019 Conference of the North American Chapter of
the Association for Computational Linguistics: Human Language Technologies}.
Minneapolis: Association for Computational Linguistics, 2019. p. 4171--4186.
Disponível em: \url{https://arxiv.org/abs/1810.04805}.

\bibitem[Dodge et al.(2020)]{dodge2020finetuning}
DODGE, Jesse et al.
Fine-tuning pretrained language models: weight initializations, data orders,
and early stopping.
\textit{arXiv preprint arXiv:2002.06305}, 2020.
Disponível em: \url{https://arxiv.org/abs/2002.06305}.

\bibitem[Efron e Tibshirani(1993)]{efron1993bootstrap}
EFRON, Bradley; TIBSHIRANI, Robert J.
\textit{An introduction to the bootstrap}.
New York: Chapman \& Hall, 1993.
DOI: \url{https://doi.org/10.1007/978-1-4899-4541-9}.

\bibitem[Google BERT(2026)]{google2026mbert}
GOOGLE BERT.
\textit{BERT multilingual base model (cased): model card}.
Hugging Face, 2026.
Disponível em:
\url{https://huggingface.co/google-bert/bert-base-multilingual-cased}.
Acesso em: 1 set. 2026.

\bibitem[Guo et al.(2017)]{guo2017calibration}
GUO, Chuan; PLEISS, Geoff; SUN, Yu; WEINBERGER, Kilian Q.
On calibration of modern neural networks.
In: \textit{Proceedings of the 34th International Conference on Machine Learning}.
Sydney: PMLR, 2017. v. 70, p. 1321--1330.
Disponível em: \url{https://arxiv.org/abs/1706.04599}.

\bibitem[Holm(1979)]{holm1979procedure}
HOLM, Sture.
A simple sequentially rejective multiple test procedure.
\textit{Scandinavian Journal of Statistics}, v. 6, n. 2, p. 65--70, 1979.

\bibitem[Gururangan et al.(2020)]{gururangan2020dontstop}
GURURANGAN, Suchin et al.
Don't stop pretraining: adapt language models to domains and tasks.
In: \textit{Proceedings of the 58th Annual Meeting of the Association for
Computational Linguistics}. Online: Association for Computational Linguistics,
2020. p. 8342--8360.
DOI: \url{https://doi.org/10.18653/v1/2020.acl-main.740}.

\bibitem[Krippendorff(2018)]{krippendorff2018content}
KRIPPENDORFF, Klaus.
\textit{Content analysis: an introduction to its methodology}.
4. ed. Thousand Oaks: SAGE, 2018.

\bibitem[Meta AI(2026)]{meta2026xlmr}
META AI.
\textit{XLM-RoBERTa base: model card}.
Hugging Face, 2026.
Disponível em: \url{https://huggingface.co/FacebookAI/xlm-roberta-base}.
Acesso em: 1 set. 2026.

\bibitem[Mitchell et al.(2019)]{mitchell2019modelcards}
MITCHELL, Margaret et al.
Model cards for model reporting.
In: \textit{Proceedings of the Conference on Fairness, Accountability, and
Transparency}. New York: ACM, 2019. p. 220--229.
DOI: \url{https://doi.org/10.1145/3287560.3287596}.

\bibitem[NeuralMind(2026)]{neuralmind2026bertimbau}
NEURALMIND.
\textit{BERTimbau Base and Large: model cards}.
Hugging Face, 2026.
Disponível em:
\url{https://huggingface.co/neuralmind/bert-base-portuguese-cased} e
\url{https://huggingface.co/neuralmind/bert-large-portuguese-cased}.
Acesso em: 1 set. 2026.

\bibitem[Röttger et al.(2021)]{rottger2021hatecheck}
RÖTTGER, Paul; VIDGEN, Bertie; NGUYEN, Dong; WASEEM, Zeerak; MARGETTS, Helen;
PIERREHUMBERT, Janet.
HateCheck: functional tests for hate speech detection models.
In: \textit{Proceedings of the 59th Annual Meeting of the Association for
Computational Linguistics and the 11th International Joint Conference on
Natural Language Processing}. Online: Association for Computational Linguistics,
2021. p. 41--58.
DOI: \url{https://doi.org/10.18653/v1/2021.acl-long.4}.

\bibitem[Saito e Rehmsmeier(2015)]{saito2015precisionrecall}
SAITO, Takaya; REHMSMEIER, Marc.
The precision-recall plot is more informative than the ROC plot when evaluating
binary classifiers on imbalanced datasets.
\textit{PLOS ONE}, v. 10, n. 3, e0118432, 2015.
DOI: \url{https://doi.org/10.1371/journal.pone.0118432}.

\bibitem[Sechidis, Tsoumakas e Vlahavas(2011)]{sechidis2011stratification}
SECHIDIS, Konstantinos; TSOUMAKAS, Grigorios; VLAHAVAS, Ioannis.
On the stratification of multi-label data.
In: \textit{Machine Learning and Knowledge Discovery in Databases}.
Berlin: Springer, 2011. p. 145--158.
DOI: \url{https://doi.org/10.1007/978-3-642-23808-6_10}.

\bibitem[Souza, Nogueira e Lotufo(2020)]{souza2020bertimbau}
SOUZA, Fábio; NOGUEIRA, Rodrigo; LOTUFO, Roberto.
BERTimbau: pretrained BERT models for Brazilian Portuguese.
In: \textit{Intelligent Systems: 9th Brazilian Conference, BRACIS 2020}.
Cham: Springer, 2020. p. 403--417.
DOI: \url{https://doi.org/10.1007/978-3-030-61377-8_28}.

\bibitem[Wolf et al.(2020)]{wolf2020transformers}
WOLF, Thomas et al.
Transformers: state-of-the-art natural language processing.
In: \textit{Proceedings of the 2020 Conference on Empirical Methods in Natural
Language Processing: System Demonstrations}. Online: Association for
Computational Linguistics, 2020. p. 38--45.
DOI: \url{https://doi.org/10.18653/v1/2020.emnlp-demos.6}.

\end{thebibliography}

%
% Pacotes recomendados no preâmbulo do documento principal:
% \usepackage[T1]{fontenc}
% \usepackage[utf8]{inputenc} % dispensável em distribuições LaTeX recentes
% \usepackage[brazil]{babel}
% \usepackage{amsmath,amssymb}
% \usepackage{booktabs,longtable,array}
% \usepackage{graphicx}
% \usepackage{url}
% \usepackage[hidelinks]{hyperref}
% \usepackage[alf]{abntex2cite}
%
% As referências foram incorporadas ao final deste arquivo para que o capítulo
% não dependa de um arquivo .bib externo. Caso o documento principal já possua
% bibliografia própria, os itens podem ser transferidos para ela sem alterar as
% chaves usadas nas citações.

\chapter{Metodologia de treinamento dos modelos BERT}
\label{cap:metodologia-bert-v01}

\section{Visão geral e objetivos}
\label{sec:visao-geral-bert}

Este capítulo estabelece o protocolo para construção, treinamento, seleção e
avaliação dos primeiros classificadores supervisionados do projeto Radar Ódio
a partir do \textit{Dataset} v0.1. Essa versão será formada pelas avaliações
humanas resolvidas que correspondam a, no mínimo, 10\% do universo elegível. O
objetivo não é tratar a v0.1 como versão definitiva do conjunto de dados, mas
como o primeiro marco reprodutível de um processo contínuo de aprendizagem e
avaliação, a ser ampliado nas versões posteriores.

A estratégia emprega modelos pré-treinados da família BERT e derivados. O uso
de representações contextuais pré-treinadas permite adaptar um modelo de língua
a uma tarefa supervisionada mediante a adição de uma camada de classificação e
o ajuste conjunto dos parâmetros, procedimento conhecido como \textit{fine-tuning}
\cite{devlin2019bert}. Para o português brasileiro, os modelos BERTimbau
constituem referências monolíngues relevantes, tendo sido pré-treinados em um
grande corpus nacional e avaliados em diferentes tarefas de processamento de
linguagem natural \cite{souza2020bertimbau}. Modelos multilíngues serão incluídos
como controles comparativos, permitindo verificar se a especialização no idioma
produz vantagem no domínio de textos oriundos de redes sociais.

O protocolo contempla duas tarefas relacionadas, porém treinadas e avaliadas
separadamente:

\begin{enumerate}
    \item \textbf{detecção binária de discurso de ódio}: classificação de cada
    texto como \texttt{hate} ou \texttt{nao\_hate};
    \item \textbf{classificação multirrótulo dos tipos de ódio}: atribuição de
    zero ou mais categorias temáticas aos registros cujo rótulo gold seja
    \texttt{hate}.
\end{enumerate}

Essa separação define a arquitetura hierárquica candidata: o primeiro modelo
decide se o conteúdo contém discurso de ódio; quando a decisão é positiva, o
segundo identifica os tipos aplicáveis. Sua adoção não será tratada como uma
vantagem presumida. Ela será comparada a uma arquitetura conjunta e terá a
propagação de erros quantificada, conforme a Seção~\ref{sec:arquiteturas}.

\section{Perguntas de pesquisa, hipóteses e contribuição}
\label{sec:perguntas-hipoteses}

Os modelos são empregados como instrumentos para investigar generalização,
desbalanceamento, calibração e propagação de erros, e não apenas como alternativas
de implementação. As perguntas de pesquisa são:

\begin{description}
    \item[RQ1:] um encoder monolíngue pré-treinado em português brasileiro
    apresenta desempenho superior a encoders multilíngues comparáveis na
    detecção de discurso de ódio?
    \item[RQ2:] a ponderação da função de perda melhora a identificação da
    classe minoritária \texttt{hate} e dos tipos menos frequentes?
    \item[RQ3:] a calibração pós-processada reduz o erro das probabilidades
    estimadas sem deteriorar a discriminação?
    \item[RQ4:] o desempenho varia entre plataformas e se mantém quando uma
    plataforma inteira é excluída do treinamento?
    \item[RQ5:] qual é a perda de desempenho causada pela propagação de erros
    entre a detecção binária e a classificação dos tipos?
    \item[RQ6:] o desempenho se mantém em dados coletados depois do período
    representado pela v0.1?
    \item[RQ7:] o pré-treinamento específico em textos de redes sociais produz
    vantagem sobre o pré-treinamento geral da web, especialmente no X/Twitter?
\end{description}

As hipóteses e os contrastes primários serão registrados antes da abertura do
teste. Para RQ1, o contraste principal será BERTimbau Base contra mBERT, pois
ambos possuem porte semelhante; BERTimbau Large e XLM-R serão comparações
secundárias. Para RQ7, o contraste secundário pré-especificado será BERTabaporu
Base contra BERTimbau Base, mantendo as mesmas partições e o mesmo protocolo de ajuste. Seja
$M$ a métrica previamente indicada para cada pergunta. Serão
testadas as seguintes hipóteses direcionais:

\begin{align}
    H_{0,1}&: M_{\mathrm{BERTimbau\ Base}}-M_{\mathrm{mBERT}} \leq 0,
    &H_{1,1}&: M_{\mathrm{BERTimbau\ Base}}-M_{\mathrm{mBERT}} > 0;\\
    H_{0,2}&: M_{\mathrm{ponderada}}-M_{\mathrm{comum}} \leq 0,
    &H_{1,2}&: M_{\mathrm{ponderada}}-M_{\mathrm{comum}} > 0;\\
    H_{0,3}&: E_{\mathrm{calibrado}}-E_{\mathrm{bruto}} \geq 0,
    &H_{1,3}&: E_{\mathrm{calibrado}}-E_{\mathrm{bruto}} < 0;\\
    H_{0,5}&: F1_{\mathrm{oracle}}-F1_{\mathrm{ponta\ a\ ponta}} \leq 0,
    &H_{1,5}&: F1_{\mathrm{oracle}}-F1_{\mathrm{ponta\ a\ ponta}} > 0.
\end{align}

Em RQ3, $E$ designa o Brier score, com o ECE como medida complementar. RQ4 e
RQ6 serão avaliadas como hipóteses de heterogeneidade e degradação de domínio,
respectivamente. RQ6 é prospectiva e somente poderá receber resposta
confirmatória quando houver um teste temporal posterior à v0.1.

Além de significância estatística, será adotada como referência uma diferença
mínima cientificamente relevante de $\delta=0{,}02$ em Macro-F1 ou F1 da classe
\texttt{hate}. Esse valor constitui uma decisão operacional prévia, e não uma
constante universal; análises de sensibilidade usarão também $\delta=0{,}01$ e
$\delta=0{,}03$. A conclusão distinguirá: ausência de evidência, diferença
estatisticamente compatível com zero e diferença superior ao limiar de relevância
prática.

\begin{longtable}{p{0.06\textwidth}p{0.25\textwidth}p{0.30\textwidth}p{0.18\textwidth}}
    \caption{Matriz pré-especificada entre pergunta, hipótese, experimento e métrica.}
    \label{tab:matriz-rq} \\
    \toprule
    \textbf{RQ} & \textbf{Contraste} & \textbf{Experimento} & \textbf{Métrica principal} \\
    \midrule
    \endfirsthead
    \toprule
    \textbf{RQ} & \textbf{Contraste} & \textbf{Experimento} & \textbf{Métrica principal} \\
    \midrule
    \endhead
    RQ1 & monolíngue versus multilíngue & BERTimbau Base versus mBERT nas mesmas partições & Macro-F1 \\
    RQ2 & perda ponderada versus comum & ablação pareada por backbone e semente & F1-\texttt{hate}; AUPRC \\
    RQ3 & probabilidade calibrada versus bruta & \textit{temperature scaling} em conjunto próprio & Brier score; ECE \\
    RQ4 & dentro versus fora do domínio & resultados por plataforma e \textit{leave-one-platform-out} & Macro-F1 \\
    RQ5 & \textit{oracle} versus ponta a ponta & execução das duas formas sobre os mesmos exemplos & $\Delta$Macro-F1 \\
    RQ6 & período histórico versus posterior & teste temporal nas versões posteriores & Macro-F1; AUPRC \\
    RQ7 & pré-treinamento social versus web & BERTabaporu Base versus BERTimbau Base, globalmente e por plataforma & Macro-F1 \\
    \bottomrule
\end{longtable}

A contribuição científica esperada é, portanto, a caracterização controlada de
como idioma de pré-treinamento, desenho amostral, desbalanceamento, calibração,
plataforma e organização hierárquica afetam classificadores de discurso de ódio
em português brasileiro. A construção do dataset e a implementação dos modelos
são meios para essa investigação. Resultados confirmatórios serão separados de
análises exploratórias, especialmente para categorias raras e plataformas com
baixa cobertura.

\section{Definição e congelamento do Dataset v0.1}
\label{sec:dataset-v01}

\subsection{Retrato da base em 10 de setembro de 2026}

Na data de referência deste protocolo, o universo continha 212.476 análises com
estado de processamento bem-sucedido. A base gold reunia 25.090 textos canônicos
com decisão humana resolvida, correspondentes a 11,81\% do universo. Desse total,
3.995 registros haviam sido classificados como \texttt{hate} e 21.095 como
\texttt{nao\_hate}. A distribuição evidencia desbalanceamento entre as classes:
aproximadamente 15,92\% dos registros gold eram positivos e 84,08\% eram
negativos.

\begin{table}[htbp]
    \centering
    \caption{Retrato da base gold em 10 de setembro de 2026.}
    \label{tab:retrato-base-gold}
    \begin{tabular}{lrr}
        \toprule
        \textbf{Indicador} & \textbf{Quantidade} & \textbf{Percentual} \\
        \midrule
        Universo elegível & 212.476 & 100,00\% \\
        Decisões gold resolvidas & 25.090 & 11,81\% \\
        Rótulo \texttt{hate} & 3.995 & 15,92\% do gold \\
        Rótulo \texttt{nao\_hate} & 21.095 & 84,08\% do gold \\
        \bottomrule
    \end{tabular}
    \par\smallskip
    {\footnotesize Fonte: elaboração própria com base nas tabelas de avaliação
    do projeto.\par}
\end{table}

As decisões gold estavam distribuídas entre sete plataformas, conforme a
Tabela~\ref{tab:distribuicao-plataforma}. Essa distribuição deverá ser recalculada
no momento do congelamento, pois o acréscimo de avaliações pode alterar tanto as
contagens quanto as proporções.

\begin{table}[htbp]
    \centering
    \caption{Distribuição preliminar das decisões gold por plataforma.}
    \label{tab:distribuicao-plataforma}
    \begin{tabular}{lr}
        \toprule
        \textbf{Plataforma} & \textbf{Quantidade} \\
        \midrule
        YouTube & 7.735 \\
        Facebook & 4.891 \\
        Reddit & 3.101 \\
        WhatsApp & 3.037 \\
        Twitter & 2.895 \\
        Telegram & 2.660 \\
        TikTok & 771 \\
        \midrule
        \textbf{Total} & \textbf{25.090} \\
        \bottomrule
    \end{tabular}
    \par\smallskip
    {\footnotesize Fonte: elaboração própria com base nas tabelas de avaliação
    do projeto.\par}
\end{table}

Entre os textos positivos, uma mesma publicação pode receber mais de um tipo de
ódio. Por isso, a soma das frequências da Tabela~\ref{tab:tipos-odio-preliminar}
é superior ao número total de textos classificados como \texttt{hate}.

\begin{table}[htbp]
    \centering
    \caption{Frequência preliminar dos tipos de ódio na base gold.}
    \label{tab:tipos-odio-preliminar}
    \begin{tabular}{llr}
        \toprule
        \textbf{Código} & \textbf{Categoria} & \textbf{Ocorrências} \\
        \midrule
        \texttt{political} & Ódio político & 1.841 \\
        \texttt{sexism} & Sexismo & 1.020 \\
        \texttt{lgbtphobia} & LGBTfobia & 725 \\
        \texttt{aporophobia} & Classismo & 316 \\
        \texttt{religious\_intolerance} & Intolerância religiosa & 255 \\
        \texttt{racism} & Racismo & 241 \\
        \texttt{xenophobia} & Xenofobia & 141 \\
        \texttt{body\_shame} & Gordofobia & 123 \\
        \texttt{ageism} & Ageísmo & 52 \\
        \texttt{capacitism} & Capacitismo & 39 \\
        \texttt{other} & Categoria histórica inativa & 1 \\
        \bottomrule
    \end{tabular}
    \par\smallskip
    {\footnotesize Fonte: elaboração própria. As frequências são multirrótulo e,
    portanto, não devem ser somadas para estimar o número de textos positivos.\par}
\end{table}

As categorias menos frequentes ainda não oferecem quantidade suficiente para
uma avaliação confirmatória estável. Esse problema não será resolvido por meio
da eliminação automática de categorias ou pela criação de exemplos artificiais
antes da separação dos dados. Em vez disso, todos os rótulos ativos serão
preservados no treinamento e os resultados das classes raras serão explicitamente
marcados como exploratórios e acompanhados de intervalos de confiança.

\subsection{Justificativa do tamanho da v0.1}

O percentual de 10\% constitui um marco operacional para iniciar o ciclo de
treinamento, e não uma justificativa estatística isolada. O retrato efetivamente
alcançado reúne 25.090 registros e 3.995 positivos, com prevalência observada de
15,92\%. Sob amostragem aleatória simples, a margem de erro aproximada de 95\%
para essa prevalência seria de 0,45 ponto percentual. Esse cálculo não será
interpretado como precisão efetiva porque a base foi estratificada, possui pesos
desiguais e contém grupos correlacionados.

Antes do congelamento será produzido um relatório de suficiência com: contagem
esperada e observada de positivos; contagem por categoria, plataforma e tipo de
conteúdo; tamanho efetivo da amostra; largura dos intervalos de confiança
bootstrap; e quantidade de positivos em cada partição. Para pesos $w_i$, o
tamanho efetivo descritivo será calculado por

\begin{equation}
    n_{\mathrm{efetivo}} =
    \frac{\left(\sum_i w_i\right)^2}{\sum_i w_i^2}.
\end{equation}

Uma análise de poder por simulação, baseada apenas na parcela de desenvolvimento
disponível, estimará a probabilidade de detectar uma diferença pareada de
$\delta=0{,}02$ nas comparações primárias, com nível de 5\% e poder-alvo de 80\%.
A simulação preservará prevalência, estratos e grupos e apresentará também a
menor diferença detectável sob o tamanho alcançado. Seus parâmetros e código
serão arquivados; nenhuma observação do teste será usada para ajustar a análise.

O estudo binário será considerado apto à análise confirmatória apenas se o teste
contiver quantidade suficiente de positivos para produzir estimativas estáveis e
se os intervalos alcançarem precisão compatível com as comparações primárias.
Para os tipos de ódio, a classificação confirmatória ou exploratória obedecerá
aos critérios da Seção~\ref{sec:tarefa-multirrotulo}. Assim, ultrapassar o limiar
de 21.248 registros autoriza o início do benchmark, mas não transforma automaticamente
todas as análises por subgrupo em confirmatórias.

\subsection{Regra de congelamento}

Dez por cento de 212.476 registros correspondem a 21.247,6 observações. Assim,
o limiar operacional da v0.1 será de 21.248 decisões gold válidas. Como as
avaliações são organizadas em lotes, o congelamento ocorrerá ao término do
primeiro lote completo que fizer a base atingir ou ultrapassar esse limiar. Não
será realizado corte dentro de um lote já iniciado, evitando uma seleção
oportunista de casos.

O congelamento deverá produzir um instantâneo imutável, acompanhado de:

\begin{itemize}
    \item data e hora da extração;
    \item consulta SQL ou código de exportação versionado;
    \item número total de registros e distribuição dos rótulos;
    \item identificador do último lote incorporado;
    \item hash SHA-256 do arquivo exportado;
    \item versão da taxonomia e do protocolo de avaliação;
    \item histórico de exclusões e correções aplicadas.
\end{itemize}

Em 10 de setembro de 2026 foi produzida uma exportação candidata com 25.090
decisões válidas. O volume já supera o limiar mínimo da v0.1. Entretanto, o lote
24 ainda constava administrativamente como \texttt{active}, embora todos os seus
953 itens possuíssem decisão resolvida. Por isso, o snapshot pode iniciar a
preparação experimental, mas seu congelamento definitivo deverá registrar o
encerramento administrativo do lote e confirmar novamente as contagens e hashes.

Uma vez congelada, a v0.1 não deverá receber alterações silenciosas. Correções
posteriores de rótulos deverão gerar uma nova versão do conjunto de dados, com
registro das diferenças em relação à versão anterior.

O manifesto de congelamento deverá identificar separadamente as versões do
dataset, da taxonomia e do manual de anotação, por exemplo:

\begin{quote}
\texttt{dataset = radarodio-v0.1}\par
\texttt{taxonomy = taxonomy-v1.2}\par
\texttt{annotation-guidelines = annotation-guidelines-v1.2}
\end{quote}

Mudanças posteriores no manual ou na taxonomia não modificarão retroativamente
o significado de decisões produzidas sob uma versão anterior.

\subsection{Qualidade da anotação e formação do gold standard}
\label{sec:qualidade-anotacao}

Não será realizada nova dupla anotação sistemática da v0.1, pois esse procedimento
é operacionalmente inviável no cronograma do projeto. O controle de confiabilidade
utilizará os arquivos históricos de dupla anotação do Reddit como estudo auxiliar,
mantido separado do teste principal. A auditoria disponível identificou 1.644
pares binários válidos, dos quais 1.512 apresentaram consenso e 132 divergência,
correspondendo a 91,97\% de concordância bruta e $\kappa$ de Cohen de aproximadamente
0,39. Os 1.512 casos consensuais somente poderão ser incorporados depois de
vinculação ao corpus, deduplicação e harmonização da taxonomia; divergências sem
adjudicação permanecerão fora do gold.

Antes da anotação definitiva, o projeto deverá arquivar uma ficha metodológica
com o número e o perfil funcional dos avaliadores, treinamento recebido, versão
do manual, exemplos de treinamento, exercício de qualificação, critérios de
aprovação e procedimento para dúvidas. Informações pessoais dos avaliadores não
entrarão no arquivo destinado ao treinamento dos modelos.

Na tarefa binária, a concordância será descrita pelo percentual bruto, pelo
kappa de Cohen e pelo alfa de Krippendorff. Na tarefa multirrótulo serão
calculados o Jaccard e o F1 entre conjuntos, além da concordância e do alfa de
Krippendorff por categoria \cite{cohen1960kappa,krippendorff2018content}. Todos
os coeficientes serão acompanhados do número de
itens, distribuição marginal e intervalo de confiança. A concordância não será
usada isoladamente como prova de validade; desacordos serão analisados por
categoria, plataforma e fenômeno linguístico.

Haverá divergência quando os avaliadores históricos discordarem sobre o rótulo
binário ou sobre ao menos um tipo ativo. Sem nova adjudicação rastreável, esses
casos não serão transformados em rótulos de treinamento. Se uma adjudicação
pontual vier a ocorrer, deverá registrar decisão, justificativa, versão das
diretrizes, data e categorias em disputa, preservando as duas respostas originais.

Na data de referência, todas as 25.090 decisões resolvidas no sistema tinham
estado \texttt{single}. A v0.1 poderá ser usada para desenvolvimento e avaliação
interna, mas a ausência de dupla anotação contemporânea e multiplataforma limita
a força das conclusões sobre confiabilidade do gold. O estudo histórico do Reddit
será reportado como evidência auxiliar, sem generalização automática para as
demais plataformas.

\subsection{Critérios de elegibilidade}

Serão incluídos somente textos canônicos cujo consenso possua um dos estados
\texttt{single}, \texttt{agreed} ou \texttt{adjudicated}. Serão excluídos:

\begin{itemize}
    \item avaliações puladas ou ainda atribuídas;
    \item conflitos que não tenham sido adjudicados;
    \item registros sem texto ou com falha de codificação não recuperável;
    \item registros sem rótulo final consistente;
    \item cópias não canônicas de textos duplicados;
    \item observações que violem as regras vigentes da taxonomia.
\end{itemize}

Na data de referência, todas as 25.090 decisões resolvidas tinham estado
\texttt{single}. Consequentemente, não é possível estimar concordância entre
avaliadores para essa parcela da base. Essa condição será registrada como uma
limitação da v0.1. Os casos históricos de revisão dupla serão mantidos em um
conjunto auxiliar com proveniência própria e não alterarão silenciosamente o
estado das decisões produzidas pelo sistema.

O código histórico \texttt{other}, atualmente inativo, não integrará o vetor de
saída da tarefa multirrótulo. O registro que ainda contém esse código poderá ser
mantido na tarefa binária, mas deverá ser excluído da tarefa de tipos até que
seja revisto e recodificado em uma categoria ativa.

Antes da exportação serão aplicadas as seguintes invariantes:

\begin{align}
    \texttt{nao\_hate} &\Rightarrow
    \texttt{gold\_hate\_types}=\varnothing,\\
    \texttt{hate} &\Rightarrow
    \left|\texttt{gold\_hate\_types}\right|\geq 1,
\end{align}

em que todos os tipos devem pertencer à versão ativa da taxonomia. A condição
\texttt{hate} sem tipo não será silenciosamente convertida em negativa: o caso
será devolvido para adjudicação ou ficará fora da tarefa multirrótulo até ser
resolvido. Também serão verificadas unicidade de \texttt{analysis\_id}, presença
dos campos obrigatórios, reprodução do \texttt{text\_hash}, ausência de tipos
inativos e inexistência de rótulos conflitantes dentro de um grupo de duplicatas.
O relatório dessas validações será armazenado com o snapshot.

\subsection{Contrato de exportação}

A extração para treinamento deverá seguir o contrato apresentado na
Tabela~\ref{tab:contrato-exportacao}. Campos administrativos e identidades dos
avaliadores não serão incorporados ao conjunto utilizado pelos modelos.

\begin{longtable}{p{0.22\textwidth}p{0.16\textwidth}p{0.48\textwidth}}
    \caption{Contrato mínimo de exportação do Dataset v0.1.}
    \label{tab:contrato-exportacao} \\
    \toprule
    \textbf{Campo} & \textbf{Tipo} & \textbf{Finalidade} \\
    \midrule
    \endfirsthead
    \toprule
    \textbf{Campo} & \textbf{Tipo} & \textbf{Finalidade} \\
    \midrule
    \endhead
    \texttt{analysis\_id} & inteiro & Identificador interno e rastreabilidade da análise. \\
    \texttt{text\_hash} & texto & Identificação de duplicatas exatas e definição dos grupos de partição. \\
    \texttt{text} & texto & Conteúdo anonimizado que será fornecido ao tokenizer. \\
    \texttt{platform} & categoria & Auditoria e avaliação estratificada por plataforma. \\
    \texttt{content\_kind} & categoria & Identificação de post ou comentário. \\
    \texttt{gold\_label} & categoria & Alvo binário: \texttt{hate} ou \texttt{nao\_hate}. \\
    \texttt{gold\_hate\_types} & lista & Conjunto dos tipos de ódio aplicáveis ao texto. \\
    \texttt{decision\_status} & categoria & Indicação de decisão simples, concordante ou adjudicada. \\
    \texttt{batch\_id} & inteiro & Identificação do lote de amostragem. \\
    \texttt{stratum} & texto & Estrato usado na seleção e no cálculo de pesos amostrais. \\
    \texttt{inclusion\_probability} & numérico & Probabilidade de inclusão, quando conhecida e auditável. \\
    \texttt{duplicate\_group\_id} & texto & Grupo transitivo de duplicatas e quase duplicatas. \\
    \texttt{context\_group\_id} & texto/nulo & Conversa, tópico ou publicação de origem compartilhada. \\
    \texttt{content\_created\_at} & data/hora/nulo & Data do conteúdo para avaliação temporal. \\
    \texttt{reviewed\_at} & data/hora & Controle da versão e avaliações temporais futuras. \\
    \texttt{taxonomy\_version} & texto & Versão da taxonomia aplicada ao rótulo. \\
    \texttt{guidelines\_version} & texto & Versão do manual de anotação. \\
    \texttt{context\_dependency} & categoria/nulo & Indício de dependência de imagem, vídeo, conversa ou evento externo. \\
    \bottomrule
\end{longtable}

Os campos \texttt{gemma\_label}, \texttt{gemma\_hate\_probability} e
\texttt{gemma\_hate\_types} poderão ser exportados em uma tabela auxiliar
restrita à auditoria, ao cálculo dos estratos e à reprodução do baseline. Eles
não poderão ser concatenados ao texto nem empregados como atributos de entrada
dos classificadores BERT. Essa separação impede que os novos modelos apenas
reproduzam os sinais do sistema que orientou a amostragem.

\section{Preparação dos textos}
\label{sec:preparacao-textos}

O pré-processamento será deliberadamente conservador. Modelos BERT utilizam
tokenização em subpalavras e foram concebidos para aprender representações a
partir de sequências próximas da linguagem observada. Remoções agressivas podem
eliminar sinais relevantes para discurso de ódio, especialmente pontuação,
emojis, hashtags, caixa, erros ortográficos e alterações intencionais de
grafia.

Cada texto passará, nesta ordem, pelas seguintes operações:

\begin{enumerate}
    \item conversão para Unicode e normalização na forma NFC;
    \item substituição de sequências inválidas de codificação quando a correção
    for inequívoca;
    \item padronização de quebras de linha e espaços consecutivos;
    \item substituição de URLs pelo marcador \texttt{<URL>};
    \item substituição de nomes de usuário e menções pelo marcador
    \texttt{<USER>};
    \item substituição de endereços de e-mail por \texttt{<EMAIL>};
    \item substituição de números de telefone por \texttt{<PHONE>};
    \item remoção de espaços nas extremidades, sem alterar o restante do texto.
\end{enumerate}

Os marcadores serão adicionados como tokens especiais indivisíveis ao tokenizer
de cada checkpoint, com redimensionamento correspondente da matriz de embeddings.
As novas linhas serão inicializadas pelo procedimento padrão definido na
configuração do checkpoint, treinadas durante o fine-tuning e controladas pelas
mesmas sementes das demais inicializações. A lista de tokens, os IDs resultantes
e o hash do tokenizer alterado serão registrados. Uma ablação no conjunto de
desenvolvimento comparará o mascaramento com a manutenção das formas originais,
após anonimização irreversível dos identificadores, para medir eventual perda de
sinais úteis.

Não serão aplicados \textit{stemming}, lematização, remoção de \textit{stopwords},
tradução, conversão geral para minúsculas ou correção ortográfica. Os modelos
sensíveis à caixa receberão a capitalização original. Expressões regulares,
dicionários e versões de qualquer reconhecedor de informação pessoal serão
versionados. Nomes de grupos sociais necessários à interpretação do alvo não
serão removidos como se fossem identificadores pessoais. O texto bruto ficará
em ambiente de acesso restrito; somente a versão anonimizada poderá integrar os
artefatos de treinamento.

O BERTabaporu é um modelo \textit{uncased}; portanto, sua normalização de caixa
será realizada pelo tokenizer do próprio checkpoint. O arquivo de entrada comum
continuará preservando a capitalização original, evitando que a inclusão desse
modelo modifique os dados recebidos pelos encoders sensíveis à caixa.

\subsection{Duplicatas e prevenção de vazamento}

Primeiramente, o mapa de duplicatas do sistema será usado para manter uma única
linha por \texttt{text\_hash} canônico. Em seguida, será criada uma representação
normalizada em 5-gramas de caracteres. Pares com similaridade de Jaccard maior
ou igual a 0,90 serão tratados como quase duplicatas para fins de partição. A
detecção de quase duplicatas não implicará exclusão automática: os registros
serão reunidos sob um mesmo identificador de grupo, garantindo que permaneçam
na mesma partição e no mesmo fold.

O limiar 0,90 será validado antes da divisão definitiva. Serão sorteados pares
dos intervalos 0,80--0,89, 0,90--0,94 e 0,95--1,00 para julgamento humano cego,
registrando precisão e revocação aproximadas do detector de quase duplicatas.
Se a inspeção indicar separação inadequada, qualquer alteração do limiar será
feita e documentada antes da abertura do teste.

O agrupamento será transitivo. Os textos serão vértices de um grafo; cada par
que satisfizer o limiar formará uma aresta; e cada componente conexo receberá
um único \texttt{duplicate\_group\_id}. Assim, se $A$ for semelhante a $B$ e
$B$ a $C$, os três permanecerão juntos mesmo quando a similaridade direta entre
$A$ e $C$ for inferior a 0,90.

Quando houver identificador confiável de conversa, tópico ou encadeamento, todos
os textos do mesmo contexto também serão mantidos no mesmo grupo. Considera-se
contexto compartilhado a mesma postagem original e seus comentários, a mesma
thread, o mesmo encadeamento de respostas ou uma conversa identificada pelo
mesmo ID de origem. Um \texttt{context\_group\_id} prevalecerá sobre a decisão
individual de partição. Essa regra reduz o risco de um comentário no treinamento
revelar o conteúdo de outro comentário relacionado presente no teste.

\subsection{Comprimento das sequências}

Antes de definir o truncamento, será calculada, separadamente para cada tokenizer,
a distribuição do comprimento dos textos. Serão reportados média, mediana,
percentis 90, 95 e 99, máximo e as proporções $P(L>128)$, $P(L>256)$ e
$P(L>512)$. O limite técnico superior será de 512 tokens, incluindo os tokens
especiais do encoder.

Três políticas serão comparadas no desenvolvimento: (A) primeiros 510 tokens de
conteúdo; (B) 255 tokens iniciais e 255 finais; e (C) janelas de até 510 tokens
com sobreposição, agregadas no nível do documento. Para a política C, tamanho da
janela, sobreposição e regra de agregação serão fixados antes do teste. Se menos
de 5\% dos textos excederem 512 tokens e não houver diferença relevante no
desenvolvimento, a política B será adotada por parcimônia. Caso contrário, a
política será tratada como hiperparâmetro de ablação. A quantidade de textos
truncados, o conteúdo médio descartado e o custo adicional do \textit{chunking}
serão registrados.

\section{Definição das tarefas de aprendizagem}
\label{sec:tarefas-aprendizagem}

\subsection{Tarefa A: detecção binária}

Para um texto $x_i$, o alvo binário será definido por

\begin{equation}
    y_i =
    \begin{cases}
        1, & \text{se o rótulo gold for \texttt{hate}},\\
        0, & \text{se o rótulo gold for \texttt{nao\_hate}}.
    \end{cases}
\end{equation}

O vetor contextual associado ao token de classificação será fornecido a uma
camada linear com duas saídas, seguida por \textit{softmax}. A função de perda
será a entropia cruzada. Serão comparadas duas condições: perda sem ponderação e
perda ponderada por classe. Na condição ponderada, o peso da classe $c$ será
calculado apenas no subconjunto de treinamento de cada fold:

\begin{equation}
    \alpha_c = \frac{N_{\mathrm{treino}}}{2N_c},
\end{equation}

em que $N_c$ é o número de exemplos da classe $c$ no fold de treinamento.

\subsection{Tarefa B: classificação multirrótulo}
\label{sec:tarefa-multirrotulo}

A tarefa multirrótulo utilizará somente os textos com rótulo gold
\texttt{hate}. Para cada texto, será construído um vetor binário
$\mathbf{y}_i \in \{0,1\}^{10}$, na seguinte ordem fixa:

\begin{enumerate}
    \item \texttt{ageism};
    \item \texttt{aporophobia};
    \item \texttt{body\_shame};
    \item \texttt{capacitism};
    \item \texttt{lgbtphobia};
    \item \texttt{political};
    \item \texttt{racism};
    \item \texttt{religious\_intolerance};
    \item \texttt{sexism};
    \item \texttt{xenophobia}.
\end{enumerate}

No banco de origem, o rótulo de sexismo ainda está armazenado sob o código
histórico \texttt{misogyny}. Durante a preparação da tarefa multirrótulo, esse
valor será mapeado deterministicamente para \texttt{sexism}, sem alterar o
arquivo bruto. A tabela de mapeamento e sua versão integrarão o manifesto do
dataset.

O cabeçote terá dez saídas independentes com ativação sigmoide. A perda será a
entropia cruzada binária com logits, implementada por
\texttt{BCEWithLogitsLoss}. Também serão comparadas uma condição sem ponderação
e outra com peso positivo por categoria. Para a categoria $\ell$, o peso será
calculado exclusivamente no fold de treinamento:

\begin{equation}
    p_{\ell} = \frac{N_{\ell}^{-}}{N_{\ell}^{+}},
\end{equation}

em que $N_{\ell}^{+}$ e $N_{\ell}^{-}$ são, respectivamente, as quantidades de
exemplos positivos e negativos para o rótulo $\ell$.

Uma categoria será considerada elegível para inferência confirmatória somente
se possuir pelo menos 100 ocorrências na v0.1 e pelo menos 20 ocorrências no
conjunto de teste. Esses valores são salvaguardas operacionais mínimas para
evitar folds praticamente vazios; não garantem, sozinhos, precisão inferencial.
Por isso, serão apresentados intervalos de confiança e uma análise de
sensibilidade exigindo 20, 30 e 50 positivos no teste. Categorias cujo intervalo
de 95\% para F1 permaneça excessivamente amplo serão descritas como exploratórias,
mesmo que ultrapassem o limite mínimo.

\section{Arquiteturas e propagação de erros}
\label{sec:arquiteturas}

O experimento principal comparará duas organizações sob o mesmo encoder,
partições, pré-processamento e orçamento de ajuste:

\begin{enumerate}
    \item \textbf{hierárquica}: um classificador binário é executado primeiro e
    o classificador multirrótulo somente produz tipos quando o primeiro estágio
    prevê \texttt{hate};
    \item \textbf{conjunta não hierárquica}: um único encoder possui onze saídas
    sigmoides, uma para \texttt{hate} e dez para os tipos. Registros
    \texttt{nao\_hate} recebem zero nos dez tipos, e nenhuma regra de bloqueio é
    aplicada antes da medição. Inconsistências, como prever um tipo sem prever
    \texttt{hate}, serão contabilizadas.
\end{enumerate}

Na arquitetura conjunta, a perda será a soma ponderada da perda binária e da
perda dos tipos,

\begin{equation}
    \mathcal{L}_{\mathrm{conjunta}} =
    \mathcal{L}_{\mathrm{binária}} +
    \gamma\mathcal{L}_{\mathrm{tipos}},
\end{equation}

com $\gamma$ definido exclusivamente no desenvolvimento. Para limitar o custo,
essa ablação será executada inicialmente com o backbone vencedor da tarefa
binária; uma extensão aos demais backbones será identificada como análise
secundária.

O classificador de tipos será avaliado tanto com a entrada binária gold
(condição \textit{oracle}) quanto no pipeline completo. A propagação de erros
será quantificada por

\begin{equation}
    \Delta F1_{\mathrm{propagação}} =
    F1_{\mathrm{oracle}}-F1_{\mathrm{ponta\ a\ ponta}},
\end{equation}

globalmente, por categoria e por plataforma. A comparação hierárquica versus
conjunta utilizará predições pareadas no mesmo teste; qualquer vantagem deverá
ser interpretada junto do custo computacional e da taxa de inconsistências.

\section{Particionamento e desenho de validação}
\label{sec:particionamento}

\subsection{Teste fixo e conjunto de desenvolvimento}

Após a formação dos grupos de duplicatas, 20\% da v0.1 serão reservados como
conjunto de teste fixo. Dos 80\% restantes, 12,5\% serão separados previamente
como conjunto de calibração, equivalentes a 10\% da base completa. Os 70\%
restantes formarão o conjunto de seleção e ajuste. A proporção global será,
portanto, 70\% para seleção/treinamento, 10\% para calibração e limiares e 20\%
para teste. Todas as divisões respeitarão os grupos e a estratificação.

O conjunto de calibração terá três usos previamente limitados: estimar o
parâmetro de calibração, definir os limiares da configuração já selecionada e
medir a estabilidade dessas escolhas. Ele não será usado para escolher backbone,
taxa de aprendizado, perda, época ou política de truncamento. O teste será
armazenado separadamente e não poderá participar de:

\begin{itemize}
    \item ajuste de hiperparâmetros;
    \item seleção de modelo ou de época;
    \item definição de pesos de classe;
    \item escolha de limiares de decisão;
    \item calibração de probabilidades ou escolha de seus métodos;
    \item análise exploratória destinada a alterar o pipeline.
\end{itemize}

A divisão será realizada com estratificação iterativa multirrótulo, abordagem
destinada a conservar distribuições de rótulos em problemas nos quais uma mesma
instância pode pertencer a várias classes \cite{sechidis2011stratification}. O
vetor de estratificação combinará o rótulo binário, os dez tipos de ódio, a
plataforma, o tipo de conteúdo e o estrato de amostragem. A unidade indivisível
será o grupo de duplicatas, e não o registro individual.

\subsection{Validação cruzada}

O conjunto de seleção e ajuste será dividido em cinco folds por meio da mesma
estratégia de agrupamento e estratificação. Em cada rodada, quatro folds serão
usados para treinamento e um para validação. Ao final das cinco rodadas, cada
registro desse conjunto terá exatamente uma predição fora da amostra,
denominada \textit{out-of-fold} (OOF).

As predições OOF serão usadas para comparar configurações e estimar sua
variabilidade sem consultar o conjunto de calibração. A métrica OOF global,
calculada após concatenar as cinco predições de validação, será o critério
principal de seleção. A média e o desvio-padrão das cinco métricas calculadas
separadamente também serão reportados, mas não substituirão o critério global.
Essa escolha é especialmente importante porque o
fine-tuning de modelos pré-treinados pode apresentar variação relevante em
função da inicialização e da ordem dos dados, sobretudo em conjuntos menores
\cite{dodge2020finetuning}.

\subsection{Correção do desenho amostral}

Os lotes de avaliação foram balanceados a partir de informações como plataforma,
rótulo previsto pelo Gemma e tipo previsto de ódio. Logo, a distribuição da
base anotada pode diferir da distribuição do universo completo. O manifesto de
amostragem registrará o universo disponível em cada lote, definição e tamanho
dos estratos, quantidade selecionada, algoritmo, semente, critérios de exclusão,
uso das previsões do Gemma e qualquer priorização manual. Para cada estrato $h$,
será calculado o peso de expansão

\begin{equation}
    w_h = \frac{N_h}{n_h},
\end{equation}

em que $N_h$ é o número de registros elegíveis do estrato no universo e $n_h$
é o número de registros selecionados do mesmo estrato. A expressão corresponde
ao inverso da probabilidade de inclusão somente quando a seleção é conhecida e
equiprovável dentro de cada estrato. Para estabilidade numérica, os pesos serão
normalizados para média 1 dentro de cada partição.

As métricas do teste serão apresentadas em duas versões: não ponderada, que
descreve diretamente o conjunto anotado, e ponderada, que busca estimar o
desempenho sob a composição do universo original. Essa segunda interpretação
depende das hipóteses de inclusão conhecível, positividade e seleção ignorável
dentro dos estratos. Seleção manual, revisão dirigida ou exclusão subjetiva
viola essas hipóteses e não é corrigida automaticamente por $w_h$; nessa situação,
a versão ponderada será qualificada como ajuste descritivo, não como estimativa
representativa da população. Os pesos não serão calculados com base no rótulo
gold do teste e jamais serão derivados de seus erros.

Sempre que operacionalmente possível, será sorteada do universo uma amostra
probabilística independente das previsões do Gemma. Ela será anotada pelo mesmo
protocolo e usada como teste externo da prevalência, dos pesos e da generalização.
Se essa amostra não existir na v0.1, a ausência será registrada como limitação e
sua construção será prioridade da versão seguinte.

\section{Modelos comparados}
\label{sec:modelos-comparados}

Cinco encoders serão submetidos ao mesmo desenho experimental:

\begin{description}
    \item[BERTimbau Base:] \texttt{neuralmind/bert-base-portuguese-cased},
    modelo BERT Base monolíngue, sensível à caixa e pré-treinado para português
    brasileiro \cite{souza2020bertimbau,neuralmind2026bertimbau};
    \item[BERTimbau Large:] \texttt{neuralmind/bert-large-portuguese-cased},
    versão de maior capacidade do mesmo projeto
    \cite{souza2020bertimbau,neuralmind2026bertimbau};
    \item[BERTabaporu Base:]
    \texttt{pablocosta/bertabaporu-base-uncased}, modelo BERT Base monolíngue
    sem distinção de caixa, pré-treinado em aproximadamente 238 milhões de
    tweets em português brasileiro, provenientes de mais de 100 mil usuários e
    totalizando cerca de 2,9 bilhões de tokens
    \cite{dacosta2023bertabaporu,dacosta2026bertabaporu};
    \item[mBERT:] \texttt{google-bert/bert-base-multilingual-cased}, modelo BERT
    multilíngue sensível à caixa, pré-treinado em 104 idiomas
    \cite{devlin2019bert,google2026mbert};
    \item[XLM-R Base:] \texttt{FacebookAI/xlm-roberta-base}, encoder
    multilíngue baseado em RoBERTa, incluído como comparador da família de
    Transformers bidirecionais derivados de BERT
    \cite{conneau2020xlmr,meta2026xlmr}.
\end{description}

O BERTabaporu permitirá testar diretamente o efeito do gênero de pré-treinamento:
BERTimbau Base representa o português geral da web, enquanto BERTabaporu Base
representa a linguagem de redes sociais. Como a v0.1 é multiplataforma, eventual
vantagem no X/Twitter não será generalizada para as demais fontes. Serão
reportados o resultado global, o resultado por plataforma e a interação entre
modelo e plataforma. A versão
\texttt{pablocosta/bertabaporu-large-uncased} ficará como extensão secundária,
condicionada ao orçamento computacional, e não substituirá a execução do Base
no benchmark principal.

Os checkpoints e tokenizers serão carregados por meio da biblioteca
Transformers \cite{wolf2020transformers}. A revisão exata de cada repositório
será fixada por hash de commit no início dos experimentos, evitando que uma
atualização remota altere silenciosamente os resultados.

\subsection{Baselines}

Três famílias de baselines serão mantidas:

\begin{enumerate}
    \item \textbf{baseline trivial}: previsão constante de \texttt{nao\_hate}
    na tarefa binária e previsão determinada apenas pela prevalência de cada
    categoria, estimada no treinamento, na tarefa multirrótulo;
    \item \textbf{TF--IDF com regressão logística}: concatenação de n-gramas de
    palavras de ordem 1--2 e n-gramas de caracteres de ordem 3--5, seguida por
    regressão logística. O parâmetro $C$ será selecionado entre 0,1, 1 e 10 no
    mesmo protocolo de cinco folds. Na tarefa multirrótulo será usado
    \texttt{OneVsRestClassifier}, com um classificador logístico independente
    para cada um dos dez tipos;
    \item \textbf{Gemma}: uso das predições originais registradas antes da
    revisão humana, sem novo ajuste. Esse baseline permitirá comparar os BERTs
    com o sistema responsável pelas classificações iniciais.
\end{enumerate}

O baseline TF--IDF será submetido às mesmas partições, agrupamentos e métricas.
Nenhum vocabulário ou transformador TF--IDF será ajustado fora do fold de
treinamento. Como as previsões do Gemma ajudaram a definir estratos e prioridades
de revisão, o desempenho desse baseline não será interpretado como proveniente
de uma amostra independente do próprio sistema. A possível dependência entre a
amostragem e seus erros será explicitada em todas as comparações.

\section{Protocolo de fine-tuning}
\label{sec:fine-tuning}

Todos os encoders receberão cabeçotes de classificação novos, inicializados
aleatoriamente. Os parâmetros do encoder e do cabeçote serão ajustados em
conjunto. A faixa inicial segue a prática de fine-tuning apresentada para BERT,
enquanto decisões adicionais serão verificadas em piloto no próprio domínio
\cite{devlin2019bert}. O protocolo comum será:

\begin{itemize}
    \item otimizador AdamW, com $\beta_1=0{,}9$, $\beta_2=0{,}999$ e
    $\epsilon=10^{-8}$;
    \item tamanho efetivo de lote igual a 32, obtido por acúmulo de gradientes
    quando necessário;
    \item limite superior de oito épocas, entendido como teto operacional e não
    como duração fixa;
    \item avaliação ao final de cada época e parada antecipada após duas épocas
    completas sem redução da perda de validação;
    \item seleção do checkpoint de menor perda de validação, antes da otimização
    de qualquer limiar;
    \item aquecimento linear durante os 10\% primeiros passos;
    \item decaimento linear da taxa de aprendizado após o aquecimento;
    \item norma máxima do gradiente igual a 1,0;
    \item precisão mista FP16 ou BF16 quando suportada pelo hardware;
    \item embaralhamento reprodutível dos exemplos a cada época.
\end{itemize}

A perda de validação foi escolhida para a parada antecipada porque independe de
um limiar de classificação; Macro-F1, F1-\texttt{hate} e AUPRC serão monitorados,
mas não mudarão o checkpoint durante essa etapa. O teto de oito épocas apenas
permite que a regra de paciência opere em tarefas multirrótulo de convergência
mais lenta. O artigo original de BERT usou buscas mais curtas, tipicamente entre
duas e quatro épocas \cite{devlin2019bert}; portanto, será realizado um piloto
com curvas de treinamento. Se as curvas mostrarem convergência sistemática antes
do teto, o custo será reduzido pela parada antecipada. Se permanecerem em melhora
na oitava época, qualquer ampliação exigirá emenda prévia ao protocolo.

Será avaliada a grade cartesiana formada por

\begin{equation}
    \eta \in \{10^{-5}, 2\times10^{-5}, 3\times10^{-5}, 5\times10^{-5}\}
\end{equation}

e

\begin{equation}
    \lambda \in \{0, 0{,}01\},
\end{equation}

em que $\eta$ representa a taxa de aprendizado e $\lambda$ o decaimento de
peso. O valor $10^{-5}$ foi acrescentado às taxas usuais de BERT para contemplar
ajustes mais conservadores, especialmente no modelo Large. A condição de perda
comum ou ponderada será um fator experimental controlado, e não uma decisão
tomada depois de observar o teste.

\subsection{Busca de hiperparâmetros em duas etapas}

Uma grade completa com cinco modelos, duas tarefas, oito combinações, cinco
folds e duas perdas produziria 800 fine-tunings antes das repetições por semente.
Para reduzir custo sem escolher configurações pelo teste, será usada a seguinte
triagem pré-especificada:

\begin{enumerate}
    \item para cada combinação de backbone, tarefa e condição de perda, as oito
    combinações de $\eta$ e $\lambda$ serão executadas por até duas épocas em uma
    divisão interna fixa do conjunto de seleção, com semente 42;
    \item as duas melhores combinações de cada célula serão continuadas até o
    teto de oito épocas, sujeitas à mesma parada antecipada;
    \item a melhor combinação de cada célula seguirá para os cinco folds. Se as
    duas primeiras diferirem menos de 0,005 em Macro-F1, ambas seguirão para
    evitar descarte baseado em ruído da triagem;
    \item a configuração será selecionada pela métrica OOF global, obtida pela
    concatenação das predições dos cinco folds com limiar 0,5. Somente depois
    dessa seleção serão definidos calibração e limiares.
\end{enumerate}

Sem empates, a triagem compreende 160 execuções curtas, 40 continuações e 100
ajustes em folds, em lugar de 800 ajustes completos. Antes da execução será
estimado o orçamento em GPU-horas a partir do piloto, incluindo as repetições
finais e a ablação de arquitetura. Se o orçamento disponível for insuficiente,
o escopo será reduzido por regra registrada antes do teste, priorizando as
perguntas de pesquisa, e não por resultados parciais convenientes.

Na tarefa binária, Macro-F1 será o critério de seleção; F1-\texttt{hate} e AUPRC
serão critérios secundários. Na tarefa multirrótulo, será usado Macro-F1 das
categorias previamente elegíveis. Empates inferiores a 0,005 serão resolvidos,
nesta ordem, por menor variância entre folds, menor perda OOF e menor taxa de
aprendizado. Número de parâmetros e tempo de inferência não serão usados para
desempatar taxas do mesmo backbone, pois permanecem essencialmente constantes.

\subsection{Repetições e definição do modelo final}

Após a seleção, cada condição experimental será treinada no conjunto de seleção
e ajuste com as sementes 13, 42, 97, 123 e 2026. O número de épocas será fixado
na mediana das melhores épocas dos cinco folds, evitando parada antecipada no
teste ou no conjunto de calibração. As probabilidades dos cinco modelos serão
preservadas para medir variabilidade de otimização.

Para a avaliação principal de cada condição, o preditor final será o ensemble
pré-especificado das cinco sementes, formado pela média dos logits antes da
calibração. Nenhuma semente será escolhida pelo desempenho no teste. Caso seja
necessário disponibilizar também um único checkpoint por restrição operacional,
a semente 42 será usada por definição prévia e seus resultados serão reportados
separadamente dos resultados do ensemble.

\section{Limiarização e calibração}
\label{sec:limiarizacao-calibracao}

Na tarefa binária, o ensemble produzirá um logit médio e uma probabilidade
estimada de \texttt{hate}. Depois de fixados backbone, hiperparâmetros, perda,
época e sementes, o limiar $\tau$ será escolhido exclusivamente no conjunto de
calibração. Serão testados valores entre 0,01 e 0,99 em incrementos de 0,01,
selecionando-se o valor que maximizar o Macro-F1. Em empate, será escolhido o
limiar mais próximo de 0,5 e, persistindo o empate, o maior limiar. O valor 0,5
também será reportado como referência.

Na tarefa multirrótulo, cada categoria $\ell$ terá um limiar próprio
$\tau_{\ell}$, escolhido na mesma grade e no mesmo conjunto para maximizar o F1
da categoria. Se houver menos de cinco exemplos positivos de uma categoria na
calibração, seu limiar será mantido em 0,5 e o resultado será tratado como
exploratório. A estabilidade será descrita reestimando os limiares em
reamostragens agrupadas da calibração e, apenas como diagnóstico, nas validações
dos cinco folds; serão reportados mediana, intervalo interquartil e amplitude.

As probabilidades da tarefa binária serão adicionalmente calibradas por
\textit{temperature scaling} ajustado nos logits médios das cinco sementes no
conjunto reservado de calibração. Essa separação evita transferir diretamente
um calibrador estimado sobre logits OOF, produzidos por modelos distintos, para
um modelo final reajustado. Redes neurais modernas podem produzir probabilidades
excessivamente confiantes, e o método oferece uma correção pós-processada simples
\cite{guo2017calibration}.

A temperatura não ponderada estimará calibração para a distribuição da amostra
anotada. Uma temperatura ponderada, obtida pela minimização da log-verossimilhança
negativa com pesos de inclusão, será estimada para a população somente se a
auditoria do desenho amostral validar esses pesos. Caso contrário, não será feita
alegação de que $P(\texttt{hate}\mid\text{amostra})$ corresponde a
$P(\texttt{hate}\mid\text{população})$. O teste receberá apenas transformações e
limiares já congelados. ECE será calculado com 15 intervalos de frequência
aproximadamente igual; também serão apresentados Brier score e diagramas de
confiabilidade antes e depois da calibração.

\section{Métricas de avaliação}
\label{sec:metricas-avaliacao}

\subsection{Tarefa binária}

O Macro-F1 será a métrica primária de comparação geral. F1-\texttt{hate} e AUPRC
serão métricas centrais para as hipóteses sobre a classe positiva, mas não
substituirão retroativamente a métrica definida para cada contraste na
Tabela~\ref{tab:matriz-rq}. Para duas classes $C$, Macro-F1 será calculado
por

\begin{equation}
    \operatorname{MacroF1} = \frac{1}{|C|}\sum_{c\in C}
    \frac{2P_cR_c}{P_c+R_c},
\end{equation}

em que $P_c$ e $R_c$ representam precisão e revocação da classe $c$. O uso de
Macro-F1 impede que o desempenho da classe majoritária domine sozinho a
comparação.

Também serão reportados:

\begin{itemize}
    \item precisão, revocação e F1 da classe \texttt{hate};
    \item área sob a curva Precisão--Revocação (AUPRC);
    \item área sob a curva ROC (AUROC);
    \item coeficiente de correlação de Matthews (MCC);
    \item acurácia balanceada;
    \item matriz de confusão;
    \item Brier score;
    \item erro esperado de calibração (ECE);
    \item curvas de confiabilidade antes e depois da calibração.
\end{itemize}

A AUPRC receberá destaque porque a classe positiva é minoritária. Em cenários
desbalanceados, curvas Precisão--Revocação são mais informativas sobre a
qualidade das predições positivas do que a curva ROC isoladamente
\cite{saito2015precisionrecall}.

\subsection{Tarefa multirrótulo}

A métrica primária será o Macro-F1 calculado sobre as categorias elegíveis para
análise confirmatória. Serão apresentados, adicionalmente:

\begin{itemize}
    \item Macro-F1 incluindo todos os rótulos ativos;
    \item Micro-F1;
    \item precisão, revocação, F1 e PR-AUC por categoria;
    \item \textit{Hamming loss};
    \item \textit{subset accuracy}, que exige correspondência exata do conjunto
    de rótulos;
    \item média da quantidade de rótulos previstos por texto;
    \item matriz de coocorrência entre tipos verdadeiros e previstos.
\end{itemize}

Serão realizadas duas avaliações complementares. A avaliação \textit{oracle}
aplicará o classificador de tipos somente aos registros cujo rótulo binário gold
seja \texttt{hate}, medindo a capacidade específica de reconhecer categorias.
A avaliação ponta a ponta executará primeiro o classificador binário e somente
atribuirá tipos aos textos previstos como positivos. A diferença entre os dois
resultados quantificará a propagação de erros da primeira etapa.

\subsection{Incerteza e comparação entre modelos}

Serão separadas duas fontes de incerteza. A variabilidade de otimização será
descrita pela média, desvio-padrão, mínimo e máximo entre as cinco sementes. A
incerteza amostral será estimada sobre a predição final pré-especificada por
2.000 reamostragens bootstrap do conjunto de teste \cite{efron1993bootstrap}. O
bootstrap será estratificado
pelos estratos de amostragem e terá como unidade o maior agrupamento entre
\texttt{duplicate\_group\_id} e \texttt{context\_group\_id}, preservando a
dependência entre textos relacionados. Na versão ponderada, os pesos serão
mantidos em cada reamostragem.

Diferenças entre modelos serão examinadas por bootstrap pareado, calculando
diretamente, em cada reamostragem, $\Delta M=M_A-M_B$ sobre os mesmos exemplos.
Não será inferida diferença pela simples sobreposição de intervalos individuais.
Os quatro contrastes primários da Tabela~\ref{tab:matriz-rq} serão apresentados
sem expansão posterior do conjunto de hipóteses. Para famílias de comparações
secundárias entre modelos, plataformas ou categorias, os valores de $p$ serão
ajustados pelo procedimento de Holm \cite{holm1979procedure}. Magnitude,
intervalo de confiança e relação
com $\delta$ serão priorizados sobre uma decisão binária de significância.

Além dos resultados globais, as métricas serão desagregadas por plataforma,
tipo de conteúdo e categoria de ódio. Subgrupos com menos de 30 exemplos
positivos não receberão interpretação comparativa; seus valores serão exibidos
apenas como evidência exploratória, acompanhados do tamanho da amostra e do
intervalo de confiança.

Para avaliar generalização entre plataformas, será realizado um experimento
\textit{leave-one-platform-out}: em cada rodada, uma plataforma será mantida
inteiramente fora do ajuste e usada como domínio de teste. A rodada somente será
confirmatória quando a plataforma excluída possuir ao menos 100 registros e 20
positivos; nos demais casos será exploratória. Os resultados serão reportados
para YouTube, WhatsApp, Telegram, Twitter/X, Facebook, Reddit e TikTok, sem
combinar artificialmente plataformas pequenas para obter significância.

Também serão registrados parâmetros, GPU-horas, memória máxima, duração do
treinamento e tempo de inferência por mil textos. A interpretação final comparará
ganho de desempenho e custo, evitando concluir que um modelo maior é preferível
quando a diferença estiver abaixo de $\delta$ e o custo for substancialmente
superior.

\section{Testes funcionais e análise de erros}
\label{sec:testes-funcionais}

Métricas agregadas em um conjunto de teste podem ocultar falhas sistemáticas de
modelos de detecção de discurso de ódio. O HateCheck demonstra a utilidade de
testes funcionais construídos para fenômenos linguísticos específicos
\cite{rottger2021hatecheck}. Com base nessa abordagem, será produzido um pequeno
conjunto diagnóstico em português brasileiro, separado de treinamento,
validação e teste estatístico.

O conjunto funcional terá, como meta inicial, pelo menos 20 exemplos por
fenômeno, equilibrando casos positivos e negativos quando isso for semanticamente
aplicável. Os exemplos poderão ser adaptados de taxonomias funcionais publicadas,
construídos de forma controlada ou extraídos do domínio, desde que origem e
licença sejam registradas. O conjunto deverá contemplar, no mínimo:

\begin{itemize}
    \item negação de uma afirmação preconceituosa;
    \item citação de discurso de ódio para denúncia ou análise;
    \item contra-discurso e defesa do grupo atacado;
    \item ironia e sarcasmo;
    \item menção neutra a grupos sociais;
    \item insulto individual sem ataque generalizante a um grupo;
    \item grafias alteradas, separação de caracteres e uso de símbolos;
    \item ameaça, incitação, desumanização e criminalização coletiva;
    \item ataques explícitos e implícitos;
    \item variações de gênero, número e referência pronominal.
\end{itemize}

Os casos serão classificados por dois avaliadores independentes e divergências
serão adjudicadas pelo protocolo da Seção~\ref{sec:qualidade-anotacao}. Antes do
uso, cada caso será comparado com treino, calibração e teste por hash, Jaccard e
busca aproximada; casos equivalentes serão removidos ou reescritos e novamente
validados. O conjunto funcional não será usado para ajustar hiperparâmetros. O
desempenho será relatado por fenômeno, sem misturá-lo à estimativa principal do
teste da v0.1.

Após a avaliação quantitativa, será conduzida análise qualitativa de erros. Para
cada modelo serão examinados, no mínimo, 50 falsos positivos e 50 falsos
negativos da tarefa binária, sorteados com semente 2026 e estratificação por
plataforma e faixa de confiança. Na tarefa multirrótulo serão sorteados erros das
categorias mais frequentes e incluídos os erros disponíveis das categorias raras,
respeitando as restrições de privacidade.

Antes de examinar as predições será congelada a seguinte taxonomia inicial:
\texttt{E01} contexto ausente; \texttt{E02} ironia; \texttt{E03} sarcasmo;
\texttt{E04} negação; \texttt{E05} alvo implícito; \texttt{E06} ambiguidade
lexical; \texttt{E07} linguagem codificada; \texttt{E08} truncamento;
\texttt{E09} possível erro de anotação; \texttt{E10} conflito entre categorias;
\texttt{E11} conhecimento externo necessário; e \texttt{E12} dependência
multimodal. Novos códigos poderão emergir, mas serão identificados como
pós-especificados. Uma parcela mínima de 20\% dos erros será codificada por dois
avaliadores, com concordância reportada.

Cada item indicará se depende de imagem, vídeo, mensagem anterior, thread,
evento externo ou conhecimento sociopolítico. Será reportada a proporção de
erros potencialmente insolúveis pelo texto isolado, distinguindo limitação do
modelo de limitação da modalidade de entrada.

\section{Reprodutibilidade, documentação e uso responsável}
\label{sec:reprodutibilidade}

Cada execução deverá gerar um manifesto contendo:

\begin{itemize}
    \item versão e hash SHA-256 do dataset;
    \item hash do código e da consulta de extração;
    \item commit do repositório, estado de alterações locais e comando executado;
    \item arquivos imutáveis de \texttt{train\_ids}, \texttt{validation\_ids},
    \texttt{calibration\_ids}, \texttt{test\_ids},
    \texttt{duplicate\_groups} e \texttt{context\_groups};
    \item checkpoint e revisão do tokenizer;
    \item versões de Python, PyTorch, Transformers e demais dependências;
    \item modelo, tarefa, hiperparâmetros e semente;
    \item número de épocas, melhor época nos folds e curvas de perda;
    \item hardware, precisão numérica, memória máxima, GPU-horas e tempo de
    treinamento e inferência;
    \item limiares, temperatura de calibração e pesos de classe;
    \item métricas globais, ponderadas e por subgrupo;
    \item caminho e hash dos checkpoints resultantes.
\end{itemize}

O ambiente de software será congelado por arquivo de dependências com versões
exatas e, preferencialmente, por imagem de contêiner. Não se dependerá apenas da
semente para reconstruir partições. As configurações de determinismo do PyTorch,
CUDA e cuDNN serão registradas, incluindo flags ativadas e operações conhecidas
como não determinísticas. Quando determinismo estrito reduzir de modo relevante
o desempenho computacional ou não estiver disponível, a exceção será documentada.

A documentação do dataset descreverá origem, composição, processo de amostragem,
população representada, critérios de anotação, limitações e usos não recomendados.
Esse procedimento é coerente com a proposta de \textit{data statements}, que
busca tornar explícitas as condições sob as quais resultados de PLN podem ou não
ser generalizados \cite{bender2018datastatements}. Cada modelo terá também uma
ficha com finalidade, desempenho, grupos avaliados, limitações, riscos e uso
pretendido, seguindo a proposta de \textit{model cards}
\cite{mitchell2019modelcards}.

Os classificadores não deverão ser descritos como substitutos da avaliação
humana. Seu uso previsto é priorizar, apoiar ou acelerar a triagem de conteúdos.
Decisões de alto impacto, estudos de prevalência e análises sobre grupos pouco
representados deverão considerar a incerteza, o desenho amostral e a revisão de
especialistas.

Exemplos divulgados em artigos, dissertação ou fichas de modelo passarão por
revisão específica de risco de reidentificação. Conforme o risco, serão usados
mascaramento adicional, paráfrase fiel, exemplo sintético controlado ou omissão.
Identificadores de origem e textos brutos não acompanharão os artefatos públicos.

\section{Continuidade do processo}
\label{sec:continuidade-processo}

O teste da v0.1 permanecerá congelado nas versões posteriores, permitindo medir
ganhos ou regressões no mesmo conjunto histórico. Contudo, após ter sido
consultado em sucessivos ciclos, ele será interpretado como benchmark histórico
ou teste de regressão, e não indefinidamente como única avaliação confirmatória.
Novos registros avaliados serão incorporados ao desenvolvimento das versões
v0.2, v0.3 e seguintes. Cada versão reservará um teste temporal formado por
conteúdos estritamente posteriores ao período de treinamento; o teste temporal
mais recente assumirá progressivamente o papel confirmatório.

Para cada atualização serão repetidos os controles de elegibilidade,
duplicidade, estratificação, treinamento e avaliação. A data de corte temporal
será definida antes do treinamento e os grupos de contexto não poderão atravessá-la.
Um modelo novo somente substituirá o anterior se não apresentar regressão
relevante no teste histórico e demonstrar desempenho aceitável no teste temporal
mais recente.

O pré-treinamento adaptativo em dados não rotulados do domínio poderá ser
investigado em uma etapa futura. Estudos anteriores mostram que a continuação do
pré-treinamento em textos do domínio ou da tarefa pode produzir ganhos mesmo em
cenários com poucos rótulos \cite{gururangan2020dontstop}. Contudo, esse
procedimento não integrará o benchmark principal da v0.1, que comparará os
checkpoints públicos originais sob o mesmo protocolo supervisionado.

\section{Estudos de ablação}
\label{sec:ablacoes}

As decisões de engenharia que possam alterar as conclusões serão transformadas
em contrastes controlados. As ablações pré-especificadas são:

\begin{itemize}
    \item perda comum versus perda ponderada;
    \item arquitetura hierárquica versus conjunta;
    \item encoder monolíngue versus multilíngue;
    \item probabilidades brutas versus calibradas;
    \item truncamento apenas no início versus início e fim e, se necessário,
    \textit{chunking};
    \item texto anonimizado com marcadores especiais versus anonimização sem os
    marcadores explícitos.
\end{itemize}

Cada ablação alterará um componente por vez, manterá partições e sementes e será
analisada de forma pareada. Perda, idioma, calibração e propagação hierárquica
correspondem às perguntas primárias; truncamento e mascaramento serão análises
secundárias no backbone vencedor, salvo se o piloto mostrar efeito suficiente
para justificar expansão previamente documentada.

\section{Fluxo experimental consolidado}
\label{sec:fluxo-experimental}

O procedimento completo seguirá a ordem abaixo. A posição do teste no fluxo
representa isolamento lógico: seus rótulos somente serão acessados depois que os
artefatos anteriores estiverem congelados.

\begin{enumerate}
    \item definir universo, estratos, probabilidades de inclusão e lotes;
    \item concluir a anotação simples, auditar a dupla anotação histórica e
    documentar concordância, divergências e exclusões;
    \item congelar dataset, taxonomia e manual, com hashes e manifests;
    \item executar anonimização, normalização e testes de consistência;
    \item detectar duplicatas, formar componentes conexos e grupos de contexto;
    \item criar e versionar as partições de seleção, calibração e teste;
    \item perfilar comprimentos e executar apenas as ablações pré-especificadas;
    \item realizar triagem de hiperparâmetros e validação cruzada;
    \item produzir e concatenar predições OOF para seleção de configuração;
    \item fixar época e treinar as cinco sementes finais;
    \item formar o ensemble, ajustar calibração e limiares no conjunto próprio;
    \item congelar checkpoint, tokenizer, calibrador, limiares e código;
    \item abrir uma única vez o teste para a avaliação principal;
    \item executar bootstrap pareado, análises por plataforma e eficiência;
    \item executar avaliações externa, temporal e funcional quando disponíveis;
    \item conduzir análise qualitativa reprodutível e publicar a documentação.
\end{enumerate}

\section{Ameaças à validade}
\label{sec:ameacas-validade}

As principais ameaças à interpretação dos resultados são:

\begin{itemize}
    \item \textbf{viés de seleção}: a amostra foi balanceada com base em sinais
    do Gemma e pode não reproduzir diretamente a prevalência do universo;
    \item \textbf{revisão simples}: a predominância de decisões
    \texttt{single} impede estimar concordância entre avaliadores na maior parte
    da v0.1;
    \item \textbf{classes raras}: algumas categorias possuem poucas dezenas de
    exemplos e não sustentam conclusões confirmatórias;
    \item \textbf{mudança temporal}: linguagem ofensiva, códigos e plataformas
    podem mudar depois do congelamento;
    \item \textbf{dependência de contexto}: o texto isolado pode não conter
    informações presentes na conversa, imagem, vídeo ou evento relacionado;
    \item \textbf{propagação de erros}: no pipeline hierárquico, um falso
    negativo binário impede a identificação dos tipos de ódio;
    \item \textbf{generalização limitada}: os resultados serão válidos para a
    composição observada de plataformas, períodos e categorias, não para todo
    conteúdo em português brasileiro.
\end{itemize}

Essas limitações deverão acompanhar qualquer divulgação de métricas ou uso do
modelo. A expansão contínua do dataset será orientada também pelos erros e pelas
lacunas de cobertura detectadas na v0.1.

\section{Tabelas para registro dos resultados}
\label{sec:tabelas-resultados}

As tabelas seguintes definem o formato de apresentação, mas não contêm valores
experimentais. Os campos marcados com ``---'' deverão ser preenchidos somente
após a execução do protocolo.

\begin{table}[htbp]
    \centering
    \caption{Modelo de tabela para a tarefa binária no teste fixo.}
    \label{tab:modelo-resultados-binarios}
    \resizebox{\textwidth}{!}{%
    \begin{tabular}{lrrrrrrr}
        \toprule
        \textbf{Modelo} & \textbf{Macro-F1} & \textbf{F1 hate} &
        \textbf{Precisão hate} & \textbf{Revocação hate} & \textbf{AUPRC} &
        \textbf{MCC} & \textbf{ECE} \\
        \midrule
        TF--IDF + regressão logística & --- & --- & --- & --- & --- & --- & --- \\
        Gemma & --- & --- & --- & --- & --- & --- & --- \\
        BERTimbau Base & --- & --- & --- & --- & --- & --- & --- \\
        BERTimbau Large & --- & --- & --- & --- & --- & --- & --- \\
        BERTabaporu Base & --- & --- & --- & --- & --- & --- & --- \\
        mBERT & --- & --- & --- & --- & --- & --- & --- \\
        XLM-R Base & --- & --- & --- & --- & --- & --- & --- \\
        \bottomrule
    \end{tabular}}
    \par\smallskip
    {\footnotesize Fonte: elaboração própria. Valores a serem preenchidos após
    os experimentos.\par}
\end{table}

\begin{table}[htbp]
    \centering
    \caption{Modelo de tabela para a tarefa multirrótulo.}
    \label{tab:modelo-resultados-multirrotulo}
    \resizebox{\textwidth}{!}{%
    \begin{tabular}{lrrrrrr}
        \toprule
        \textbf{Modelo} & \textbf{Macro-F1 elegível} &
        \textbf{Macro-F1 total} & \textbf{Micro-F1} & \textbf{PR-AUC} &
        \textbf{Hamming loss} & \textbf{Subset accuracy} \\
        \midrule
        TF--IDF + regressão logística & --- & --- & --- & --- & --- & --- \\
        BERTimbau Base & --- & --- & --- & --- & --- & --- \\
        BERTimbau Large & --- & --- & --- & --- & --- & --- \\
        BERTabaporu Base & --- & --- & --- & --- & --- & --- \\
        mBERT & --- & --- & --- & --- & --- & --- \\
        XLM-R Base & --- & --- & --- & --- & --- & --- \\
        \bottomrule
    \end{tabular}}
    \par\smallskip
    {\footnotesize Fonte: elaboração própria. Valores a serem preenchidos após
    os experimentos.\par}
\end{table}

\begin{table}[htbp]
    \centering
    \caption{Modelo de tabela para resultados por categoria de ódio.}
    \label{tab:modelo-resultados-categoria}
    \resizebox{\textwidth}{!}{%
    \begin{tabular}{lrrrrrr}
        \toprule
        \textbf{Categoria} & \textbf{$n$ teste} & \textbf{Precisão} &
        \textbf{Revocação} & \textbf{F1} & \textbf{IC95\% F1} & \textbf{PR-AUC} \\
        \midrule
        Ageísmo & --- & --- & --- & --- & --- & --- \\
        Classismo & --- & --- & --- & --- & --- & --- \\
        Gordofobia & --- & --- & --- & --- & --- & --- \\
        Capacitismo & --- & --- & --- & --- & --- & --- \\
        LGBTfobia & --- & --- & --- & --- & --- & --- \\
        Ódio político & --- & --- & --- & --- & --- & --- \\
        Racismo & --- & --- & --- & --- & --- & --- \\
        Intolerância religiosa & --- & --- & --- & --- & --- & --- \\
        Sexismo & --- & --- & --- & --- & --- & --- \\
        Xenofobia & --- & --- & --- & --- & --- & --- \\
        \bottomrule
    \end{tabular}}
    \par\smallskip
    {\footnotesize Fonte: elaboração própria. Categorias abaixo do limiar
    confirmatório deverão ser identificadas como exploratórias.\par}
\end{table}

\begin{table}[htbp]
    \centering
    \caption{Modelo de tabela para desempenho e eficiência computacional.}
    \label{tab:modelo-eficiencia}
    \resizebox{\textwidth}{!}{%
    \begin{tabular}{lrrrrr}
        \toprule
        \textbf{Modelo} & \textbf{Macro-F1} & \textbf{Parâmetros} &
        \textbf{GPU-horas} & \textbf{Memória máxima} &
        \textbf{Tempo/1.000 textos} \\
        \midrule
        BERTimbau Base & --- & --- & --- & --- & --- \\
        BERTimbau Large & --- & --- & --- & --- & --- \\
        BERTabaporu Base & --- & --- & --- & --- & --- \\
        mBERT & --- & --- & --- & --- & --- \\
        XLM-R Base & --- & --- & --- & --- & --- \\
        \bottomrule
    \end{tabular}}
    \par\smallskip
    {\footnotesize Fonte: elaboração própria. Valores a serem preenchidos após
    os experimentos.\par}
\end{table}

\section{Critérios de conclusão da etapa v0.1}
\label{sec:criterios-conclusao}

A etapa de treinamento da v0.1 será considerada metodologicamente concluída
quando:

\begin{enumerate}
    \item o primeiro lote completo que alcance ao menos 21.248 decisões válidas
    tiver sido congelado e versionado;
    \item a suficiência amostral tiver sido examinada e as análises confirmatórias
    e exploratórias estiverem separadas;
    \item a dupla anotação histórica, as medidas de concordância e as exclusões
    por divergência tiverem sido documentadas, e a ausência de dupla anotação
    contemporânea estiver explicitamente declarada como limitação;
    \item as partições tiverem sido verificadas quanto a duplicatas e vazamento;
    \item todos os baselines e cinco encoders tiverem sido executados nas mesmas
    divisões;
    \item os hiperparâmetros tiverem sido definidos por OOF e os limiares e a
    calibração pelo conjunto próprio, sem uso do teste;
    \item as cinco sementes tiverem sido executadas para cada configuração
    selecionada;
    \item o ensemble final tiver sido fixado sem seleção posterior de semente;
    \item as métricas globais, ponderadas, por plataforma e por categoria
    tiverem sido calculadas com intervalos de confiança e comparações pareadas;
    \item a análise funcional e a análise qualitativa de erros tiverem sido
    documentadas;
    \item os manifests, checkpoints, hashes, tabelas e fichas de documentação
    tiverem sido arquivados de forma reprodutível.
\end{enumerate}

% -----------------------------------------------------------------------------
% Referências incorporadas ao próprio capítulo
% -----------------------------------------------------------------------------

\begin{thebibliography}{99}

\bibitem[Bender e Friedman(2018)]{bender2018datastatements}
BENDER, Emily M.; FRIEDMAN, Batya.
Data statements for natural language processing: toward mitigating system bias
and enabling better science.
\textit{Transactions of the Association for Computational Linguistics}, v. 6,
p. 587--604, 2018.
DOI: \url{https://doi.org/10.1162/tacl_a_00041}.

\bibitem[Cohen(1960)]{cohen1960kappa}
COHEN, Jacob.
A coefficient of agreement for nominal scales.
\textit{Educational and Psychological Measurement}, v. 20, n. 1, p. 37--46,
1960.
DOI: \url{https://doi.org/10.1177/001316446002000104}.

\bibitem[Conneau et al.(2020)]{conneau2020xlmr}
CONNEAU, Alexis et al.
Unsupervised cross-lingual representation learning at scale.
In: \textit{Proceedings of the 58th Annual Meeting of the Association for
Computational Linguistics}. Online: Association for Computational Linguistics,
2020. p. 8440--8451.
Disponível em: \url{https://arxiv.org/abs/1911.02116}.

\bibitem[da Costa et al.(2023)]{dacosta2023bertabaporu}
DA COSTA, Pablo Botton; PAVAN, Matheus Camasmie; DOS SANTOS, Wesley Ramos;
DA SILVA, Samuel Caetano; PARABONI, Ivandre.
BERTabaporu: assessing a genre-specific language model for Portuguese NLP.
In: \textit{Proceedings of the 14th International Conference on Recent Advances
in Natural Language Processing}. Varna: INCOMA, 2023. p. 217--223.
DOI: \url{https://doi.org/10.26615/978-954-452-092-2_024}.

\bibitem[da Costa(2026)]{dacosta2026bertabaporu}
DA COSTA, Pablo Botton.
\textit{BERTabaporu Base Uncased: model card}.
Hugging Face, 2026. Disponível em:
\url{https://huggingface.co/pablocosta/bertabaporu-base-uncased}.
Acesso em: 11 set. 2026.

\bibitem[Devlin et al.(2019)]{devlin2019bert}
DEVLIN, Jacob; CHANG, Ming-Wei; LEE, Kenton; TOUTANOVA, Kristina.
BERT: pre-training of deep bidirectional transformers for language understanding.
In: \textit{Proceedings of the 2019 Conference of the North American Chapter of
the Association for Computational Linguistics: Human Language Technologies}.
Minneapolis: Association for Computational Linguistics, 2019. p. 4171--4186.
Disponível em: \url{https://arxiv.org/abs/1810.04805}.

\bibitem[Dodge et al.(2020)]{dodge2020finetuning}
DODGE, Jesse et al.
Fine-tuning pretrained language models: weight initializations, data orders,
and early stopping.
\textit{arXiv preprint arXiv:2002.06305}, 2020.
Disponível em: \url{https://arxiv.org/abs/2002.06305}.

\bibitem[Efron e Tibshirani(1993)]{efron1993bootstrap}
EFRON, Bradley; TIBSHIRANI, Robert J.
\textit{An introduction to the bootstrap}.
New York: Chapman \& Hall, 1993.
DOI: \url{https://doi.org/10.1007/978-1-4899-4541-9}.

\bibitem[Google BERT(2026)]{google2026mbert}
GOOGLE BERT.
\textit{BERT multilingual base model (cased): model card}.
Hugging Face, 2026.
Disponível em:
\url{https://huggingface.co/google-bert/bert-base-multilingual-cased}.
Acesso em: 1 set. 2026.

\bibitem[Guo et al.(2017)]{guo2017calibration}
GUO, Chuan; PLEISS, Geoff; SUN, Yu; WEINBERGER, Kilian Q.
On calibration of modern neural networks.
In: \textit{Proceedings of the 34th International Conference on Machine Learning}.
Sydney: PMLR, 2017. v. 70, p. 1321--1330.
Disponível em: \url{https://arxiv.org/abs/1706.04599}.

\bibitem[Holm(1979)]{holm1979procedure}
HOLM, Sture.
A simple sequentially rejective multiple test procedure.
\textit{Scandinavian Journal of Statistics}, v. 6, n. 2, p. 65--70, 1979.

\bibitem[Gururangan et al.(2020)]{gururangan2020dontstop}
GURURANGAN, Suchin et al.
Don't stop pretraining: adapt language models to domains and tasks.
In: \textit{Proceedings of the 58th Annual Meeting of the Association for
Computational Linguistics}. Online: Association for Computational Linguistics,
2020. p. 8342--8360.
DOI: \url{https://doi.org/10.18653/v1/2020.acl-main.740}.

\bibitem[Krippendorff(2018)]{krippendorff2018content}
KRIPPENDORFF, Klaus.
\textit{Content analysis: an introduction to its methodology}.
4. ed. Thousand Oaks: SAGE, 2018.

\bibitem[Meta AI(2026)]{meta2026xlmr}
META AI.
\textit{XLM-RoBERTa base: model card}.
Hugging Face, 2026.
Disponível em: \url{https://huggingface.co/FacebookAI/xlm-roberta-base}.
Acesso em: 1 set. 2026.

\bibitem[Mitchell et al.(2019)]{mitchell2019modelcards}
MITCHELL, Margaret et al.
Model cards for model reporting.
In: \textit{Proceedings of the Conference on Fairness, Accountability, and
Transparency}. New York: ACM, 2019. p. 220--229.
DOI: \url{https://doi.org/10.1145/3287560.3287596}.

\bibitem[NeuralMind(2026)]{neuralmind2026bertimbau}
NEURALMIND.
\textit{BERTimbau Base and Large: model cards}.
Hugging Face, 2026.
Disponível em:
\url{https://huggingface.co/neuralmind/bert-base-portuguese-cased} e
\url{https://huggingface.co/neuralmind/bert-large-portuguese-cased}.
Acesso em: 1 set. 2026.

\bibitem[Röttger et al.(2021)]{rottger2021hatecheck}
RÖTTGER, Paul; VIDGEN, Bertie; NGUYEN, Dong; WASEEM, Zeerak; MARGETTS, Helen;
PIERREHUMBERT, Janet.
HateCheck: functional tests for hate speech detection models.
In: \textit{Proceedings of the 59th Annual Meeting of the Association for
Computational Linguistics and the 11th International Joint Conference on
Natural Language Processing}. Online: Association for Computational Linguistics,
2021. p. 41--58.
DOI: \url{https://doi.org/10.18653/v1/2021.acl-long.4}.

\bibitem[Saito e Rehmsmeier(2015)]{saito2015precisionrecall}
SAITO, Takaya; REHMSMEIER, Marc.
The precision-recall plot is more informative than the ROC plot when evaluating
binary classifiers on imbalanced datasets.
\textit{PLOS ONE}, v. 10, n. 3, e0118432, 2015.
DOI: \url{https://doi.org/10.1371/journal.pone.0118432}.

\bibitem[Sechidis, Tsoumakas e Vlahavas(2011)]{sechidis2011stratification}
SECHIDIS, Konstantinos; TSOUMAKAS, Grigorios; VLAHAVAS, Ioannis.
On the stratification of multi-label data.
In: \textit{Machine Learning and Knowledge Discovery in Databases}.
Berlin: Springer, 2011. p. 145--158.
DOI: \url{https://doi.org/10.1007/978-3-642-23808-6_10}.

\bibitem[Souza, Nogueira e Lotufo(2020)]{souza2020bertimbau}
SOUZA, Fábio; NOGUEIRA, Rodrigo; LOTUFO, Roberto.
BERTimbau: pretrained BERT models for Brazilian Portuguese.
In: \textit{Intelligent Systems: 9th Brazilian Conference, BRACIS 2020}.
Cham: Springer, 2020. p. 403--417.
DOI: \url{https://doi.org/10.1007/978-3-030-61377-8_28}.

\bibitem[Wolf et al.(2020)]{wolf2020transformers}
WOLF, Thomas et al.
Transformers: state-of-the-art natural language processing.
In: \textit{Proceedings of the 2020 Conference on Empirical Methods in Natural
Language Processing: System Demonstrations}. Online: Association for
Computational Linguistics, 2020. p. 38--45.
DOI: \url{https://doi.org/10.18653/v1/2020.emnlp-demos.6}.

\end{thebibliography}
